
\documentclass[a4paper,11pt]{article}
\usepackage{./jheppub}
\usepackage{tikz-cd}
\usepackage{amsmath}
\usepackage{array,booktabs}
\usepackage{hyperref}
\usepackage{amsthm}
\newtheorem{theorem}{Theorem}
\newtheorem{lemma}[theorem]{Lemma}
\newtheorem{corollary}[theorem]{Corollary}
\newtheorem{proposition}[theorem]{Proposition}
\theoremstyle{definition}
\newtheorem{definition}{Definition}[section]
\newtheorem{prop}{Conjecture}
\newtheorem{cor}{Corollary}
\usepackage{enumerate}
\usepackage[utf8]{inputenc}
\usepackage{amssymb,amsbsy,mathtools,amsmath}
\usepackage{braket}
\usepackage{graphicx}
\usepackage{xcolor}
\usepackage{hyperref}
\usepackage{natbib}
\usepackage{bm}
\usepackage{soul}
\newcommand{\E}{\mathbb{E}}
\newcommand{\dd}{\,\mathrm{d}}
\newcommand{\cH}{\mathcal{H}}
\newcommand{\KLL}{D_{\mathrm{KL}}}
\newcommand{\Mop}{\mathsf{M}}

\usepackage[nottoc]{tocbibind}

\usepackage{dsfont}
\let\mathbbm\mathds
\usepackage{amsthm}
\usepackage{amssymb}
\newcommand{\id}{\mathbb{I}}
\newcommand{\Tr}{\text{Tr}}
\newcommand{\Div}{\operatorname{div}}
\usepackage{amsmath}
\usepackage{braket}
\usepackage{graphicx}
\usepackage{caption}
\newcommand{\mA}{\mathcal{A}}
\newcommand{\mB}{\mathcal{B}}
\newcommand{\mP}{\mathcal{P}}
\newcommand{\mR}{\mathcal{R}}
\newcommand{\nn}{\nonumber}
\newcommand{\mN}{\mathcal{N}}
\newcommand{\mI}{\mathcal{I}}
\newcommand{\mL}{\mathcal{L}}
\newcommand{\lb}{\left(}
\newcommand{\rb}{\right)}
\newcommand{\mH}{\mathcal{H}}
\usepackage{subcaption}
\newcommand{\tOO}{\tilde{\mathcal{O}}}
\newcommand{\mO}{\mathcal{O}}
\newcommand{\mK}{\mathcal{K}}
\newcommand{\ep}{\epsilon}
\newcommand{\mbR}{\mathbb{R}}
\newcommand{\mM}{\mathcal{M}}
\newcommand{\mT}{\mathcal{T}}
\newcommand{\NL}[1]{\textcolor{blue}{#1}}
\newcommand{\p}{\partial}
\newcommand{\tr}{\text{tr}}
\newcommand{\tT}{\tilde{T}}
\newcommand{\Ad}{\operatorname{Ad}}
\newcommand{\Dom}{\operatorname{Dom}}
\newcommand{\mE}{\mathcal{E}}
\newcommand{\relcomm}[2]{#1'\cap #2}

\newcommand{\SO}[1]{\textcolor{purple}{#1}}
\newcommand{\KL}[1]{\textcolor{brown}{#1}}
\newcommand{\YC}[1]{\textcolor{teal}{#1}}
\newcommand{\cL}{\mathcal{L}}
\setlength{\emergencystretch}{2em}
\begin{document}

\title{Subsystem Flow: Quantum Error Correction,\\ Petz Recovery, and Ergodic Inclusions}

\author{Yidong Chen$^1$,}
\affiliation{$^1$Department of Mathematics, University of Illinois at Urbana-Champaign, Illinois, IL 61801, USA}
\author{Nima Lashkari$^2$,}
\author{Kwing Lam Leung$^2$}
\affiliation{$^2$Department of Physics and Astronomy, Purdue University, West Lafayette, IN 47907, USA}
\emailAdd{yidongc2@illinois.edu}
\emailAdd{nima@purdue.edu}
\emailAdd{leung60@purdue.edu}

\abstract{
When the closed unitary dynamics of a bipartite quantum system $A=BR$ leaves subsystem $B$ invariant, we obtain a Quantum Error Correction (QEC) code: The code subspace created by observables in $B$ is correctable against the erasure of $R$. We define an algebraic flow called \emph{subsystem flow} and argue that it captures the failure of this error-correction property and the mixing of the retained and discarded observables. We generalize this to arbitrary quantum subsystems (von Neumann Subfactors).

\NL{In the opposite regime, ergodic modular flows are those where no subsystem is invariant.} We translate key notions of the ergodic hierarchy such as generation, mixing, invariant cores and forgetfulness into the language of QEC. The analytic continuation of the subsystem flow defines Petz and universal recovery maps and a class of positive \emph{inclusion parent Hamiltonians}. We argue that the inclusions with properties such as
 common cyclic and separating vectors, and $L^2$-nuclearity similar to those of subregions of quantum field theory, provide examples of highly ergodic subsystems in this hierarchy.}

\maketitle
\section{Introduction}\label{sec:intro}

A subsystem in QFT need not correspond to a tensor factor of the Hilbert space.  The observables accessible in a spacetime region, or observables accessible along an observer's timeline, form a von Neumann algebra $\mN$ (with trivial center) embedded in a larger algebra $\mM$.  The inclusion $\mN\subset\mM$
is therefore a natural algebraic notion of a quantum subsystem.  This viewpoint is particularly useful in continuum QFT, where local algebras are special infinite-dimensional von Neumann algebras (typically type III factors), and neither a density matrix nor a Hilbert-space tensor factorization is intrinsic to a region.

In quantum mechanics, we associate to every density matrix $\rho$ a unitary dynamical flow $\rho^{it}$ called modular flow. More generally, in a von Neumann algebra $\mM$, every faithful normal state $\omega$ defines modular dynamics, $
    \sigma_t^{\omega;\mM}
    =
    \operatorname{Ad}\Delta_{\mM}^{it}$,
while the restriction of the same state to $\mN$ has its own intrinsic modular evolution $\rho_t^{\omega,N}$. There are thus two natural notions of subsystem dynamics: evolve $\mN$ using the ambient modular flow of $\mM$, or evolve it using the modular flow inferred from the reduced state.  The mismatch between these two motions is \emph{subsystem flow} that we study in this paper: 
\begin{align}
    \mathcal{D}_{\mN\subset\mM}(t)=\sigma_{-t}^{\omega;\mM}\circ \sigma_{t}^{\omega;\mN}
\end{align}

In vacuum QFT, modular dynamics of wedge algebras has a geometric interpretation. The Bisognano--Wichmann theorem identifies vacuum wedge modular flow with Lorentz boosts~\cite{BisognanoWichmann1975}, while for a ball in a CFT it is generated by a conformal Killing flow~\cite{CasiniHuertaMyers2011}. For suitable half-sided modular inclusions the difference between modular evolutions reconstructs null translations~\cite{Wiesbrock1993,Borchers1992}. These are familiar examples in which the algebraic evolution studied below can be interpreted geometrically.

This viewpoint is also important in holography, where modular flow and relative entropy are related to entanglement-wedge reconstruction and QEC~\cite{AlmheiriDongHarlow2015,JafferisLewkowyczMaldacenaSuh2016}. The present paper develops the general information-theoretic and operator-algebraic mechanism. Explicit wedge, double-cone, and AdS constructions, and their implications for modular chaos and emergent spacetime time evolution, are developed in the companion paper~\cite{companionQFT}.

In this paper, under the closed dynamics of a general bipartite system, what information does a quantum subsystem preserve or lose? 
To formalize this problem mathematically and as generally as possible, we consider a von Neumann subalgebra of $\mM$ represented on the GNS Hilbert space of $\mH_\Omega$, where $\Omega$ is cyclic and separating for $\mM$:
\begin{equation}
    \mH_{\mN}=\overline{\mN\Omega},
    \qquad
    V:\mH_{\mN}\hookrightarrow\mH_{\mM} ,
\end{equation}
and $V$ is the embedding map, and let $\Delta_{\mN}$ be the modular operator of the restricted state on $\mH_{\mN}$.  We define the {\it subsystem flow}
\begin{equation}
    D_{\mN\subset\mM}(t)
    =
    \overline{\Delta_{\mM}^{-it}V\Delta_{\mN}^{it}}:\mH_{\mN}\to \mH_{\mM}
    \label{eq:intro-compressed-flow}
\end{equation}
When $\Omega$ is cyclic for both algebras, $\mH_{\mN}=\mH_{\mM}$ and this reduces to Borchers' characteristic function \cite{revolutionizing}
\begin{equation}
    D_{\mN\subset\mM}(t)
    =
    \Delta_{\mM}^{-it}\Delta_{\mN}^{it},
    \label{eq:intro-subsystem-flow}
\end{equation}
We will argue that by factoring out the intrinsic modular motion, the subsystem flow isolates precisely the failure of the subsystem to evolve autonomously; i.e., the failure of the encoding $V$ (code subspace is $V\mH_{\mN}$) to protect against the restriction to $\mN$ as an erasure QEC code.

\paragraph{Quantum error correction and the modular ergodic hierarchy:} We denote the space of normal states of algebra $\mM$ by $\mM_*$. Then, the restriction of states
\begin{equation}
    \mathcal T:\mM_*\to\mN_*,
    \qquad
    \mathcal T(\varphi)=\varphi|_{\mN},
\end{equation}
is the algebraic analogue of erasing the information complementary to $\mN$.  A normal conditional expectation $\mE:\mM\to\mN$ provides an encoding in the opposite direction, projection operators in $\mM$ down to $\mN$. It is the Heisenberg picture analog of a recovery map \cite{furuya2022real,Faulkner2020}. As we will review, a classic result in operator algebras (Takesaki's theorem) characterizes the necessary and sufficient condition for the existence of such a QEC code: a reference-preserving conditional expectation exists precisely when the subsystem is invariant under the modular dynamics of $\omega_{\mM}$:
\begin{equation}
    \sigma_t^{\omega;\mM}(\mN)=\mN
    \qquad\forall\,t
    \label{eq:intro-autonomy}
\end{equation}
\cite{Takesaki1972,Faulkner2020,furuya2022real}.  This is also equivalent to
$\Delta_{\mM}^{it}V=V\Delta_{\mN}^{it}$ which implies $D_{\mN\subset\mM}(t)=V$ for all $t$.  When a closed dynamics of $\mM$ preserves subsystem $\mN$, we say $\mN$ evolves autonomously, and $\mN$ forms a subsystem QEC code. We have well-defined dynamics on the subsystem that we will refer to, as a subtheory. Takesaki's theorem states that modular autonomy, a reference-preserving conditional expectation, exact error correction against restriction to $\mN$, and trivial subsystem flow are equivalent descriptions of the same algebraic structure.  From the point of view of the inclusion, the relative commutant $\mR=\mN'\cap \mM$, if it exists, plays the role of the environment degrees of freedom we trace over (erasure). We will see that when $\mR$ exists, and $\mR\vee \mN=\mM$ (conormal inclusion) the condition of invariant subsystem becomes a Knill--Laflamme condition.  The relation with equality and approximate equality in data processing connects this picture directly with Petz sufficiency and recovery~\cite{Petz1986,Petz1988,JungeRennerSutterWildeWinter2018,FaulknerHollandsSwingleWang2022}.

It is important to point out that the invariance condition in (\ref{eq:intro-autonomy}) does not mean that modular flow acts trivially in $\mN$ \footnote{When that happens, we say $\mN$ is the fixed point subalgebra of the dynamics.}. In the operator algebraic literature, e.g., \cite{revolutionizing}, sometimes the condition in (\ref{eq:intro-autonomy}) is referred to as modular covariance, instead of invariance. 

If subsystem $\mN$ is not exactly preserved, the mixing of $\mN$ in $\mM$ generates entropy. Our subsystem flow is, in fact, a generating function for this entropy. To see this, consider the vector $\ket{\Omega}$ and its corresponding reduced states $\omega$ that define modular flows, and another unit vector in the code subspace $V\Psi$ with $\Psi\in\mH_{\mN}$ and its corresponding reduced state $\varphi$. Then, under the appropriate domain and finiteness assumptions,
\begin{align}
 -i\left.
 \partial_t
 \langle V\Psi|D_{\mN\subset\mM}(t)|\Psi\rangle
 \right|_{t=0}
 &=
 \langle\Psi|(V^\dagger K_{\mM}V-K_{\mN})|\Psi\rangle
 \nonumber\\
 &=\bigl[S_{\mM}(\varphi\|\omega)-S_{\mN}(\varphi\|\omega)\bigr]
 \nonumber\\
 &\quad+\bigl[S_{\mN'}(\varphi\|\omega)-S_{\mM'}(\varphi\|\omega)\bigr],
 \label{eq:intro-entropy}
\end{align}
where $K_{\mM}=-\log\Delta_{\mM}$ and $K_{\mN}=-\log\Delta_{\mN}$ are modular Hamiltonians, and
 the $\mN'$ relative entropy is understood on the support $e_{\mN}=VV^\dagger$. The two nonnegative terms are the Petz sufficiency gaps for the inclusion and its commutant. Thus, the quadratic form associated with the infinitesimal subsystem flow measures failure of recoverability state by state.

Away from exact QEC of invariant $\mN$, we organize subsystem dynamics through the history algebra and the invariant core,
\begin{equation}
    \mN_{\rm hist}
    =
    \bigvee_{t\in\mathbb R}\sigma_t^\omega(\mN),
    \qquad
    \mN_{\rm core}
    =
    \bigcap_{t\in\mathbb R}\sigma_t^\omega(\mN).\label{eq:intro-core}
\end{equation}
They are respectively the smallest modularly invariant algebra containing $\mN$ and the largest modularly invariant algebra contained in it. Their commutant duality relates generation by $\mN$ to triviality of the invariant core of its environment. In a truly chaotic flow, we would like to have $\mN_{\rm hist}=\mM$ and $\mN_{\rm core}=\mathbb C1$. We will show that common cyclicity gives generation, while standard inclusions also have trivial invariant core. These algebraic conditions are distinct from asymptotic forgetfulness or mixing, which requires a late-time limit.

\paragraph{Positive inclusion parent Hamiltonians:}
Analytic continuation of subsystem flow to the strip $0\leq\operatorname{Im}z\leq1/2$ produces contractions even when there is no state-preserving conditional expectation. For $0<\alpha<1/2$ we define
\begin{equation}
 \rho_{\alpha,\mN\subset\mM}=D(i\alpha)^\dagger D(i\alpha),\qquad
 Q_{\alpha,\mN\subset\mM}
 =-\frac{1}{2\sin(2\pi\alpha)}\log\rho_{\alpha,\mN\subset\mM}\geq0 .
 \label{eq:intro-Halpha}
\end{equation}
Here the logarithm is defined on the support of $\rho_\alpha$. Since $D(i\alpha)$ fixes $\Omega$, the reference vector is a zero-energy state. The first key result is that the full ground-state space, and not just the ground state among vectors $a\Omega$, is determined by the largest invariant subtheory:
\begin{equation}
 \ker(1-\rho_{\alpha,\mN\subset\mM})
 =\overline{\mN_{\rm core}^{\,\mM}\Omega}.
 \label{eq:intro-ground-core}
\end{equation}
Thus a trivial invariant core is equivalent to a unique ground state when $\rho_\alpha$ has full support. This is an algebraic distinction between a subsystem that retains a closed subtheory and one whose observables are genuinely mixed by the ambient modular dynamics.

Our second main result is that, when this core is trivial and $\rho_\alpha$ has full support, the positive inclusion parent Hamiltonian generates an automorphism group $\gamma_t^\alpha=\operatorname{Ad}e^{itQ_\alpha}$ on the subsystem GNS Hilbert space, with
\begin{equation}
 \boxed{\qquad
 \bigvee_{t\in\mathbb R}\gamma_t^\alpha(\mN)=B(\mH_{\mN}) .\qquad}
 \label{eq:intro-full-generation}
\end{equation}
The original modular orbit of $\mN$ stays inside $\mM$, whereas the orbit generated by $Q_\alpha$ can explore operators outside that ambient algebra. For a standard inclusion $\mH_{\mN}=\mH_{\mM}$, and the same result gives the whole $B(\mH_{\mM})$. In particular, neither an assumption of half-sided modular nesting nor an $L^2$-nuclearity hypothesis is needed for this two-sided generation statement. The stronger future- and past-generation results use additional spectral assumptions.

\paragraph{Petz maps, canonical endomorphisms and inclusion spectra:}
The upper boundary $D(i/2)$ gives the Petz--Accardi--Cecchini dual, a generalized conditional expectation that exists even when a state-preserving conditional expectation does not. Its rotated versions yield the universal recovery channel. For a common cyclic and separating reference vector the Petz dual becomes the canonical endomorphism implemented by $\Gamma=J_{\mN}J_{\mM}$. At the symmetric point $\alpha=1/4$, the positive operator $\rho_{1/4}$ is its symmetric $L^2$ transfer operator; its logarithm therefore provides a canonical positive generator without taking a branch-dependent logarithm of $\Gamma$. We discuss Longo's mixing theorem for the discrete canonical evolution, the spectral consequences of the three-lines theorem, and the role of nuclearity. Neither the unitary orbit generated by $Q_{1/4}$ nor its global-generation property is, by itself, a claim about geometric time evolution.

Section~\ref{section:basics} reviews quantum subsystems and the operator-algebraic setting, and Section~\ref{sec:QEC} explains modular autonomy as QEC. Section~\ref{sec:ergodic} relates Petz recovery, subsystem flow and the ergodic hierarchy. Section~\ref{sec:probing-outside-M} constructs inclusion parent Hamiltonians, proves their ground-state and generation theorems, and studies the symmetric $L^2$ transfer operator and nuclearity. Further algebraic definitions, finite-dimensional examples and the classical ergodic analogy are given in the appendices. The companion QFT paper~\cite{companionQFT} addresses half-sided modular chaos, geometric translations and emergent time in concrete QFT and holographic inclusions.

\section{Operator algebra background}\label{section:basics}

The observables of a physical quantum system form a von Neumann algebra. The center of this algebra describes its classical observables (superselection sectors). A quantum system with no non-trivial classical superselection observables is a von Neumann \emph{factor}. For a finite quantum system, such factors are
finite matrix algebras\footnote{The factor associated to a $d$-level quantum system is the algebra of $d\times d$ complex matrices, $M(d,\mathbb{C})$}. In many-body finite quantum systems, the algebra of observables is the tensor product of single-body algebras. For a composite system $A=BR$ with $\mathcal H_{BR}=\mathcal H_B\otimes\mathcal H_R$ and density matrices $\rho_{BR},\rho_B,\rho_R$, the algebras are $\mN=B(\mathcal H_B)\otimes1$, $\mR=1\otimes B(\mathcal H_R)$ and $\mM=\mN\vee\mR=B(\mathcal H_{BR})$.\footnote{Throughout the paper an inclusion of observable algebras is written $\mN\subset\mM$, $\mR=\mN'\cap\mM$ is its relative commutant, and $\mM'$, $\mN'$ are the commutants on the Hilbert space. Roman capitals are reserved for physical systems.} 
The algebra of $R$ is the \emph{relative commutant} of this inclusion, $\mR=\mN'\cap\mM$.
The partial trace over $R$ in a fixed uncorrelated state $\chi_R$ of the environment, read in the Heisenberg picture, is a unital completely positive projection (conditional expectation)
\begin{align}
    \mE_{\chi_R}:\mM\to\mN,\qquad
    \mE_{\chi_R}(b\otimes r)=\chi_R(r)\,(b\otimes 1)\ ,
    \label{eq:slice-expectation-intro}
\end{align}
and a simple example of a coarse-graining to a subsystem. In infinite-dimensional quantum systems (von Neumann algebras) an inclusion of factors (subfactors) $\mN\subset\mM$ plays the role of a subsystem. However, before generalizing to infinite dimensions, it is helpful to remind ourselves of finite-dimensional (matrix) algebra subalgebras of $B(\mH)$. Following the QEC terminology, we classify them as:
\begin{enumerate}
 \item {\bf Subspace inclusions:} Let $\mathcal H=P\mathcal H\oplus(1-P)\mathcal H$.
 The unital subalgebra is
 \[
  \mN_P=P B(\mathcal H)P\oplus\mathbb C(1-P).
 \]
 Its center is $\mathbb CP\oplus\mathbb C(1-P)$. For a fixed state $\chi_{1-P}$
 on the complementary corner, the conditional expectation is
 \[
  \mE_{\chi_{1-P}}(a)=PaP+
  \chi_{1-P}\bigl((1-P)a(1-P)\bigr)(1-P).
 \]
 It preserves precisely the block-diagonal states whose complementary block
 is proportional to the density matrix of $\chi_{1-P}$.
 Block diagonality alone does not suffice for a fixed expectation.
In the Hilbert space of a single quantum particle moving on some ambient space $X$, a net of inclusion of regions in $X$, $I_i\subset I_{i+1}$, defines a net of inclusion of projections $P_i\leq P_{i+1}$ and a net of subspace inclusions $\mN_{P_i}\subseteq\mN_{P_{i+1}}$. 
 
 \item {\bf Subsystem (subfactor) inclusions:} In $\mathcal H=\mathcal H_B\otimes\mathcal H_R$ we consider the subalgebra $B(\mathcal H_B)\otimes1_R$.
 Its expectations are those of \eqref{eq:slice-expectation-intro}, which preserve all uncorrelated states of the form $\rho_A=\chi_R\otimes \rho_B$. In the Hilbert space of a many-body quantum system defined on some lattice $X$, a net of inclusion of regions, $I_i\subset I_{i+1}$, gives a net of inclusion of conditional expectations $\chi_i \leq\chi_{i+1}$, and a net of subsystem inclusions $\mN_{\chi_i}\subset \mN_{\chi_{i+1}}$. We will see in Section~\ref{subsec:quantum-coarse} that in some continuum spaces $X$ like background spacetime in QFT, the net of subfactors (subsystems) exists, even though normal conditional expectations might not exist. This has to do with the absence of uncorrelated states of touching regions.
 
 \item {\bf Mixed inclusions:} A general unital finite-dimensional subalgebra takes the block diagonal form
 \[
  \mathcal H=\bigoplus_a\mathcal H_B^{(a)}\otimes\mathcal H_R^{(a)},\qquad
  \mN=\bigoplus_a B(\mathcal H_B^{(a)})\otimes1_R^{(a)}.
 \]
 The center records the block label $a$. A choice of normalized hidden state $\chi_R^{(a)}$ specifies an
 expectation in each block. The states preserved under the projections by conditional expectation take the Markovian form
 \[
  \rho=\bigoplus_a p_a\,\rho_B^{(a)}\otimes\chi_R^{(a)}.
 \]
 With the label stored in a classical register $Z$, these have $I(B:R|Z)=0$.
\end{enumerate}
Note that the classification above is central to the algebraic formulation of QEC where the subalgebra $\mN$ is the set of unitarily correctible operators under the erasure of $R$ \cite{BenyKempfKribs2007} \footnote{The subspace construction resolves alternatives among microstates, while a tensor
subsystem retains simultaneously specified degrees of freedom. The same distinction
organizes the classical constructions in Appendix~\ref{sec:classical}.
}. In QEC jargon, the corresponding codes are called {\it subspace code}, {\it subsystem code}, and {\it mixed codes}, respectively.

In all inclusions above, the relative commutant includes both $Z(\mN)$ the center of $\mN$, and the algebra of $R$:
\begin{align}
    Z(\mN)\vee B(\mH_R)\subset \mR=\mN'\cap \mM\ .
\end{align}
Since $\mM_A=B(\mH_A)$ and $\mN$ is a von Neumann subalgebra, in all the examples above, its double relative commutant is equal to itself:
\begin{align}
    \mN=(\mN'\cap \mM)'\cap\mM
\end{align}
Such an inclusion is called {\it normal}. However, when $\mM$ has a center, the inclusion need not be normal, as is the case of the inclusion of Abelian algebras. 

However, when the algebra $\mN$ has a non-trivial center (cases 1 and 3), the algebra and its relative commutant do not additively generate the whole $\mM$. Whereas in the second case, we have $\mM=\mN\vee \mR$ and we say the inclusion is {\it conormal}.

An interesting special case of mixed inclusions is when all $\mH_R^{(a)}$ are one-dimensional, and all $\mH_B^{(a)}$ have the same dimension; e,g; the subalgebra $\mN=\oplus_{a=1}^d B(\mH^{(a)})=B(\mH^{(a)})\otimes D_{d\times d}$, where $D_{d\times d}$ is the algebra of diagonal matrices. Then, permutations of superselection labels $a\to \sigma(a)$ by the unitaries
 \begin{align}
    \forall b\in\mN: \alpha_\sigma(b)=U_\sigma b U_\sigma^\dagger,\qquad U_{\sigma}=\sum_a \ket{i_{\sigma(a)}}\bra{i_a}
 \end{align}
act as outer automorphisms of the subalgebra $\mN$; meaning that $U_\sigma$ do not belong to $\mN$. In these cases, we can always enlarge the subalgebra by adding $U_\sigma$ to it. In Mathematical terminology, this is called taking the crossed product with the outer automorphism permutation group $S_d$: $\mM=\mN\rtimes_\alpha S_d$. In this case, the relative commutant is the Abelian (the center of $\mB$). In finite dimensions, we cannot have an inclusion of factors with trivial relative commutant. All automorphisms of finite-dimensional factors are inner. However, as we will see in Section \NL{\ref{subsec:types-modular}}, some infinite-dimensional factor algebras (type II and III) can have outer automorphisms, and taking the crossed product with the automorphism group results in inclusions with trivial relative commutant. In the operator algebraic language, they are called {\it Irreducible inclusions} or {\it singular inclusions}. We remind the readers not to confuse this notion of irreducibility with irreducibility of an algebra under a dynamical law introduced in later Sections.

In this work, we are primarily interested in the infinite-dimensional factors describing
subsystems of many-body quantum systems in the thermodynamic and continuum limit.
In these systems, inclusions can have features that are impossible in
finite-dimensional tensor-product examples: 
\begin{itemize}
\item We can have inclusion of factors with trivial relative commutant. This automatically means that the inclusion is not normal.

\item We can have inclusion of factors with non-trivial relative commutant $\mR$ and conormal $\mM=\mN\vee \mR$, but no normal conditional expectations. This can be interpreted as inclusions as the absence of uncorrelated states of $\mN$ and $\mR$.

\item We can have inclusions with common cyclic and separating vectors, or even stronger, we can {\it standard inclusions} meaning that $\mN$, $\mM$ and $\mR$ all share a common cyclic and separating vector.
\end{itemize}
We introduce the modular notions needed for these examples briefly below.
For completeness, the general definitions of von Neumann algebras, the
types of factors, and the distinction between arbitrary unitary dynamics
and modular dynamics are reviewed in Appendix~\ref{app:extra-operator-algebra}.

\subsection{Modular theory in brief}
For a von Neumann algebra $\mM$ with a faithful normal state $\omega$, the GNS representation has a cyclic and separating vector $\ket{\Omega}$ and the Hilbert space $\mH_\mM=\overline{\mM\ket{\Omega}}$. Linear superoperator $\Phi:\mM\to \mM$ are represented as operators on the GNS Hilbert space:
\begin{align}
    F a\ket{\Omega}=\Phi(a)\ket{\Omega},\qquad a\in\mM
\end{align}
The Tomita super-operator $\mathcal{S}(a)=a^*$ is represented by the Tomita operator $S_\omega a\ket{\Omega}=a^*\ket{\Omega}$ with polar decomposition $S_\omega=J_\omega\Delta_\omega^
{1/2}$; $\Delta_\omega$ its modular operator and $J_\omega$. The modular Hamiltonian is $K_\omega=-\log\Delta_\omega$ which generates the modular flow $\Delta^{is}$ on the GNS Hilbert space corresponding to the superoperator (modular automorphism)
\begin{align}
\sigma_t^\omega(a)=\Delta_\omega^{it}a\Delta_\omega^{-it},\qquad a\in\mM
\end{align}
Modular conjugation maps the algebra to its commutant:
\begin{align}
&j(a)=J_\omega a J_\omega\in \mM',\qquad a\in\mM
\end{align}
 Modular flow is the unique automorphism of the algebra $\mM$ that satisfies the KMS condition. In a finite matrix algebra, it reduces to conjugation by $\rho_\omega^{it}$, whereas local QFT algebras are typically type III$_1$ factors for which no intrinsic density matrix is available.

The relative entropy of normal states is defined using relative modular operators, and decreases upon restriction to a subalgebra~\cite{Araki1976,Araki1977}. This monotonicity is the information-theoretic input for the sufficiency decomposition below. Readers familiar with these notions can go directly to the inclusion theory in the next subsection. We have retained the detailed definitions, the types of factors, and the discussion of which dynamics can be modular in Appendix~\ref{app:extra-operator-algebra}.

\subsection{General quantum subsystems (subfactors)}
\label{subsec:quantum-coarse}

Subfactors play a central role in modular ergodic theory. The modular flow of a faithful state $\omega$ on $\mM$ is a closed, automorphic dynamics of $\mM$, but from the point of view of the retained subalgebra $\mN$ it is an open dynamics: the flow can transport $\mN$ outside itself and move it around $\mM$. In tensor-product or conormal examples\footnote{An inclusion is conormal when the subalgebra $\mN$ and the relative commutant $\mR$ generate the larger algebra $\mM$.} the relative commutant $\mR$ plays the role of the environment inside $\mM$. However, for a general non-conormal inclusion, such as singular inclusions (trivial relative commutant), the relative commutant does not contain all discarded information. In this paper, we often use the intuitive language of mixing the degrees of freedom in the subalgebra $\mN$ and the relative commutant $\mR$, which strictly speaking holds only in conormal inclusions. Note that the inclusion of algebras of QFT, associated to the inclusion of regions $A=BR$ where the boundaries of regions $B$ and $R$ touch, is believed to be conormal. In QFT literature, this is often called the strong additivity condition \cite{Wiesbrock1993b,Guido1998}.

In this subsection, we review the basic notions of subfactors and conditional expectations in general von Neumanna algebras. We introduce the notion of subsystem flow, which is a key ingredient in our study of modular ergodic theory.


\paragraph{Conditional expectations:} An inclusion $\mN\subseteq\mM$ specifies the retained observables.
When it exists, a \emph{normal conditional expectation}
$\mE_{\mN}:\mM\to\mN$ is a completely positive, unital projection
with the bimodule property
\begin{equation}
 \mE_{\mN}(bac)=b\,\mE_{\mN}(a)\,c,
 \qquad a\in\mM,\quad b,c\in\mN;
 \label{eq:CE-bimodular}
\end{equation}
it is \emph{faithful} if $\mE_{\mN}(a^*a)=0$ implies $a=0$. A conditional expectation specifies how a state on $\mN$ is extended to $\mM$ \footnote{Unlike the classical conditional expectation in Appendix~\ref{sec:classical}, which exists for every coarse algebra,
a normal expectation onto an infinite-dimensional subfactor need not exist, as we show below.}. The two results below establish that when conditional expectations exist (for conormal inclusions), there is a tensor product decomposition $\mM=\mN\otimes \mR$ and there are uncorrelated states of the form $\omega=\omega_\mN\otimes\omega_\mR$.

\begin{lemma}\label{lem:expectation-slice}
For subfactors $\mN\subset \mM$, there is a one-to-one correspondence between normal faithful conditional expectation $\mE:\mM\to\mN$ (if they exists), and faithful normal states $\chi_{\mR}$ on
the relative commutant $\mR$ such that
\begin{equation}
 \mE(c)=\chi_{\mR}(c)\mathbbm 1,
 \qquad \mE(b c)=\chi_{\mR}(c)b,
 \qquad b\in\mN,\quad c\in\mR.
\end{equation}
\end{lemma}
The proof is given in Appendix~\ref{app:expectation-proofs}.

\begin{lemma}[conormality and tensor product]\label{prop:conormal-factorization}
Under the hypotheses of Lemma~\ref{lem:expectation-slice}, multiplication
extends to a normal $*$-isomorphism
\begin{equation}
 m:\mN\overline\otimes\mR
 \overset{\cong}{\longrightarrow}\mN\vee\mR,
 \qquad m(b\otimes c)=b c.
\end{equation}
If the inclusion is conormal, then
\begin{equation}\label{eq:conormal-factorization}
\mM\cong
 \mN\overline\otimes\mR .
\end{equation}
\end{lemma}
The proof is given in Appendix~\ref{app:expectation-proofs}.

\paragraph{Subsystem flow:} Let $\ket{\Omega}$ be the cyclic and separating vector corresponding to the state $\omega$ on $\mM$ in the GNS Hilbert space $\mH_\mM$. Consider the subspace $\mH_\mN=\overline{\mN\ket{\Omega}}=e_\mN\mH_\mM$ and the orthogonal projection $e_\mN$ onto it. In operator algebraic terminology, $e_\mN$ is often called the Jones projection. The restriction of the state $\omega$ to $\mN$ is denoted by $\omega_\mN$. The modular operator $\Delta_{\omega;\mM}$ is defined on $\mH_\mM$, and the modular operator $\Delta_{\omega;\mN}$ is defined on $\mH_\mN$. 

Unital CP maps of $\mM$ define a contraction on the GNS Hilbert space $\mH_\mM$. However, not every contraction on $\mH_\mM$ comes from a unital CP map on $\mM$. The projection $e_\mN$ is a contraction on $\mH_\mM$, and defines the superoperator
\begin{align}
\mE_\mN(a)=e_\mN a e_\mN
\end{align}
but it is not a unital CP map on $\mM$. In general, $\mE(a)$ is not in $\mM$, and the map is not unital. Instead, the projection $e_\mN$ defines a unital CP map on the commutant $\mM'$ \cite  {AccardiCecchini1982}:
\begin{align}
\mE'_{\mN}(a')=e_\mN a' e_\mN:\mM'\to \mN',\qquad a'\in\mM'
\end{align}
Therefore, we can always define a unital CP map called the Petz recovery map \cite{Petz1986,Petz1988} from $\mM$ to $\mN$ as 
\begin{align}
&\gamma_\mN(a)=j_\mN\circ \mE'_{\mN}\circ j_\mM(a):\mM\to \mN\nn\\
&\gamma_\mN(a)\ket{\Omega}=J_\mN e_\mN J_\mM a\ket{\Omega}=e_\mN J_\mN J_\mM a\ket{\Omega}\in e_\mN\mH_\mM,\qquad a\in\mM
\end{align}
where we have used $[J_\mN,e_\mN]=0$.

Write $\Delta_{\mM}=\Delta_{\omega;\mM}$, $\Delta_{\mN}=\Delta_{\omega;\mN}$. The modular operator $\Delta_{\mN}$ is defined on $\mH_{\mN}$, and the modular operator $\Delta_{\mM}$ is defined on $\mH_{\mM}$. The embedding map $V$ intertwines the two modular operators in a nontrivial way, which we will study in detail. They correspond to the following superoperators
\begin{align}
\Delta_{\mN}^{it}:\mH_\mN\to\mH_\mN
\end{align}


\begin{definition}[Subsystem flow]\label{def:subsystem-flow}
For an inclusion $\mN\subset\mM$ with reference vector $\Omega$, define the subsystem flows
\begin{align}
 D_{\mN\subset\mM}(t)&:=\Delta_{\mM}^{-it}e_\mN\Delta_{\mN}^{it}
 =\Delta_{\mM}^{-it}\Delta_{\mN}^{it}e_\mN
 :\mH_{\mN}\longrightarrow\mH_{\mM},\qquad t\in\mathbb R.
 \label{eq:subsystem-flow-general}
\end{align}
\end{definition}
For real $t$, $D(t)$ is an isometry fixing $\ket{\Omega}$. If $\Omega$ is also cyclic for $\mN$ on $\mH_{\mM}$, then $e_\mN=1$ and $D(t)$ is unitary. It follows from the monotonicity of modular operator that $\Delta_{\mN}\;\geq\;e_\mN\Delta_{\mM}e_\mN$, therefore, analytically continuing $t\to i\alpha$ gives a contraction for $0\le\alpha\le1/2$:
\begin{align}
 &|D_{\mN\subset\mM}(i\alpha)|\le1\nn\\
 &D_{\mN\subset\mM}(t+i/2)=J_\mM D_{\mN\subset\mM}(t)J_\mN:\mH_{\mN}\longrightarrow\mH_{\mM},\qquad t\in\mathbb R.
 \label{eq:subsystem-flow-contraction}
\end{align}
At $\operatorname{Im}z=1/2$ is the contraction dual to the Petz recovery map:
\begin{align}
\braket{a\Omega|D_{\mN\subset\mM}(i/2)b\Omega}=\braket{\gamma_\mN(a)\Omega|b\Omega},\qquad a\in\mM,\ b\in\mN.
\end{align}
The subsystem flow is generating operator for
\begin{align}
-i\partial_t D_{\mN\subset\mM}(t)\big|_{t=0}
 =(K_\mM-K_\mN)e_\mN\geq 0
 \label{eq:subsystem-flow-generator}
\end{align}
and the generating function for relative entropies
\begin{align}
-i\partial_t\braket{a\Omega|D_{\mN\subset\mM}(t)a\Omega}
 =&S_\mM(e_\mN a\Omega\|\Omega)-S_\mN(e_\mN a\Omega\|\Omega)\nn\\
 &+S_{\mN'}(e_\mN a\Omega\|\Omega)-S_{\mM'}(e_\mN a\Omega\|\Omega)\\
 &+\braket{(1-e_\mN)a\Omega|K_\mM e_\mN a\Omega}\nn
\end{align}

We will see in Section~\ref{subsec:analytic-subsystem-flow} that subsystem flow can be analytically continued to the strip $0\le\operatorname{Im}z\le1/2$
The analytic continuation is a key ingredient in this work.


As we stated earlier, infinite-dimensional von Neumann algebras allow for more exotic inclusions of factors than those in finite-dimensional matrix algebras. Here, we review a few of these inclusions that will be relevant for our study of modular ergodic theory.

\paragraph{Common cyclic and separating vector:} For an inclusion we
call $\ket{\Omega}$ a \emph{common} cyclic and separating vector if $e_N=1$. These inclusions can occur in algebras of all types, however they are more commonplace in type III algebras \footnote{Every $\sigma$-finite type $III$ factor admits a standard representation~\cite{Blackadar2006}, and under the usual separability hypotheses the set of vectors that are simultaneously cyclic and separating for finitely many standardly represented algebras is a dense $G_\delta$ set~\cite{dixmiermarechal}}. In QFT, the vacuum is cyclic and separating for the local algebras by the Reeh--Schlieder theorem.

\paragraph{Standard and split inclusions:}
An inclusion $\mN\subset\mM$ is
\emph{standard} with respect to $\Omega\in\mathcal H$ if $\Omega$ is cyclic
and separating for $\mM$, $\mN$ and the relative commutant $\mR=\mN'\cap \mM$. It is \emph{split} if a type I factor $\mathcal F$ exists with
$\mN\subset\mathcal F\subset\mM$.

\begin{theorem}[Split inclusions and normal expectations]\label{thm:split-no-expectation}
If $\mN$ is a type III factor and
$\mN\subset\mM$ is split, then no normal conditional expectation
$\mE:\mM\to\mN$ exists.
\end{theorem}
See Appendix~\ref{app:expectation-proofs} for a proof. The split property is common in QFT, and the theorem shows that normal conditional expectations are not always available.


\section{Invariant subsystems as QEC codes}\label{sec:QEC}

Theory of Quantum Error Correction (QEC) is often formulated in the Schrödinger picture for finite-dimensional quantum systems. The noise map is a Completely Positive (CP) trace-preserving map acting on density matrices. For concreteness, consider a system $A=BR$ where $R$ plays the role of the environment that purifies the noise quantum channel on states of $B$ (Stinespring dilation of the channel):
\begin{align}
    \forall \rho_B:\qquad \Phi(\rho_B)=\text{tr}_R(U_A (\rho_B\otimes \sigma_R)U_A^\dagger)
\end{align}
We say this code subspace is correctible if there exists a recovery quantum channel $\mR$ satisfying the Knill-Laflamme condition
\begin{align}
    \mR\circ \Phi\circ P_C\sim P_C
\end{align}
In a subsystem code, we cannot protect all of $P_C$, instead $P_C\mH_B$ 


In the most general case, assume that the Hilbert space of $B$ decomposes as $\mH_B=\oplus_r \mH_{B_1}^{(r)}\otimes \mH_{B_2}^{(r)}$, and the code subspace is spanned by orthonormal vectors $\ket{\Psi_i^{(r)}}_B=\ket{i}^{(r)}_{B_1}\otimes \ket{0}_{B_2}^{(r)}$ for some range of $r$. Then, this code subspace forms a mixed subsystem-subspace code that satisfies the correction equation
\begin{align}
    \mR\circ \Phi\circ P_C=\oplus_r(1\otimes \mathcal{F}_{B_2}^{(r)}
\end{align}


Assume that there exists an orthonormal set of code states $\ket{\Psi_i}\in \mH_B$ that span a $d$-dimensional code subspace of $P_C\mH_B$:
\begin{align}
    P_C=\sum_{i=1}^d \ket{\Psi_i}\bra{\Psi_i}
\end{align}
We say this code subspace is correctible if there exists a recovery quantum channel $\mR$ satisfying 
\begin{align}
    \forall \rho_B\in B(\mH_B), \rho_R\in B(\mH_R):\qquad (\mR\circ \Phi)(\rho_B\otimes \rho_R)=\rho_B\otimes \mathcal{F}(\rho_R)
\end{align}
for some quantum channel $\mathcal{F}$ on $R$.



In the Heisenberg picture, the noise map is described by the trace-dual  
\begin{align}
    \tr(\Phi(\rho) a)=\tr(\rho \Phi^*(a))=
\end{align}
which is a unital CP map. 


As we saw, type III algebras neither have density matrices nor a trace. Therefore, to formulate QEC, we need to work in the Heisenberg picture (a unital CP noise map) which is realized as a contraction map $\Phi:\mM\to \mM$ on the GNS Hilbert space of physical operators $\mH_\mM=\overline{\mA\ket{\rho^{1/2}}}$:
\begin{align}
    \Phi(a)\ket{\rho^{1/2}}=F a \ket{\rho^{1/2}}
\end{align}
In general, we are only interested in a code subspace of encoded states 
\begin{align}
    \mH_\mN=\overline{\mN\ket{\rho^{1/2}}}=e_\mN\mH_\mM
\end{align}
This code subspace is the range of an isometric encoding of logical operators 
\begin{align}
    e_\mN=VV^\dagger, \qquad \mN=\iota_V(b)=V b V^\dagger\ .
\end{align}
For instance, in finite quantum systems this encoding can be written as 
\begin{align}
    \ket{\Psi_i}_A=V\ket{i}_B=U_A(1_B\otimes \ket{0}_R)\ket{i}_B
\end{align}
and the projection to the code subspace is
\begin{align}
    e_\mN=\sum_i \ket{\Psi_i}\bra{\Psi_i}_A\ .
\end{align}




When subsystem flow acts trivially on the subalgebra $\mN$: $D_{\mN\subset\mM}(t)=e_\mN$, we say that subsytem $\mN$ is invarinat under the modular flow of $\mM$ (in operator algebraic terminology \emph{modular autonomous}):
\begin{align}
    \sigma_t^\omega(\mN)=\mN,\qquad \forall t\in\mathbb R.
\end{align}
This means that 
\begin{align}
\Delta_\mM^{it}e_\mN=e_\mN \Delta_\mN^{it},\qquad \forall t\in\mathbb R, \qquad J_\mM e_\mN=e_\mN J_\mN
\end{align}
This has a Quantum Error Correction (QEC) interpretation: the codespace is $e_\mN\mH_\mM$, the logical algebra is $\mN$, and the physical algebra is $\mM$. The restriction channel $\mE_\mN:\mM\to\mN$ is exactly correctable on the code of $\mE_\mN$ using the , and the Petz recovery map $\gamma_\mN:\mM\to\mN$ is a universal recovery channel for general inclusions.

the restriction channel $\mE_\mN:\mM\to\mN$ is exactly correctable on the code of $\mE_\mN$, and the Petz recovery map $\gamma_\mN:\mM\to\mN$ is a universal recovery channel for general inclusions.

In this section we show that modular autonomy is equivalent to exact quantum error correction of the restriction channel. We also show that the Petz recovery map is a canonical conditional expectation for modular autonomous inclusions, and that it is a universal recovery channel for general inclusions.

We now relate modular autonomy to exact correction of the quantum channel that restricts states $\rho_\mM$ down to $\rho_\mN$. A matrix example introduces the interpretation, and Theorem~\ref{thm:correctable} collects the modular, conditional-expectation and Knill--Laflamme formulations~\cite{Faulkner2020,FaulknerSperanza2024}.

\subsection{QEC and subsystem flow in matrix algebras}
A finite composite system has observable algebra $\mM=\mN\otimes\mR$: in this subsection $\mN=B(\mathcal H_B)$ and $\mR=B(\mathcal H_R)$ denote the single-system algebras, and $\mN_0$, $\mR_0$ their embeddings in $\mM$. A general unitary flow on $\mM$ sends $\mN_0\equiv \mN\otimes 1_R$ and $\mR_0\equiv 1_B\otimes \mR$ to a new embedding
\begin{align}
&\iota_t(\mN)\equiv U_t(\mN\otimes 1_R)U_t^\dagger\subset\mM  \nn\\
&\iota_t^c(\mR)\equiv U_t(1\otimes \mR)U_t^\dagger
\end{align}
where $\iota_t^c$ is the corresponding embedding of the complementary factor. It results in new decompositions $\mM=\iota_t(\mN)\vee\iota_t^c(\mR)$ with $[\iota_t(\mN),\iota_t^c(\mR)]=0$.
Writing $U_t=e^{itH}$, every such flow is the modular flow of the faithful state $\rho_{BR}:=e^{H}/\tr(e^{H})$, since $\rho_{BR}^{it}=U_t$ up to a phase (equivalently, $U_{-t}$ is the modular flow of the Gibbs state $e^{-H}/\tr(e^{-H})$):
\begin{align}
    \iota_t(\mN)=\rho_{BR}^{it}\mN_0\rho_{BR}^{-it}\ .
\end{align}
For this dynamics to preserve $\mN$ it must be that $\iota_t(\mN)=\mN_0$ for all time, which occurs if and only if the state $\rho_{BR}$ is uncorrelated between $\mN$ and $\mR$:
\begin{align}
 &   \rho_{BR}=\rho_B\otimes \rho_R:\qquad     \rho_{BR}^{it}(b\otimes 1)\rho_{BR}^{-it}=\rho_B^{it}b\rho_B^{-it}\otimes 1_R,
\end{align}
\NL{The slice expectation of \eqref{eq:slice-expectation-intro}, with hidden state $\rho_R$, preserves this product state.}

We also have
\begin{align}\label{KnillLaff}
    \mE_{\rho_R}(1_B\otimes r)=\mE_{\rho_R}(1_B\otimes 1_R)\tr(\rho_R r)=(1_B\otimes 1_R)\tr(\rho_R r)
\end{align}
The condition for autonomous evolution of subalgebra $\mN_0$ identifies an uncorrelated state $\rho_{BR}=\rho_B\otimes \rho_R$. In the language of QEC, this is the condition that the erasure of $\mR_0$ is exactly correctable on the code of $\mE_{\rho_R}$, whose logical algebra is $\mE_{\rho_R}(\mM)=\mN_0$, and the equation \eqref{KnillLaff} is the algebraic version of the Knill-Laflamme condition for subsystem QEC codes. It also implies
\begin{align}
    \mE_{\rho_R}(\rho_{BR}^{it}(1\otimes r)\rho_{BR}^{-it})=\mE_{\rho_R}(1\otimes r)=(1_B\otimes 1_R)\tr_R(\rho_R r)
\end{align}
which is a totally forgetful map, erasing all the information in $R$.

The leakage of information from $\mN_0$ into $\mR_0=1_B\otimes \mR$ can be attributed to correlations in the state $\rho_{BR}$ between $B$ and $R$. The information--disturbance theorem of QEC says that the encoding $(\mE_{\rho_R})_*$ is a code for $\mN_0$ if and only if its complementary channel, the restriction to the relative commutant $\mR_0$, is forgetful on the code (Section~\ref{subsec:QEC-KL}).

\NL{For a fixed hidden state $\chi_R$, membership in the product-state code also requires the hidden marginal to equal $\chi_R$:}

\begin{equation}\label{eq:quantum-coarse-entropy-decomposition}
 S(\rho_{BR}\|\rho_B\otimes\chi_R)=I_\rho(B:R)+S(\rho_R\|\chi_R).
\end{equation}
\NL{The two terms separate correlations between the logical and hidden factors, and a hidden marginal different from the reference $\chi_R$.}

\subsection{Modular invariant subsystem as exact correctability}\label{subsec:QEC-takesaki}

Here, we generalize the discussion of matrix algebras above to arbitrary von Neumann algebras.

\paragraph{The erasure channel.}
In the Schr\"odinger picture, restriction of normal states
\begin{equation}
 \mathcal T:\mM_*\to\mN_*,\qquad \varphi\mapsto\varphi|_{\mN},
 \label{eq:erasure-channel}
\end{equation}
is the \emph{erasure channel} associated with the inclusion: an observer confined to $\mN$ has lost access to every observable not in $\mN$. In the Heisenberg picture it is the inclusion map $\mN\hookrightarrow\mM$. When $\mM=\mN\vee\mR$ and multiplication gives a spatial tensor factorization $\mM\cong\mN\overline\otimes\mR$, this is the ordinary erasure of the factor $\mR$ and $\varphi\mapsto\varphi|_{\mR}$ is its complementary channel. Without such a factorization, the relative commutant need not account for every discarded observable, and the erasure is simply the restriction \eqref{eq:erasure-channel}. Section~\ref{subsec:QEC-KL} identifies the complementary channel for a general inclusion and states correctability in Knill--Laflamme form.

\paragraph{Encoding by a conditional expectation.}
A normal conditional expectation $\mE:\mM\to\mN$ specifies how a state on $\mN$ is extended to $\mM$:
\begin{equation}
 \mE_*:\mN_*\to\mM_*,\qquad \varphi_{\mN}\mapsto\varphi_{\mN}\circ\mE .
 \label{eq:encoding}
\end{equation}
Since $\mE$ is a projection onto $\mN$, $(\varphi_{\mN}\circ\mE)|_{\mN}=\varphi_{\mN}$, so $\mathcal T\circ\mE_*=\operatorname{id}_{\mN_*}$. Thus $\mE_*$ is an \emph{encoding} of the logical states $\mN_*$ into $\mM_*$ that is exactly correctable against the erasure \eqref{eq:erasure-channel}, with recovery map $\mE_*$ itself:
\begin{equation}
 \mathcal C_{\mE}:=\mE_*(\mN_*)=\{\varphi\in\mM_*:\ \varphi\circ\mE=\varphi\}
 \label{eq:code}
\end{equation}
is the \emph{code} of $\mE$. It is a norm-closed convex set of normal states, and the logical algebra is $\mN$ itself: $\mE(n)=n$ for $n\in\mN$. In the Heisenberg picture $\mE(n)=n$ says that the logical observables themselves are untouched; this much is automatic for a restriction, and the content of correctability is that the state on all of $\mM$, correlations included, is reconstructed from its restriction (Section~\ref{subsec:QEC-KL}). Which subalgebras of $\mM$ admit such an expectation is not a kinematic question; by Takesaki's theorem it is the dynamical question of modular autonomy, and it is this that distinguishes the quantum from the classical theory, where every subalgebra is expected with respect to the reference measure (Appendix~\ref{sec:classical}).

By Lemma~\ref{lem:expectation-slice}, every faithful normal expectation onto a factor is determined on $\mN\vee\mR$ by a faithful normal state $\chi$ on $\mR$ through $\mE(nr)=\chi(r)\,n$. In the tensor-product case the code \eqref{eq:code} is the set of product states $\varphi_{\mN}\otimes\chi$ with the hidden factor fixed; different choices of $\chi$ give different codes with the same logical algebra. 

Note that in passive QEC (noiseless subsystems) the errors, here the elements of $\mR$, commute with the logical algebra, so no active recovery is needed: the encoding $\mE_*$ itself serves as the recovery map.

\paragraph{State preservation and zero coarse-graining entropy.}
A state $\omega$ lies in the code of $\mE$ precisely when
\begin{equation}
 S(\omega\|\omega\circ\mE)=0\quad\Longleftrightarrow\quad\omega=\omega\circ\mE .
 \label{eq:QEC-entropy}
\end{equation}
This is equality of a state with its specified extension. It does not by itself establish a code: exact error correction requires a channel, a recovery map, and a family of states on which their composition is the identity. The family here is \eqref{eq:code}, and \eqref{eq:QEC-entropy} says that $\omega$ belongs to it.

\NL{The compressed flow $C(t)$ of Definition~\ref{def:subsystem-flow} gives a Hilbert-space criterion for exact correctability. Its generator and the compression identity are recorded in Appendix~\ref{app:compressed-flow}.}

The following theorem collects the equivalent forms of the correctable case.
\begin{theorem}[Correctable inclusions]\label{thm:correctable}\label{prop:KL}
Let $\mN\subset\mM$ be an inclusion of factors, $\Omega$ cyclic and separating for $\mM$ and separating for $\mN$, and $\omega=\braket{\Omega|\cdot\,\Omega}$. The following are equivalent.
\begin{enumerate}
 \item[\rm(i)] \emph{Modular autonomy:} $\sigma^{\omega}_t(\mN)=\mN$ for all $t\in\mathbb R$.
 \item[\rm(ii)] \emph{Expected inclusion:} there is a normal conditional expectation $\mE:\mM\to\mN$ with $\omega\circ\mE=\omega$.
 \item[\rm(iii)] \emph{Trivial subsystem flow:} $C_{\mN\subset\mM}(t)=1_{\mH_{\mN}}$ for all $t\in\mathbb R$.
\item[\rm(iv)] \NL{\emph{Knill--Laflamme compression:} $e_{\mN}\mM e_{\mN}\subseteq\mN e_{\mN}$.}
\end{enumerate}
When these hold, $\mE$ is unique, its GNS projection is $e_{\mN}=VV^*$, and $\mathcal T\circ\mE_*=\operatorname{id}$ on $\mN_*$ while $\mE_*\circ\mathcal T=\operatorname{id}$ on the code $\mathcal C_{\mE}\ni\omega$.

\NL{For every $x\in\mM$, the expectation is characterized by $e_{\mN}xe_{\mN}=\mE(x)e_{\mN}$.}
\end{theorem}

\begin{proof}
(i)$\Leftrightarrow$(ii) is Takesaki's theorem~\cite{Takesaki1972}; uniqueness and the identification of the GNS projection with $e_{\mN}$ are part of the same theorem, since $\mE(m)\Omega=e_{\mN}m\Omega$ for an $\omega$-preserving expectation.

(ii)$\Rightarrow$(iii). For an $\omega$-preserving expectation the modular flows intertwine
\begingroup\color{blue}
\begin{align}
 \sigma^{\omega;\mM}_t(n)&=\sigma^{\omega;\mN}_t(n),\qquad n\in\mN,\nn\\
 \Delta_{\mM}^{it}V&=V\Delta_{\mN}^{it}.
 \label{eq:takesaki-intertwining}
\end{align}
\endgroup Hence
$C_{\mN\subset\mM}(t)=V^*V\Delta_{\mN}^{-it}\Delta_{\mN}^{it}=1$.

(iii)$\Rightarrow$(i). If $C(t)=1$ then $V^*\Delta_{\mM}^{-it}V=\Delta_{\mN}^{-it}$ is unitary on $\mH_{\mN}$, so for $\psi\in\mH_{\mN}$,
$\|e_{\mN}\Delta_{\mM}^{-it}V\psi\|=\|V\psi\|=\|\Delta_{\mM}^{-it}V\psi\|$,
which forces $\Delta_{\mM}^{-it}V\psi\in\operatorname{Ran}V$. Thus $\Delta_{\mM}^{-it}$ leaves $\mH_{\mN}$ invariant and restricts there to $V\Delta_{\mN}^{-it}V^*$. For $n\in\mN$,
\[
 \sigma^{\omega}_{-t}(n)\Omega=\Delta_{\mM}^{-it}n\Omega
 =V\Delta_{\mN}^{-it}n\Omega=\sigma^{\omega;\mN}_{-t}(n)\Omega .
\]
Since $\Omega$ is separating for $\mM$, $\sigma^{\omega}_{-t}(n)=\sigma^{\omega;\mN}_{-t}(n)\in\mN$ for all $t$, which is (i).

The final statements follow from $\mE_*$ being a right inverse of $\mathcal T$ on $\mN_*$ and the identity $\varphi\circ\mE=\varphi$ on the code.

\NL{(ii)$\Rightarrow$(iv).} For $n,n'\in\mN$ and $x\in\mM$,
\[
 \langle n\Omega|\,x\,n'\Omega\rangle=\omega(n^*xn')=\omega\bigl(\mE(n^*xn')\bigr)=\omega\bigl(n^*\mE(x)n'\bigr)=\langle n\Omega|\,\mE(x)\,n'\Omega\rangle,
\]
by invariance of $\omega$ and the bimodule property. Since $\mE(x)\in\mN$ commutes with $e_{\mN}\in\mN'$, this is $e_{\mN}xe_{\mN}=\mE(x)e_{\mN}$.

\NL{(iv)$\Rightarrow$(ii).} Define $\mE(x)\in\mN$ by $e_{\mN}xe_{\mN}=\mE(x)e_{\mN}$. It is well defined: the central support of $e_{\mN}$ in $\mN'$ is the projection onto $\overline{\mN'\mN\Omega}\supseteq\overline{\mN'\Omega}=\mathcal H$, since $\Omega$ is separating for $\mN$ and hence cyclic for $\mN'$, so $n\mapsto ne_{\mN}$ is a normal $*$-isomorphism of $\mN$ onto $\mN e_{\mN}$. Thus $\mE$ is the composition of the normal completely positive map $x\mapsto e_{\mN}xe_{\mN}$ with the inverse of this isomorphism: it is linear, normal, completely positive and unital. It is $\mN$-bimodular because $e_{\mN}\in\mN'$ gives $e_{\mN}\,nxn'\,e_{\mN}=n\,(e_{\mN}xe_{\mN})\,n'$, and it preserves $\omega$ because $\omega(\mE(x))=\langle\Omega|\mE(x)e_{\mN}\Omega\rangle=\langle\Omega|e_{\mN}xe_{\mN}\Omega\rangle=\omega(x)$.
\end{proof}

Two remarks limit what Theorem~\ref{thm:correctable} asserts. First, it concerns the specific channel \eqref{eq:erasure-channel} and the specific family \eqref{eq:code}; failure of (ii) for the state $\omega$ does not exclude expectations preserving another faithful normal state, nor error correction for a different channel or family of states. Second, neither the Jones projection $e_{\mN}$ nor a product of modular conjugations is, by itself, a recovery channel: recovery is implemented by $\mE_*$, which exists only under the hypotheses of the theorem.

\paragraph{\NL{Knill--Laflamme conditions and the forgetful complement.}}
\label{subsec:QEC-KL}
\NL{The finite-dimensional channel and complementary-channel dictionary is recalled in Appendix~\ref{app:finite-KL}. The restriction channel has the following modular interpretation.}

The system is $\mM$ and the logical algebra is $\mN$. The errors are the elements of the relative commutant $\mR=\mN'\cap\mM$: in the tensor-product case $\mM=\mN\overline\otimes\mR$, erasure of the factor $\mR$ has Kraus operators $\bra{k}$ on the erased factor, and the operators $E_l^\dagger E_k=\ket{k}\!\bra{l}$ of \eqref{eq:KL} are the matrix units of $\mR$. The complementary channel of the erasure on a vector state $\psi$ is restriction to the commutant $\mN'$, the Stinespring environment of the inclusion on the GNS space; it contains $\mR$, the part of the environment inside $\mM$, and $\mM'$, the reference that purifies $\mM$. The code subspace is $\mH_{\mN}=\overline{\mN\Omega}$, with the Jones projection $e_{\mN}=VV^*$ as code projection. 

\NL{The compression identity in Theorem~\ref{thm:correctable} is \eqref{eq:KL} with an operator-valued coefficient.} For the errors themselves, $c\in\mR$, the coefficient $\mE(c)=\chi(c)$ is a scalar by Lemma~\ref{lem:expectation-slice}, and the condition is the Knill--Laflamme condition in its original form. For a general $x\in\mM$ the coefficient is the logical operator $\mE(x)$, which expresses the Schr\"odinger-picture requirement of Section~\ref{subsec:QEC-takesaki} that the whole state on $\mM$, and not only the expectation values of the logical algebra, be reconstructed from the restriction to $\mN$. The Heisenberg-picture requirement that the expectation values of logical observables survive is automatic for the restriction to $\mN$; the content of exact correctability here lies in the correlations between $\mN$ and the discarded observables, and this is what the operator-valued coefficient records.

\begin{corollary}[Conormal inclusions: Knill--Laflamme is a forgetful complement]\label{cor:KL-conormal}
Assume in addition that the inclusion is conormal, $\mN\vee\mR=\mM$. Then \NL{{\rm(i)--(iv)}} are further equivalent to each of the following, for a faithful normal state $\chi$ on $\mR$.
\begin{enumerate}
 \item[\rm(a)] \emph{Knill--Laflamme for the gauge algebra:} $e_{\mN}\,c\,e_{\mN}=\chi(c)\,e_{\mN}$ for all $c\in\mR$.
 \item[\rm(b)] \emph{Forgetful complement:} $\langle\psi|c\,\psi\rangle=\chi(c)$ for every unit vector $\psi\in\mH_{\mN}$ and every $c\in\mR$; the complementary channel $\psi\mapsto\psi|_{\mR}$ is constant on the code subspace.
 \item[\rm(c)] \emph{No logical--gauge correlations in the reference:} $\omega(nc)=\omega(n)\chi(c)$ for all $n\in\mN$, $c\in\mR$.
\end{enumerate}
\end{corollary}

\begin{proof}
(a)$\Leftrightarrow$(b) is polarization. (b)$\Leftrightarrow$(c): for $n,n'\in\mN$ and $c\in\mR$, $\langle n\Omega|c\,n'\Omega\rangle=\omega(n^*cn')=\omega(c\,n^*n')$ because $c$ commutes with $\mN$, so (b) reads $\omega(cn'')=\chi(c)\omega(n'')$ for all $n''\in\mN$, with $\chi=\omega|_{\mR}$, faithful since $\omega$ is. (i)$\Rightarrow$(c) is Lemma~\ref{lem:expectation-slice}, applied to the expectation of Theorem~\ref{thm:correctable}, which is faithful because $\omega$ is. (c)$\Rightarrow$(i): $\mR$ is a factor, since an element of its centre commutes with $\mN$ and with $\mR$, hence with $\mM$. Thus $\omega$ is a faithful normal product state on the commuting factors $\mN$ and $\mR$, which generate $\mM$; as in the proof of Proposition~\ref{prop:conormal-factorization}, multiplication extends to a spatial isomorphism $\mN\overline\otimes\mR\cong\mM$ carrying $\omega|_{\mN}\otimes\chi$ to $\omega$~\cite{DoplicherLongo1984}, and $\mE=\operatorname{id}\otimes\chi$ is a normal conditional expectation preserving $\omega$.
\end{proof}

Read together with Theorem~\ref{thm:correctable}, the corollary says that for conormal inclusions Takesaki's condition \emph{is} the Knill--Laflamme condition for erasure of the gauge algebra, and that the forgetful complement is its environment-side reading. The intuition is the one behind the Takesaki condition of Section~\ref{sec:intro}: the modular Hamiltonian $K_{\mM}=-\log\Delta_{\mM}$ generates the dynamics for which $\omega$ is an equilibrium state; correlations between $\mN$ and $\mR$ in $\omega$ are an interaction term in $K_{\mM}$ coupling logical to gauge degrees of freedom; and such a term is precisely what transports logical observables out of $\mN$. In finite dimensions, where $\mN=\bigoplus_aM_{n_a}\otimes1_{m_a}$, all of \NL{(i)--(iv) and (a)--(c)} reduce to $\rho=\bigoplus_ap_a\,\rho_a\otimes\tau_a$, the standard form of a code state of an operator-algebra code (Section~\ref{section:basics}).

\paragraph{Non-conormal inclusions.}
Without conormality the gauge algebra $\mR$ does not account for the discarded observables, and \NL{(a)--(b)} are not sufficient for (i). For an irreducible subfactor, $\mR=\mathbb C1$ and (a), (b) hold trivially, while (i) fails for every proper inclusion with $\Omega$ cyclic for $\mN$ (Corollary~\ref{cor:common-cyclic-noncorrectable}). The environment of the erasure is then the full commutant $\mN'$, which contains $\mR\vee\mM'=(\mR'\cap\mM)'$ and equals it only when $\mN=\mR'\cap\mM$; and ``constant on the code'' cannot hold on $\mN'$, because the reference $\mM'$ varies over the code by design. Its correct replacement is: for every $\psi\in\mH_{\mN}$ the restriction $\mN'\to\mM'$ is sufficient for the pair $(\psi,\omega)$, i.e.\ the environment holds no information distinguishing a code vector from the reference beyond what the reference system $\mM'$ already holds. This follows from (i). The support of $\omega$ on $\mN'$ is $e_{\mN}$, so for $\psi\in\mH_{\mN}$ the relative entropies on $\mN'$ reduce to the compression $e_{\mN}\mN'e_{\mN}$, which is the commutant of $\mN$ on $\mH_{\mN}$; the decomposition of Proposition~\ref{prop:sufficiency} below therefore holds for the compressed generator $V^*K_{\mM}V-K_{\mN}$ of Section~\ref{subsec:QEC-takesaki} on $\mH_{\mN}$, and (i) gives $C_{\mN\subset\mM}(t)\equiv1$, hence $V^*K_{\mM}V=K_{\mN}$ on $\mH_{\mN}$ and the vanishing of both nonnegative gaps.

\subsection{Symmetries, QEC and sub-theories}\label{subsec:symmetries-QEC-subtheories}

The preceding discussion interpreted an inclusion $\mN\subset\mM$ as a
coarse-graining and related its modular autonomy to exact recovery of a
reference-preserving code. Symmetries supply an important class of such
inclusions. They also distinguish three notions that need not coincide:
the preservation of an operator subalgebra by a dynamics, the correctability
of a family of states after restriction to that algebra, and the preservation
of a conventional error-correcting code subspace. The first two are tied by
Takesaki's theorem for the modular dynamics of a faithful reference state;
the third involves an additional choice of errors and a code projector.

\paragraph{Invariant algebras and symmetry averaging.}
Let a compact group $G$ act normally and continuously on $\mM$ by
$*$-automorphisms $\alpha_g$, and let the faithful normal state $\omega$
be $G$-invariant, $\omega\circ\alpha_g=\omega$. The neutral, or
fixed-point, algebra and its Haar expectation are
\begin{align}
 \mN=\mM^G
 &=\{x\in\mM:\alpha_g(x)=x\ \text{for all }g\in G\},\nonumber\\
 \mE_G(x)&=\int_G\!\dd g\,\alpha_g(x),
 \qquad \omega\circ\mE_G=\omega .
 \label{eq:symmetry-Haar-expectation}
\end{align}
The modular automorphism group of a faithful state is unique. Therefore
\begin{align}
 \alpha_g\circ\sigma_t^{\omega;\mM}
 &=\sigma_t^{\omega;\mM}\circ\alpha_g,
 \qquad
 \sigma_t^{\omega;\mM}(\mN)=\mN,
 \qquad
 \mE_G\circ\sigma_t^{\omega;\mM}
 =\sigma_t^{\omega;\mM}\circ\mE_G .
 \label{eq:symmetry-modular-intertwining}
\end{align}
The fixed-point observables need not be modular zero modes: they can evolve
nontrivially into other elements of $\mN$. By Theorem~\ref{thm:correctable},
the expectation $\mE_G$ defines an exactly recoverable family
$\mathcal C_{\mE_G}=\{\varphi:\varphi=\varphi\circ\mE_G\}$ under
restriction to $\mN$. In other words, information encoded in these
symmetry-invariant states is retained by the neutral algebra. For an
arbitrary state outside this code, Haar averaging discards charged
coherences and is not a recovery operation.

The state-preserving hypothesis is essential. Compactness produces a
normal conditional expectation, but does not ensure
$\omega\circ\mE_G=\omega$ for a general state $\omega$. In a matrix
algebra $\mM=B(\cH)$, with $\alpha_g=\operatorname{Ad}U_g$ and faithful
density matrix $\rho$, the condition reads $[\rho,U_g]=0$ for every $g$;
then $K_{\mM}=-\log\rho$ commutes with the symmetry representation. If the
representation has the isotypic decomposition
\begin{align}
 \cH=\bigoplus_{\lambda} V_\lambda\otimes\mathbb C^{m_\lambda},
 \qquad
 \mM^G\simeq\bigoplus_{\lambda}
       \mathbf 1_{V_\lambda}\otimes B(\mathbb C^{m_\lambda}),
 \label{eq:symmetry-isotypic}
\end{align}
where $V_\lambda$ are inequivalent irreducible representations, Schur's
lemma shows that an invariant faithful density matrix has the form
$\rho=\bigoplus_\lambda\mathbf1_{V_\lambda}\otimes r_\lambda$,
with $r_\lambda>0$ and $\sum_\lambda\dim(V_\lambda)\Tr r_\lambda=1$.
Modular evolution acts within the multiplicity factors. The neutral
algebra may have a nontrivial center, so a fixed-point algebra is not
necessarily a subfactor; Takesaki's theorem applies to von Neumann
subalgebras in this more general form.

Stabilizer codes, topological codes and noiseless subsystems give concrete finite-dimensional illustrations of this distinction. The technical examples, including non-Abelian noiseless subsystems, are collected in Appendix~\ref{app:extra-qec-examples} so that the modular and QEC arguments can be read without interruption.

\paragraph{From symmetry sectors to local sub-theories.}
For a vacuum QFT with a compact unbroken internal zero-form symmetry,
let $\mathcal F(O)$ be a field algebra and $U(g)\Omega=\Omega$. Internal
symmetry preserves every region, so
\begin{align}
 \mN(O)=\mathcal F(O)^G,\qquad
 \mE_O(x)=\int_G\dd g\,\operatorname{Ad}U(g)(x)
 \label{eq:internal-neutral-net}
\end{align}
are vacuum-preserving, modular-covariant inclusions. Moreover,
Reeh--Schlieder for the field net implies
\begin{align}
 \overline{\mN(O)\Omega}=P_G\cH,
 \qquad P_G=\int_G\dd g\,U(g),
 \label{eq:internal-coherent-projection}
\end{align}
independent of $O$. The family is therefore coherent in the sense of
Borchers' construction of local sub-theories~\cite{revolutionizing}:
its restriction to $P_G\cH$ is a Poincar\'e-covariant neutral theory.
\NL{The proper neutral algebra is not cyclic on the full field Hilbert space, consistently with Corollary~\ref{cor:common-cyclic-noncorrectable}.}

Geometric symmetries give a weaker, region-dependent example. For a
ball $B$ in the QFT vacuum, rotations preserve both $\mathcal A(B)$ and
$\Omega$. Thus $\mathcal A(B)^{SO(d-1)}$ is modular invariant in
$d$ spacetime dimensions, even when the ball modular Hamiltonian is
nonlocal. The associated singlet projector depends on the center of the
ball, however, and generally fails Borchers' coherence condition across
translated regions. Rotational singlets are therefore an autonomous
\emph{algebra for a fixed ball}, but not automatically a relativistic
sub-theory. Similarly, for disconnected regions, the diagonal neutral
algebra can contain charge--anticharge intertwiners connecting the
components, while independent charge rotations of each component need
not preserve the entangled vacuum.

Finally, a higher-form symmetry should not be identified with ordinary
symmetry averaging on every regional algebra. A one-form symmetry acts
on extended lines rather than pointlike observables; for nontrivial
topology it can produce a strict inclusion of additive and maximal
algebras. If $\Omega$ is cyclic for both the additive and maximal algebras and separating
for the maximal one, such a strict inclusion cannot admit a
vacuum-preserving expectation. Consequently its additive algebra is
\emph{not} modular invariant under the maximal algebra's vacuum flow,
even if a compact dual action or an expectation erasing winding sectors
exists~\cite{GaiottoKapustinSeibergWillett2015,CasiniHuertaMaganPontello2021}.
The distinction illustrates the point of this subsection: symmetry is a
source of exact modular QEC only when it acts on the specified algebra
and preserves the specified reference state. The companion paper
examines the corresponding geometric and topological QFT examples in
more detail~\cite{companionQFT}.

\section{Subsystem flow and ergodic subsystems}\label{sec:ergodic}


The modular flow of $\omega$ transports the subalgebra $\mN$ to
$\sigma^{\omega;\mM}_t(\mN)$. Since in general
$\sigma^{\omega;\mM}_t(\mN)$ need not be contained in $\mN$, whenever a conditional expectation $\mE:\mM\to\mN$ is available, the effective dynamics seen in $\mN$ is given by the unital completely positive map
\begin{align}
    \Phi^{\omega}_t:=\mE\circ\sigma^{\omega;\mM}_t\big|_{\mN}:\mN\to\mN\ ,
    \label{eq:reduced-dynamics}
\end{align}
which depends on the state $\omega$ and on the choice of $\mE$. The map $\Phi^{\omega}_t$ is dual to a quantum channel. Intuitively, we can view the existence of normal conditional expectations as the existence of states uncorrelated between $\mN$ and $\mR$\footnote{Strictly speaking, this characterization is exact for conormal inclusions of factors ($\mN \vee \mR = \mM$), where any faithful normal conditional expectation $\mE: \mM \to \mN$ induces states $\omega_0 = \phi_{\mN} \circ \mE$ that factorize across $\mN$ and $\mR$, $\omega_0(nr) = \phi_{\mN}(n)\psi_{\mR}(r)$. At the opposite extreme, for a singular inclusion ($\mR = \mathbb{C}1$), factorization between $\mN$ and $\mR$ holds trivially for every state, even when no conditional expectation exists. When a singular inclusion \emph{does} admit a normal conditional expectation (which is then necessarily unique), we have the linear decomposition $\mM=\mN\oplus \ker{\mE}$.  The induced state $\omega_0 = \phi_{\mN} \circ \mE$ is instead uncorrelated between $\mN$ and $\ker\mE$ in the sense that for all $n\in\mN$ and $x\in \ker\mE$ we have $\omega_0(n x)=\omega_0(n) \omega_0(x)=0$. Note that in general the kernel $\ker(\mE)$ of a faithful normal conditional expectation is never an algebra. For example, the conditional expectation onto diagonal matrices has the off-diagonal matrices in its kernel, but multiplying two off-diagonal matrices can yield a diagonal matrix.}. From this perspective, the quantum channel dual to $\Phi^{\omega}_t$ can be viewed as starting with an uncorrelated state of $\mN$ and $\mR$, coupling them via the modular flow of $\omega$, and then projecting $\mR$ out.

\subsection{Subsystem flow}
\label{subsec:analytic-subsystem-flow}

\NL{The real-time comparison of Definition~\ref{def:subsystem-flow} extends to imaginary modular time:}

\begin{equation}
    D_{\mN\subset\mM}(z)
    :=
    \overline{
    \Delta_{\mM}^{-iz}V\Delta_{\mN}^{iz}}
    :
    \mH_{\mN}\longrightarrow\mH_{\mM} .
    \label{eq:Wz-definition}
\end{equation}
We abbreviate $D_{\mN\subset\mM}(z)$ to $D(z)$ and put
$D_\alpha:=D(i\alpha)$. \NL{This is the analytic continuation of Definition~\ref{def:subsystem-flow}; $C(z)=V^*D(z)$ denotes its compression.} Modular monotonicity, reviewed in
Lemma~\ref{lem:modular-monotonicity}, gives a bounded analytic continuation
to the strip
\begin{equation}
    \mathcal S_{1/2}
    :=
    \left\{
    z\in\mathbb C:
    0<\operatorname{Im}z<\frac12
    \right\},
\end{equation}
strongly continuous on its closure. Writing
\begin{equation}
    z=t+i\alpha,
    \qquad
    0\leq\alpha\leq\frac12 ,
\end{equation}
one has
\begin{equation}
    D(t+i\alpha)
    =
    \Delta_{\mM}^{-it}
    D_\alpha
    \Delta_{\mN}^{it},
    \qquad
    D_\alpha
    :=
    \overline{
    \Delta_{\mM}^\alpha V\Delta_{\mN}^{-\alpha}},
    \label{eq:W-alpha-strip}
\end{equation}
and
\begin{equation}
    \boxed{
    \|D(z)\|\leq1,
    \qquad
    0\leq\operatorname{Im}z\leq\frac12 .
    }
    \label{eq:W-strip-bound}
\end{equation}
\NL{By Lemma~\ref{lem:modular-monotonicity}, both boundaries are isometric, with $D_0=V$ and $D_{1/2}=J_{\mM}VJ_{\mN}$.}


\NL{When $\Omega$ is common cyclic and separating, $V=1$ and $D(t)$ is Borchers\textquotesingle{} characteristic function~\cite{revolutionizing}.}

For $b\in\mN$, its adjoint action and the transported algebra are
\begin{equation}\label{eq:characteristic-algebra-transport}
 \operatorname{Ad}D_{\mN\subset\mM}(t)(b)=\sigma^{\omega;\mM}_{-t}\bigl(\sigma^{\omega;\mN}_t(b)\bigr),
 \qquad \operatorname{Ad}D_{\mN\subset\mM}(t)(\mN)=\sigma^{\omega;\mM}_{-t}(\mN).
\end{equation}
This compares the intrinsic subsystem motion with the outer modular motion.
It requires no conditional expectation. In general $D(t)$ is a cocycle rather than
a one-parameter group; it need not belong to $\mM$.

\begingroup\color{blue}
\begin{align}
 D(t+s)&=\Delta_{\mM}^{-it}D(s)\Delta_{\mM}^{it}D(t),\nn\\
 D(t+i/2)&=J_{\mM}D(t)J_{\mN},\qquad D(i/2)=J_{\mM}J_{\mN}.
 \label{eq:subsystem-flow-boundary}
\end{align}
\endgroup
\NL{The full characterization of an inclusion by its characteristic function is recalled in Appendix~\ref{app:characteristic-function}. The upper endpoint supplies the recovery maps below and the canonical tower in Section~\ref{subsec:interpolation-tunnel}.}

\subsection{Subsystem flow as a sufficiency witness}\label{subsec:QEC-sufficiency}

\NL{Throughout this subsection $\Omega$ is a common cyclic and separating vector for $\mN\subset\mM$. The generator of the subsystem flow quantifies, state by state, the failure of sufficiency under restriction.} The relevant notion is Petz's \emph{sufficiency}~\cite{Petz1986,Petz1988,OhyaPetz1993}. Let $\alpha:\mN\to\mM$ be a unital normal completely positive map (here the inclusion) and $\varphi,\omega$ faithful normal states on $\mM$. The restriction $\alpha$ is \emph{sufficient} for the pair $(\varphi,\omega)$ if there is a unital normal completely positive map $\beta:\mM\to\mN$ with $\varphi\circ\alpha\circ\beta=\varphi$ and $\omega\circ\alpha\circ\beta=\omega$; equivalently, both states are recovered from their restrictions to $\mN$ by one and the same channel.

\begin{theorem}[Petz]\label{thm:Petz}
For faithful normal states $\varphi,\omega$ on $\mM$ with $S(\varphi\|\omega)<\infty$, the restriction to $\mN$ is sufficient for $(\varphi,\omega)$ if and only if
\[
 S_{\mM}(\varphi\|\omega)=S_{\mN}(\varphi|_{\mN}\|\omega|_{\mN}) .
\]
In that case a recovery channel is given by the Petz map of the inclusion with respect to $\omega$. Equivalently, the Connes cocycle $[D\varphi:D\omega]_t=\Delta_{\varphi|\omega}^{it}\Delta_{\omega}^{-it}$ belongs to $\mN$ for every $t\in\mathbb R$~\cite{Petz1988,JencovaPetz2006}.
\end{theorem}

\paragraph{The Petz dual and generalized conditional expectations.}
Even if no $\omega$-preserving conditional expectation exists, the inclusion
has a canonical Petz dual $\mathsf E_\omega:\mM\to\mN$, a normal unital
completely positive map preserving the reference state. With the
embedding $V:\mH_{\mN}\to\mH_{\mM}$ and the modular conjugations of the
restricted and ambient states, it is
\begin{equation}
 \boxed{\quad
 \mathsf E_\omega(m)=\Gamma_c m\Gamma_c^*,\qquad
 \Gamma_c=J_{\mN}V^*J_{\mM}
 =D_{\mN\subset\mM}(i/2)^*.\quad}
 \label{eq:QEC-Petz-dual}
\end{equation}
The upper-strip value $D(i/2)$ is an isometry between the two GNS
Hilbert spaces; its adjoint implements the Heisenberg-picture Petz map,
also called the Accardi--Cecchini generalized conditional expectation
\cite{AccardiCecchini1982,Petz1988}. In general $\mathsf E_\omega$ is
not a projection onto $\mN$, and $\mathsf E_\omega(n)\ne n$ for some
$n\in\mN$. It becomes the $\omega$-preserving conditional expectation
precisely when $\mN$ is modularly autonomous, as in
Theorem~\ref{thm:correctable}.

\paragraph{Universal recovery from the upper boundary.}
The Petz dual in \eqref{eq:QEC-Petz-dual} is the $t=0$ member of a family
of rotated Petz maps. For $m\in\mM$ define
\begin{align}
 \mathsf E_{\omega,t}(m)
 &:=\sigma_{-t/2}^{\omega;\mN}
 \!\circ\mathsf E_\omega\circ
 \sigma_{t/2}^{\omega;\mM}(m)\nn\\
 &=D_{\mN\subset\mM}\!\left(\frac{t+i}{2}\right)^*
 m\,D_{\mN\subset\mM}\!\left(\frac{t+i}{2}\right),
 \qquad t\in\mathbb R.
 \label{eq:rotated-Petz-D}
\end{align}
Indeed, $D((t+i)/2)^*=\Delta_{\mN}^{-it/2}
D(i/2)^*\Delta_{\mM}^{it/2}$ by the upper-strip boundary
identity. Each $\mathsf E_{\omega,t}:\mM\to\mN$ is normal, unital
and completely positive. The universal (twirled) recovery map is the
following average in the ultraweak topology:
\begin{align}
 \mathsf E_\omega^{\mathrm{univ}}(m)
 &=\int_{-\infty}^{\infty}\!dt\,\beta_0(t)\,
 D\!\left(\frac{t+i}{2}\right)^*m
 D\!\left(\frac{t+i}{2}\right),\nn\\
 \beta_0(t)&=\frac{\pi}{2[1+\cosh(\pi t)]}.
 \label{eq:universal-recovery-D}
\end{align}
Thus universal recovery is a weighted integral of the \emph{CP maps}
induced by the upper boundary of $D(z)$, not a linear integral of
$D(t)$ itself. The Schr\"odinger-picture dual
$(\mathsf E_\omega^{\mathrm{univ}})_*:\mN_*\to\mM_*$ depends only
on the inclusion and $\omega$, recovers $\omega$ exactly, and for
normal states $\varphi$ of finite relative entropy obeys the
approximate recovery bound
\begin{equation}
 S_{\mM}(\varphi\|\omega)-S_{\mN}(\varphi\|\omega)
 \geq -2\log F_{\mM}\!\left(
 \varphi,\,\varphi|_{\mN}\circ
 \mathsf E_\omega^{\mathrm{univ}}\right),
 \label{eq:universal-recovery-bound}
\end{equation}
where $F_{\mM}$ denotes root fidelity
\cite{JungeRennerSutterWildeWinter2018,FaulknerHollandsSwingleWang2022}.
The formula holds also without common cyclicity if $D$ is read as the
inclusion isometry between the two GNS spaces. It identifies the
universal recovery channel as another use of the same analytic
subsystem flow that supplies the inclusion parent Hamiltonian below.

The cocycle condition is the form of Petz's theorem closest to the subsystem flow. The Connes cocycle compares the modular groups of two states on one algebra, $\sigma^{\varphi}_t=\operatorname{Ad}[D\varphi:D\omega]_t\circ\sigma^{\omega}_t$; the subsystem flow $D_{\mN\subset\mM}(t)=\Delta_{\mM}^{-it}\Delta_{\mN}^{it}$ compares the modular groups of one state on two algebras. Both are cocycles rather than one-parameter groups, and the generator of the second produces exactly the entropy differences that appear in Theorem~\ref{thm:Petz}. The link is the following identity, which expresses the modular Hamiltonian of a state as the imbalance between the information held by an algebra and by its commutant.

\begin{lemma}[Modular energy as an entropy imbalance]\label{lem:modular-energy}
Let $\mathcal C$ be a von Neumann algebra on $\mathcal H$, let $\Omega$ and $\Psi$ be unit vectors that are cyclic and separating for $\mathcal C$, with induced states $\omega$ and $\varphi$, and suppose that $S_{\mathcal C}(\varphi\|\omega)$ and $S_{\mathcal C'}(\varphi\|\omega)$ are finite. Then $t\mapsto\langle\Psi|\Delta_{\omega;\mathcal C}^{it}|\Psi\rangle$ is differentiable at $t=0$ and
\begin{equation}
 -i\,\partial_t\langle\Psi|\Delta_{\omega;\mathcal C}^{it}|\Psi\rangle\big|_{t=0}
 =S_{\mathcal C'}(\varphi\|\omega)-S_{\mathcal C}(\varphi\|\omega).
 \label{eq:modular-energy}
\end{equation}
If moreover $\Psi$ lies in the form domain of $\log\Delta_{\omega;\mathcal C}$, the left side is the quadratic form $\langle\Psi|\log\Delta_{\omega;\mathcal C}|\Psi\rangle$.
\end{lemma}

\begin{proof}
Write $\omega',\varphi'$ for the states induced on $\mathcal C'$. For cyclic and separating vectors $\Phi,\Psi$ the relative modular operators $\Delta_{\Phi|\Psi;\mathcal C}$ of \eqref{eq:relative-modular} satisfy~\cite{Araki1976,Araki1977}
\begin{equation}
 \Delta_{\Phi|\Psi;\mathcal C}^{it}=[D\Phi:D\Psi]^{\mathcal C}_t\,\Delta_{\Psi;\mathcal C}^{it},
 \qquad
 \Delta_{\Phi|\Psi;\mathcal C'}=\Delta_{\Psi|\Phi;\mathcal C}^{-1},
 \label{eq:relative-modular-cocycles}
\end{equation}
where $[D\Phi:D\Psi]^{\mathcal C}_t\in\mathcal C$ is the Connes cocycle of the two states on $\mathcal C$. The first identity changes the state in the first slot; applied on $\mathcal C'$ and transported by the second identity, it changes the state in the second slot: $\Delta_{\Phi|\Psi;\mathcal C}^{it}=\Delta_{\Phi;\mathcal C}^{it}\,[D\Phi':D\Psi']^{\mathcal C'}_t$. Applying both to $\Delta_{\omega;\mathcal C}=\Delta_{\omega|\omega;\mathcal C}$ gives
\begin{equation}
 \Delta_{\omega;\mathcal C}^{it}=[D\omega:D\varphi]^{\mathcal C}_t\;\Delta_{\varphi;\mathcal C}^{it}\;[D\varphi':D\omega']^{\mathcal C'}_t .
 \label{eq:three-factor}
\end{equation}
Set $a_t=[D\omega:D\varphi]^{\mathcal C}_t\in\mathcal C$ and $b_t=\Delta_{\varphi;\mathcal C}^{it}[D\varphi':D\omega']^{\mathcal C'}_t\Delta_{\varphi;\mathcal C}^{-it}\in\mathcal C'$, both unitary with $a_0=b_0=1$. Since $\Delta_{\varphi;\mathcal C}^{it}\Psi=\Psi$, \eqref{eq:three-factor} gives
\begin{align}
 \langle\Psi|\Delta_{\omega;\mathcal C}^{it}|\Psi\rangle&=\langle\Psi|a_tb_t|\Psi\rangle,\nn\\
 \langle\Psi|a_t|\Psi\rangle&=\varphi\bigl([D\omega:D\varphi]_t\bigr),
 \qquad
 \langle\Psi|b_t|\Psi\rangle=\varphi'\bigl([D\varphi':D\omega']_t\bigr).
\end{align}
By the first identity in \eqref{eq:relative-modular-cocycles}, $\varphi([D\omega:D\varphi]_t)=\langle\Psi|\Delta_{\omega|\varphi;\mathcal C}^{it}|\Psi\rangle=\int e^{it\lambda}\,d\mu(\lambda)$, with $\mu$ the spectral measure of $\log\Delta_{\omega|\varphi;\mathcal C}$ in $\Psi$. Here $\int e^{\lambda}d\mu=\|\Delta_{\omega|\varphi;\mathcal C}^{1/2}\Psi\|^2=\|\Omega\|^2=1$ bounds the positive part of $\lambda$, and $S_{\mathcal C}(\varphi\|\omega)=-\int\lambda\,d\mu<\infty$ then bounds the negative part, so $\int|\lambda|\,d\mu<\infty$ and $\varphi([D\omega:D\varphi]_t)$ is differentiable at $0$ with derivative $-iS_{\mathcal C}(\varphi\|\omega)$~\cite{OhyaPetz1993}. The same argument on $\mathcal C'$ gives $\partial_t\varphi'([D\omega':D\varphi']_t)|_{t=0}=-iS_{\mathcal C'}(\varphi\|\omega)$, and since $[D\varphi':D\omega']_t=[D\omega':D\varphi']_t^*$, the derivative of $\langle\Psi|b_t|\Psi\rangle$ at $0$ is $+iS_{\mathcal C'}(\varphi\|\omega)$. Finally
\begin{equation}
 \langle\Psi|a_tb_t|\Psi\rangle-1
 =\bigl(\langle\Psi|a_t|\Psi\rangle-1\bigr)+\bigl(\langle\Psi|b_t|\Psi\rangle-1\bigr)
 +\langle(a_t^*-1)\Psi|(b_t-1)\Psi\rangle,
\end{equation}
and $\|(a_t^*-1)\Psi\|^2=\|(a_t-1)\Psi\|^2=2\operatorname{Re}\bigl(1-\langle\Psi|a_t|\Psi\rangle\bigr)=o(t)$, because the derivative of $\langle\Psi|a_t|\Psi\rangle$ at $0$ is purely imaginary; likewise for $b_t$. The cross term is therefore $o(t)$, and
$-i\,\partial_t\langle\Psi|\Delta_{\omega;\mathcal C}^{it}|\Psi\rangle|_{t=0}=-S_{\mathcal C}(\varphi\|\omega)+S_{\mathcal C'}(\varphi\|\omega)$. The last statement is the spectral theorem for $\log\Delta_{\omega;\mathcal C}$ on its form domain. In finite dimensions $\log\Delta_{\omega;\mathcal C}=\log\rho^\omega_{\mathcal C}\otimes1-1\otimes\log\rho^\omega_{\mathcal C'}$ on $\mathcal H_{\mathcal C}\otimes\mathcal H_{\mathcal C'}$, and \eqref{eq:modular-energy} reduces to the equality of the von Neumann entropies of the two marginals of a pure state.
\end{proof}

\begin{proposition}[Sufficiency decomposition of the subsystem-flow generator]\label{prop:sufficiency}
Let $\Omega$ be a common cyclic and separating vector for $\mN\subset\mM$, and let $\Psi$ be a unit vector, cyclic and separating for $\mN$ and for $\mM$, for which the four relative entropies $S_{\mM}(\varphi\|\omega)$, $S_{\mN}(\varphi\|\omega)$, $S_{\mN'}(\varphi\|\omega)$, $S_{\mM'}(\varphi\|\omega)$ are finite and $\Psi\in\operatorname{Dom}K_{\mM}\cap\operatorname{Dom}K_{\mN}$, where $\varphi=\braket{\Psi|\cdot\,\Psi}$. Then
\begin{align}
 -i\,\partial_t\langle\Psi|D_{\mN\subset\mM}(t)|\Psi\rangle\big|_{t=0}
 &=\langle\Psi|K_{\mM}-K_{\mN}|\Psi\rangle \nonumber\\
 &=\underbrace{\bigl[S_{\mM}(\varphi\|\omega)-S_{\mN}(\varphi\|\omega)\bigr]}_{\text{gap for }\mN\subset\mM}
 \nonumber\\
 &\quad+\underbrace{\bigl[S_{\mN'}(\varphi\|\omega)-S_{\mM'}(\varphi\|\omega)\bigr]}_{\text{gap for }\mM'\subset\mN'} ,
 \label{eq:sufficiency-decomposition}
\end{align}
with both gaps nonnegative. The generator vanishes on $\Psi$ if and only if the restriction $\mM\to\mN$ is sufficient for $(\varphi,\omega)$ and the restriction $\mN'\to\mM'$ is sufficient for $(\varphi,\omega)$.
\end{proposition}

\begin{proof}
On $\operatorname{Dom}K_{\mM}\cap\operatorname{Dom}K_{\mN}$ the generator of $D_{\mN\subset\mM}(t)=\Delta_{\mM}^{-it}\Delta_{\mN}^{it}$ at $t=0$ is $K_{\mM}-K_{\mN}$, which is the first equality. Lemma~\ref{lem:modular-energy}, applied to $\mathcal C=\mM$ and to $\mathcal C=\mN$, gives $\langle\Psi|K_{\mM}|\Psi\rangle=S_{\mM}(\varphi\|\omega)-S_{\mM'}(\varphi\|\omega)$ and $\langle\Psi|K_{\mN}|\Psi\rangle=S_{\mN}(\varphi\|\omega)-S_{\mN'}(\varphi\|\omega)$; subtracting and regrouping gives the second equality. Nonnegativity of each bracket is monotonicity of relative entropy under $\mN\subset\mM$ and $\mM'\subset\mN'$. Vanishing of the sum of two nonnegative terms is vanishing of each, and by Theorem~\ref{thm:Petz} each vanishes exactly when the corresponding restriction is sufficient for $(\varphi,\omega)$.
\end{proof}

\begin{corollary}\label{cor:positive-generator}
Under the hypotheses of Proposition~\ref{prop:sufficiency}, if $\mN\subset\mM$ is proper then the quadratic form $\Psi\mapsto\langle\Psi|K_{\mM}-K_{\mN}|\Psi\rangle$ is nonnegative and does not vanish identically on any dense subspace of $\operatorname{Dom}K_{\mM}\cap\operatorname{Dom}K_{\mN}$ that is a core for $K_{\mM}$: there exist states $\varphi$ for which the erasure \eqref{eq:erasure-channel} is not sufficient, and the subsystem flow registers them.
\end{corollary}
\begin{proof}
Nonnegativity is \eqref{eq:sufficiency-decomposition}. If the form vanished on a dense subspace $D\subseteq\operatorname{Dom}K_{\mM}\cap\operatorname{Dom}K_{\mN}$, polarization would give $K_{\mM}\Psi=K_{\mN}\Psi$ for $\Psi\in D$; if $D$ is a core for $K_{\mM}$, the closure of $K_{\mM}|_D$ is $K_{\mM}$ and is contained in the self-adjoint operator $K_{\mN}$, so $K_{\mM}=K_{\mN}$, hence $\Delta_{\omega;\mM}=\Delta_{\omega;\mN}$ and $\sigma^{\omega}_t(\mN)=\mN$ for all $t$, contradicting Corollary~\ref{cor:common-cyclic-noncorrectable}.
\end{proof}

\paragraph{Failure of Knill--Laflamme is measured by the subsystem flow.}
When conditions \NL{(i)--(iv)} of Theorem~\ref{thm:correctable} fail, the two gaps of Proposition~\ref{prop:sufficiency}  are the two sides of \eqref{eq:KL}--\eqref{eq:forgetful} made quantitative: $S_{\mM}(\psi\|\omega)-S_{\mN}(\psi\|\omega)$ measures the failure of recovery from the retained algebra, and $S_{\mN'}(\psi\|\omega)-S_{\mM'}(\psi\|\omega)$ measures what the environment has learned. They are independent as functionals of $\psi$, and no inequality relates them pair by pair. For $\psi=u\Omega$ with $u$ a unitary of $\mN$ the second gap vanishes identically, since $\psi$ agrees with $\omega$ on $\mN'$, while the Connes cocycle of $(\psi,\omega)$ on $\mM$ is $u\,\sigma^\omega_t(u)^*$, so by Theorem~\ref{thm:Petz} the first gap at $u\Omega$ vanishes exactly when $\sigma^\omega_t(u)\in\mN$ for all $t$: Petz sufficiency, quantified over the logical unitaries, is Takesaki's condition. On the logical orbit $\{u\Omega\}$ the quadratic form of Corollary~\ref{cor:positive-generator} therefore reduces to a single gap, and it is this gap that detects non-correctability. This is no conflict with the information--disturbance trade-off of~\cite{KretschmannSchlingemannWerner2008}, which quantifies over a code subspace rather than over a pair of states: the span of the logical orbit is all of $\mH_{\mN}$, on which the environment does learn. The trade-off survives in the modular setting only at the level of the algebra, as the equivalence of \NL{Theorem}~\ref{thm:correctable}, and the subsystem flow is its generating function state by state.

\subsection{History, invariant cores, generation and forgetfulness}\label{subsec:QEC-history-memory}

For a state with correlations between $\mN$ and its relative commutant, so that $\sigma^{\omega;\mM}_t(\mN)\ne\mN$, the orbit of $\mN$ under the modular flow determines a history algebra and an invariant core. If the relative commutant $\mN_{\rm hist}^c=(\mN_{\rm hist})'\cap \mM$ is non-trivial, it is itself an invariant subalgebra, while $\mN_{\rm hist}$, being modular invariant, admits an $\omega$-preserving expectation and is exactly correctable in the sense of Theorem~\ref{thm:correctable}. The information--disturbance relation is reflected in the identities $\mN_{\rm hist}^c=\mR_{\rm core}$ and $\mR_{\rm hist}\subset \mN_{\rm core}^c$ proved below.

The ergodic content of the subsystem flow enters through the orbit of $\mN$ under the modular flow. Let $\alpha_t=\sigma^{\omega}_t$, which is implemented on $\mathcal H$ by $\Delta_{\mM}^{it}$ and preserves $\mM$ and $\mM'$.

\begin{definition}[History and invariant-core algebras]\label{def:history-memory}
For a subalgebra $\mN\subset\mM$, define
\begin{align}
 \mN_{\rm hist}&:=\bigvee_{t\in\mathbb R}\alpha_t(\mN), &
 \mN_{\rm hist}^{+}&:=\bigvee_{t>0}\alpha_t(\mN), &
 \mN_{\rm hist}^{-}&:=\bigvee_{t<0}\alpha_t(\mN),\nn\\
 \mN_{\rm core}&:=\bigcap_{t\in\mathbb R}\alpha_t(\mN), &
 \mN_{\rm core}^{+}&:=\bigcap_{t>0}\alpha_t(\mN), &
 \mN_{\rm core}^{-}&:=\bigcap_{t<0}\alpha_t(\mN).
\end{align}
We call $\mN_{\rm hist}^{+}$ and $\mN_{\rm hist}^{-}$ the \emph{future} and \emph{past history}, and $\mN_{\rm core}^{+}$ and $\mN_{\rm core}^{-}$ the \emph{future} and \emph{past invariant cores}. The inclusion is \emph{generating}, \emph{future generating}, or \emph{past generating} when $\mN_{\rm hist}=\mM$, $\mN_{\rm hist}^{+}=\mM$, or $\mN_{\rm hist}^{-}=\mM$, respectively. We say it has \emph{trivial invariant core}, \emph{trivial future invariant core}, or \emph{trivial past invariant core} when the corresponding core is $\mathbb C1$.\footnote{For a one-sided nested family, the corresponding one-sided invariant core is commonly called the \emph{tail algebra}; triviality of this tail is the standard terminology in the definition of a $K$-system~\cite{CornfeldFominSinai1982,ouseph2024local}.} Future or past generation implies generation, and a trivial future or past invariant core implies a trivial invariant core; the converses need not hold without one-sided nesting.

Separately, \emph{forgetfulness} is an asymptotic dynamical property, not merely an intersection condition. For a normal unital completely positive family $\Phi_t:\mN\to\mN$ with faithful invariant state $\omega_{\mN}$, we call the dynamics \emph{future forgetful} if
\begin{equation}
 \lim_{t\to+\infty}\|\varphi\circ\Phi_t-\omega_{\mN}\|_{\mN_*}=0
 \qquad\text{for every normal state }\varphi\in\mN_* ,
 \label{eq:forgetful-dynamics}
\end{equation}
and define past forgetfulness analogously. For iterates of a quantum channel or for a quantum Markov semigroup this is the usual mixing property; in finite dimensions, when the stationary state is faithful, it is equivalent to primitivity. Triviality of an invariant core by itself does not imply this asymptotic property.
\end{definition}

The two-sided algebras $\mN_{\rm hist}$ and $\mN_{\rm core}$ are invariant under $\alpha_t$ by construction. By Takesaki's theorem each therefore admits an $\omega$-preserving normal expectation: $\mN_{\rm core}$ is the largest subalgebra of $\mN$ that is exactly correctable in the sense of Theorem~\ref{thm:correctable}, and $\mN_{\rm hist}$ is the smallest exactly correctable algebra containing $\mN$. In the language of noiseless subsystems, $\mN_{\rm core}$ is the noiseless subalgebra of $\mN$ for the composite channel ``evolve by the modular flow, then restrict to $\mN$'': the observables of $\mN$ that remain accessible from $\mN$ at every modular time. These statements are algebraic and should not be confused with the asymptotic forgetfulness defined above.

\begin{lemma}[History--core duality]\label{lem:history-memory}
Let $\alpha_t=\operatorname{Ad}U_t$ be a one-parameter group of automorphisms of $B(\mathcal H)$ with $U_t\mM U_t^*=\mM$, hence also $U_t\mM'U_t^*=\mM'$, and let history and invariant cores be defined as in Definition~\ref{def:history-memory} for any von Neumann algebra on $\mathcal H$. For every subalgebra $\mN\subset\mM$:
\begin{enumerate}
 \item[\rm(a)] On $B(\mathcal H)$ the duality is exact,
 \begin{equation}
  (\mN_{\rm hist})'=(\mN')_{\rm core},\qquad (\mN_{\rm core})'=(\mN')_{\rm hist}.
  \label{eq:history-memory-BH}
 \end{equation}
 The same identities hold separately for the future and past algebras,
 $(\mN_{\rm hist}^{\pm})'=(\mN')_{\rm core}^{\pm}$ and
 $(\mN_{\rm core}^{\pm})'=(\mN')_{\rm hist}^{\pm}$.
 \item[\rm(b)] Inside $\mM$, with $\mR=\mN'\cap\mM$,
 \begin{equation}
  (\mN_{\rm hist})'\cap\mM=\mR_{\rm core},\qquad
  \mR_{\rm hist}\subseteq(\mN_{\rm core})'\cap\mM ,
  \label{eq:history-memory}
 \end{equation}
 with equality in the second relation when $\mN$ and $\mR_{\rm hist}$ each coincide with their double relative commutant in $\mM$; this holds in particular when $\mM=\mN\overline\otimes\mR$ and $\mR$ is autonomous.
\end{enumerate}
\end{lemma}
\begin{proof}
(a) For von Neumann algebras the commutant of a join is the intersection of the commutants and the commutant of an intersection is the join of the commutants, and $\operatorname{Ad}U_t$ commutes with taking commutants; hence $(\bigvee_t\alpha_t(\mN))'=\bigcap_t\alpha_t(\mN')$ and $(\bigcap_t\alpha_t(\mN))'=\bigvee_t\alpha_t(\mN')$.
(b) Intersecting the first identity of (a) with $\mM=\alpha_t(\mM)$ gives $\bigcap_t\alpha_t(\mN')\cap\mM=\bigcap_t\alpha_t(\mN'\cap\mM)=\mR_{\rm core}$. For the inclusion, $\alpha_t(\mN'\cap\mM)\subseteq\alpha_t(\mN')\cap\mM\subseteq(\bigcap_s\alpha_s(\mN))'\cap\mM$ for every $t$; take the join over $t$. For the equality statement write $\mathcal P^c=\mathcal P'\cap\mM$. The first relation of (b) applied to $\mR$ in place of $\mN$ reads $(\mR_{\rm hist})^c=(\mR^c)_{\rm core}$, and $\mR^c=\mN$ when $\mN=\mN^{cc}$; taking relative commutants once more, $(\mR_{\rm hist})^{cc}=(\mN_{\rm core})^c$, which equals $\mR_{\rm hist}$ when $\mR_{\rm hist}=(\mR_{\rm hist})^{cc}$. In the tensor-product case $\mN=\mN^{cc}$, and if $\mR$ is autonomous then $\mR_{\rm hist}=\mR=\mR^{cc}$.
\end{proof}

\begin{corollary}[Generation forces a trivial environment invariant core]\label{cor:generating-forgetful}
If $\mM$ is a factor and $\mN\subset\mM$ is generating, then $\mR_{\rm core}=\mathbb C1$: the relative commutant has trivial invariant core under the modular flow. In particular this holds whenever $\Omega$ is a common cyclic and separating vector for $\mN$ and $\mM$ {\rm(Lemma~\ref{lem:common-cyclicity-generating})}.
\end{corollary}
\begin{proof}
$\mR_{\rm core}=(\mN_{\rm hist})'\cap\mM=\mM'\cap\mM=\mathbb C1$.
\end{proof}

Corollary~\ref{cor:generating-forgetful} is the modular form of the decoupling principle of error correction: if the logical algebra $\mN$ dynamically generates the full algebra, then its environment $\mR$ has no nontrivial invariant core. Since the Stinespring environment of the erasure is the full commutant $\mN'$ (Section~\ref{subsec:QEC-KL}), part (a) of the lemma completes the picture: $(\mN')_{\rm core}=(\mN_{\rm hist})'\supseteq\mM'$, with equality exactly when $\mN$ is generating. The reference $\mM'$ is the frozen part of the environment, fixed setwise by the modular flow of $\omega$ and therefore always contained in its invariant core; $\mR$ is the part the flow can move. The operators in $(\mN_{\rm core})'\cap\mM$ that are not in $\mR_{\rm hist}$ commute with the permanent invariant core of $\mN$ but cannot be generated from within $\mR$; they are the modular analogue of boundary degrees of freedom. This corollary is algebraic: it does not imply asymptotic forgetfulness or a decay rate for connected correlations.

\paragraph{\NL{Common cyclicity and standard inclusions.}}
\label{subsec:common-cyclic}
\begin{lemma}[Common cyclicity implies generation]\label{lem:common-cyclicity-generating}
If $\mN\subset\mM$ has a common cyclic and separating vector $\Omega$, then
$\bigvee_{t\in\mathbb R}\sigma^{\omega;\mM}_t(\mN)=\mM$.
\end{lemma}
\begin{proof}
Let $\mP=\bigvee_t\sigma^{\omega;\mM}_t(\mN)$. It is modular invariant, so Takesaki's theorem
gives an $\Omega$-preserving expectation $\mE_{\mP}:\mM\to\mP$.
Since $\mN\subset\mP$ and $\Omega$ is cyclic for $\mN$, the projection onto
$\overline{\mP\Omega}$ is $1$. Hence $\mE_{\mP}(a)\Omega=a\Omega$ for every $a\in\mM$.
Separatingness for $\mM$ gives $\mE_{\mP}(a)=a$.
\end{proof}
\begin{corollary}\label{cor:common-cyclic-noncorrectable}
If, in addition to the hypotheses of Theorem~\ref{thm:correctable}, $\Omega$ is cyclic for $\mN$ and the inclusion is proper, then none of \NL{{\rm(i)--(iv)}} of that theorem holds: the erasure \eqref{eq:erasure-channel} is not exactly correctable by any $\omega$-preserving expectation on any family of states containing $\omega$.
\end{corollary}

\begin{proof}
With $e_{\mN}=1$, (ii) of Theorem~\ref{thm:correctable} would give $\mE(m)\Omega=m\Omega$ and, by separation, $\mE(m)=m$ for all $m$, forcing $\mN=\mM$ (Lemma~\ref{lem:common-cyclicity-generating}).
\end{proof}

\paragraph{Local algebras of QFT (no normal conditional expectations):}
In a local conformal net $\{\mathcal{A}(I)\}$ on the circle, let $I_1,I_2$ be disjoint
open intervals with closures sharing exactly one endpoint, and set $I=\operatorname{int}\overline{I_1\cup I_2}$.
Haag duality and strong additivity give $\mathcal{A}(I_1)'\cap\mathcal{A}(I)=\mathcal{A}(I_2)$
and $\mathcal{A}(I_1)\vee\mathcal{A}(I_2)=\mathcal{A}(I)$, so the inclusion
$\mN:=\mathcal{A}(I_1)\subset\mM:=\mathcal{A}(I)$ is conormal with
$\mR=\mathcal A(I_2)$. The adjacent pair has no spatial tensor-product
factorization in a nontrivial conformal net~\cite{Wiesbrock1993}.
Proposition~\ref{prop:conormal-factorization} therefore excludes faithful normal expectations, in particular any expectation preserving a faithful normal reference state. This is the conclusion needed below. The obstruction concerns normal product structure for this inclusion; it does not exclude error correction for other channels or chosen families of states.

\begin{lemma}[Standard inclusions of factors are generating with trivial invariant core]
\label{lem:standard-forgetful}
Let $\mN \subset \mM$ be a standard inclusion of factors. Then the inclusion is generating and has trivial invariant core, $\mN_{\rm core} = \mathbb C1$.
\end{lemma}
\begin{proof}
Since in a standard inclusion $\mN$, $\mM$ and the relative commutant $\mR = \mN' \cap \mM$ share a common cyclic and separating vector, we know that the inclusion is generating. Since $\mR \subset \mM$, the modular orbit of $\mR$ generates the full ambient algebra, giving $\mR_{\rm hist} = \mM$.

By Lemma~\ref{lem:history-memory}(b), the history of the relative commutant is contained within the commutant of the invariant core relative to $\mM$:
\begin{equation}
    \mR_{\rm hist} \subseteq \mN_{\rm core}' \cap \mM.
\end{equation}
Substituting $\mR_{\rm hist} = \mM$ yields $\mM \subseteq \mN_{\rm core}' \cap \mM$, which forces the equality $\mN_{\rm core}' \cap \mM = \mM$. This equality dictates that every element of $\mN_{\rm core}$ commutes with every element of $\mM$.

Because $\mN_{\rm core}$ is a subalgebra of $\mM$ (via $\mN$), it must reside entirely within the center of $\mM$:
\begin{equation}
    \mN_{\rm core} \subseteq \mM' \cap \mM = \mathcal Z(\mM).
\end{equation}
Since $\mM$ is a factor, its center is trivial ($\mathcal Z(\mM) = \mathbb C1$). This restricts the invariant core to $\mN_{\rm core} = \mathbb C1$.
\end{proof}

To make the discussion less abstract, below, we give a few examples of the concepts introduced in this subsection:

\paragraph{Finite-dimensional example (generating inclusion)}
Everything in this section already happens for two qubits. Take $\mM=M_2\otimes M_2$, $\mN=M_2\otimes1$, $\mR=1\otimes M_2$, and the reference density matrix $\rho=(1-\epsilon)\ket{\Phi^+}\bra{\Phi^+}+\epsilon\,\mathbbm 1/4$ with $\ket{\Phi^+}$ maximally entangled and $0<\epsilon<1$, purified by a vector $\Omega$ that is cyclic and separating for $\mM$. Its spectrum $\{1-3\epsilon/4,\epsilon/4,\epsilon/4,\epsilon/4\}$ is not a product of two two-point spectra, so $\rho\ne U(\rho_B\otimes\rho_R)U^\dagger$ for every unitary $U$: no encoding of one qubit into the two has $\rho$ as a code state, and Theorem~\ref{thm:correctable} fails for $\mN$ and for every $U\mN U^\dagger$. The modular Hamiltonian $K_{\mM}=-\log\rho$ differs from a multiple of the identity by a multiple of $\ket{\Phi^+}\bra{\Phi^+}=\tfrac14(\mathbbm 1+X\otimes X-Y\otimes Y+Z\otimes Z)$, a pure interaction term, and $i[K_{\mM},Z\otimes1]\propto X\otimes Y+Y\otimes X$ leaves $\mN$: this is the leakage the subsystem flow measures. Since $\rho_B=\rho_R=\mathbbm 1/2$, a logical unitary $u\otimes1$ changes neither marginal, so for $\Psi=(u\otimes1)\Omega$ the whole sufficiency gap of Proposition~\ref{prop:sufficiency} sits in the correlations between $\mN$ and $\mR$. Finally the inclusion is generating: if $1\otimes r\in\mR_{\rm core}$ then $i[K_{\mM},1\otimes r]\in\mR$, and averaging this commutator over the unitaries of $\mN$, which fix $\mR$ pointwise and average $\ket{\Phi^+}\bra{\Phi^+}$ to $\mathbbm 1/4$, shows that $1\otimes r$ commutes with $\ket{\Phi^+}\bra{\Phi^+}$, hence $r\propto1$. Thus $\mR_{\rm core}=\mathbb C1$, so by Lemma~\ref{lem:history-memory} the algebra $\mN_{\rm hist}$, which is of the form $M_2\otimes\mathcal P$, has trivial relative commutant and equals $\mM$.

\paragraph{Tensor amplification (an expected, nongenerating inclusion):}
Let $\mN$ be a type $III_1$ factor and set $\mM=\mN\overline\otimes M_n(\mathbb C)$,
with $n>1$. Then $\mR=1\otimes M_n(\mathbb C)$ and the inclusion is conormal.
The slice expectations \eqref{eq:slice-expectation-intro} exist, and the minimal Jones--Kosaki index
is $n^2$. Finite index by itself neither specifies a matrix relative commutant
nor implies conormality.
For a faithful product state $\omega_{\mM}=\omega_{\mN}\otimes\omega_{\mR}$, the outer modular
flow preserves $\mN$, so its modular orbit generates only $\mN$.
The site spectral projections of $\omega_{\mR}$ lie in the centralizer of $\omega_{\mM}$.
In the product GNS representation the compressed subsystem flow is trivial, $C_{\mN\subset\mM}(t)=1$, in accordance with Theorem~\ref{thm:correctable}. The reference vector is not cyclic for $\mN$ on the full Hilbert
space, consistently with Lemma~\ref{lem:common-cyclicity-generating}.

\subsection{Modular ergodic hierarchy}\label{subsec:QEC-hierarchy}

The properties relevant to this work can now be summarized in the language of modular ergodic theory.
\begin{center}
\begin{tabular}{@{}>{\raggedright\arraybackslash}p{0.20\linewidth}>{\raggedright\arraybackslash}p{0.33\linewidth}>{\raggedright\arraybackslash}p{0.39\linewidth}@{}}
 \textit{Property} & \textit{Modular formulation} & \textit{Error-correction reading}\\[4pt]
 Ergodicity & $\mM^\omega=\mathbb C1$ & no observable is conserved by the dynamics\\
 Mixing & $C_{a,b}^\omega(t)\to0$ for all $a,b\in\mM$ & connected correlations decay\\
 Autonomy of $\mN$ & $\sigma^\omega_t(\mN)=\mN$ for all $t$ & $\mN$ exactly correctable (Thm.~\ref{thm:correctable})\\
 Generation & $\mN_{\rm hist}=\mM$ & full algebra generated by the orbit\\
 Trivial invariant core & $\mN_{\rm core}=\mathbb C1$ & no nontrivial invariant subalgebra remains inside $\mN$\\
 Half-sided inclusion & $\sigma^\omega_t(\mN)\subseteq\mN$ for one sign of $t$ & one-sided loss of access
\end{tabular}
\end{center}
Here $C^\omega_{a,b}(t)$ is the connected modular correlation function \eqref{eq:modular-correlation}; mixing requires its decay for every pair, is a property of the state on $\mM$ rather than of the inclusion, and implies ergodicity. Asymptotic forgetfulness of a reduced CP dynamics is defined separately in Definition~\ref{def:history-memory}; for channels or Markov semigroups it is the usual mixing/primitive notion. The classical $K$-property of Appendix~\ref{subsec:classical-ergodic-hierarchy} is the conjunction of half-sided inclusion, one-sided generation, and a trivial one-sided invariant core (tail). Its modular counterpart is the \emph{modular $K$-system}, built on
Wiesbrock's half-sided modular inclusions~\cite{Wiesbrock1993}.
The construction of positive-energy translations and its relation to
modular chaos are given together in the companion QFT analysis~\cite{companionQFT}.

\section{Probing outside of $\mM$ with inclusion parent Hamiltonians}
\label{sec:probing-outside-M}

The analytic continuation of the subsystem flow defines a family of positive
operators intrinsic to an inclusion $\mN\subset \mM$. Exponentiating these operators gives a dynamics on the GNS
Hilbert space of the retained subsystem $\mH_{\mN}$. Unlike the modular flow of
$\mM$, this dynamics need not remain inside $\mM$. We prove that the ground space is the GNS closure of the
largest modular-invariant subalgebra inside $\mN$. Thus a trivial invariant
core makes $\Omega$ the unique ground state. Under the full-support
assumption, we also show
that if the ambient modular action on $\mN$ has trivial invariant core, then
the two-sided orbit of $\mN$ under this new positive-energy flow
generates the entire type-I algebra $B(\mH_{\mN})$. For a genuine compression,
its ambient corner necessarily contains operators outside $\mM$. Imaginary modular time therefore provides a canonical positive operator whose unitary orbit may probe beyond the original ambient algebra.
The three-lines theorem further singles out the midpoint of the modular strip
as the point of maximal spectral attenuation; its implication for the energy
gap depends on how the parent Hamiltonian is normalized. We also establish
the covariance of the interpolation tunnel, canonical towers, and their
operator-system analogues under the parent-Hamiltonian flow, distinguishing
this kinematic covariance from a nontrivial locality or propagation property.

\NL{We use the GNS spaces, embedding $V$, and analytic inclusion map $D(z)$ of Definition~\ref{def:subsystem-flow} and Section~\ref{subsec:analytic-subsystem-flow}. Throughout this section $\Omega$ is cyclic and separating for $\mM$, without assuming cyclicity for $\mN$ on $\mH_{\mM}$.}

\subsection{Positive inclusion parent Hamiltonians}

For $0<\alpha<\frac12$ define the positive inclusion operator
\begin{equation}
    \boxed{
    \rho_{\alpha,\mN\subset\mM}
    :=
    D_\alpha^*D_\alpha
    }
    \label{eq:rho-alpha-general}
\end{equation}
on $\mH_{\mN}$. By \eqref{eq:W-strip-bound},
\begin{equation}
    0\leq
    \rho_{\alpha,\mN\subset\mM}
    \leq
    1_{\mH_{\mN}}.
    \label{eq:rho-alpha-contraction}
\end{equation}
Since all modular operators fix the reference vector,
\begin{equation}
    D_\alpha\Omega=\Omega,
\end{equation}
and hence
\begin{equation}
    \boxed{
    \rho_{\alpha,\mN\subset\mM}\Omega=\Omega .
    }
    \label{eq:rho-alpha-vacuum}
\end{equation}

We use the normalization
\begin{equation}
    \boxed{
    Q_{\alpha,\mN\subset\mM}
    :=
    -\frac{1}{2\sin(2\pi\alpha)}
    \log
    \rho_{\alpha,\mN\subset\mM},
    \qquad
    0<\alpha<\frac12 .
    }
    \label{eq:H-alpha-general}
\end{equation}
The logarithm is defined by the spectral calculus. If
$\rho_{\alpha,\mN\subset\mM}$ has full support, then
$Q_{\alpha,\mN\subset\mM}$ is a positive self-adjoint operator on the dense
domain. If $\rho_\alpha$ has a kernel, the same definition gives a self-adjoint
operator on its support; vectors in the kernel have formally infinite
inclusion energy. Whenever we discuss the unitary group generated by
$Q_{\alpha,\mN\subset\mM}$ on all of $\mH_{\mN}$, we assume it has full support. 

There is no logarithmic branch ambiguity in
\eqref{eq:H-alpha-general}: $\rho_\alpha$ is positive and its spectrum lies
in $[0,1]$. Therefore
\begin{equation}
    \boxed{
    Q_{\alpha,\mN\subset\mM}\geq0,
    \qquad
    Q_{\alpha,\mN\subset\mM}\Omega=0.
    }
    \label{eq:H-alpha-positive}
\end{equation}
Imaginary modular time thus converts the contractive analytic interpolation
of the inclusion into a canonical positive inclusion parent Hamiltonian.

The normalization in \eqref{eq:H-alpha-general} is chosen so that
$Q_{\alpha,\mN\subset\mM}$ is independent of $\alpha$ for a half-sided
modular inclusion, as shown below. At the symmetric point $\alpha=1/4$,
\begin{equation}
    Q_{\frac14,\mN\subset\mM}
    =
    -\frac12
    \log\left(D_{1/4}^*D_{1/4}\right).
    \label{eq:H-quarter}
\end{equation}

\paragraph{Properties of inclusion parent Hamiltonians.}
The family $Q_{\alpha,\mN\subset\mM}$ has the following properties.
\begin{enumerate}[(a)]
\item \emph{Positivity and reference vacuum.}
The contraction property of $D_\alpha$ and
$D_\alpha\Omega=\Omega$ imply
\begin{equation}
    Q_{\alpha,\mN\subset\mM}\geq0,
    \qquad
    Q_{\alpha,\mN\subset\mM}\Omega=0.
    \label{eq:Q-parent-positive-property}
\end{equation}
Hence $Q_\alpha$ is a parent Hamiltonian for $\Omega$, although
neither uniqueness of its ground state nor locality of its unitary
flow follows from positivity.

\item \emph{Reflection across the strip.}
Using \eqref{eq:characteristic-bounded},
\begin{align}
    J_{\mM}D_\alpha J_{\mN}
    &=
    \Delta_{\mM}^{-\alpha}
    (J_{\mM}VJ_{\mN})
    \Delta_{\mN}^\alpha
    \nonumber\\
    &=
    \Delta_{\mM}^{1/2-\alpha}
    V\Delta_{\mN}^{-(1/2-\alpha)}
    =D_{\frac12-\alpha}.
    \label{eq:W-reflection}
\end{align}
Consequently,
\begin{equation}
    \boxed{
    J_{\mN}\rho_{\alpha,\mN\subset\mM}J_{\mN}
    =\rho_{\frac12-\alpha,\mN\subset\mM},
    }
    \label{eq:rho-reflection}
\end{equation}
and, because $\sin(2\pi(1/2-\alpha))=\sin(2\pi\alpha)$,
\begin{equation}
    \boxed{
    J_{\mN}Q_{\alpha,\mN\subset\mM}J_{\mN}
    =Q_{\frac12-\alpha,\mN\subset\mM}.
    }
    \label{eq:H-reflection}
\end{equation}
At the midpoint,
\begin{equation}
    J_{\mN}Q_{\frac14,\mN\subset\mM}J_{\mN}
    =Q_{\frac14,\mN\subset\mM}.
    \label{eq:H-quarter-real}
\end{equation}
This is an antiunitary time-reflection symmetry of the midpoint flow,
not by itself a statement that the flow is a geometric null translation.

\item \emph{Action of $J_{\mM}$.}
The modular conjugation of the ambient algebra acts naturally on
$\mH_{\mM}$, whereas $\rho_\alpha$ and $Q_\alpha$ act on
$\mH_{\mN}$. For a general compressed inclusion, the well-defined
ambient identity following from \eqref{eq:W-reflection} is
\begin{equation}
    \boxed{
    J_{\mM}(D_\alpha D_\alpha^*)J_{\mM}
    =D_{\frac12-\alpha}D_{\frac12-\alpha}^*.
    }
    \label{eq:Q-JM-ambient}
\end{equation}
If $\Omega$ is cyclic for both algebras, so that
$\mH_{\mN}=\mH_{\mM}$, put
$\widetilde\mN=J_{\mM}\mN J_{\mM}\subset\mM'$. Then
\begin{equation}
    \boxed{
    J_{\mM}Q_{\alpha,\mN\subset\mM}J_{\mM}
    =Q_{\alpha,\widetilde\mN\subset\mM'}.
    }
    \label{eq:Q-JM-mirror}
\end{equation}
This identifies the parent Hamiltonian with that of the reflected
inclusion, rather than asserting $J_{\mM}Q_\alpha J_{\mM}=Q_\alpha$.
In the common-cyclic representation, setting
$\Gamma=J_{\mM}J_{\mN}$ also gives
$J_{\mM}Q_\alpha J_{\mM}
=\Gamma Q_{1/2-\alpha}\Gamma^*$.

\item \emph{Monotonicity under further restriction.}
Consider $\mN_2\subset\mN_1\subset\mM$ with the same reference
vector. Writing $D_\alpha^{\mB\leftarrow\mA}$ for the characteristic
contraction of $\mA\subset\mB$, analyticity gives the composition law
\begin{equation}
    D_\alpha^{\mM\leftarrow\mN_2}
    =D_\alpha^{\mM\leftarrow\mN_1}
     D_\alpha^{\mN_1\leftarrow\mN_2}.
    \label{eq:Q-inclusion-composition}
\end{equation}
With $C=D_\alpha^{\mN_1\leftarrow\mN_2}$, a contraction,
\begin{equation}
    \rho_{\alpha,\mN_2\subset\mM}
    =C^*\rho_{\alpha,\mN_1\subset\mM}C
    \leq C^*C\leq\mathbf1.
    \label{eq:Q-inclusion-restriction}
\end{equation}
In particular, when the characteristic operators are compact, their
ordered singular values obey
\begin{equation}
    s_k(D_\alpha^{\mM\leftarrow\mN_2})
    \leq s_k(D_\alpha^{\mM\leftarrow\mN_1}).
    \label{eq:Q-spectral-monotonicity}
\end{equation}
Thus the ordered inclusion-energy levels cannot decrease when the
retained algebra shrinks; the corresponding Schatten norms and traces,
when finite, are also nonincreasing. This is a \emph{spectral}
monotonicity statement, not in general the Loewner-order inequality
$Q_{\alpha,\mN_2\subset\mM}\geq
Q_{\alpha,\mN_1\subset\mM}$: the two operators act on different
GNS subspaces, and even with common cyclicity the intervening
contraction need not preserve operator order.

\item \emph{The limit $\alpha\to0$.}
Let $K_\mA=-\log\Delta_\mA$ denote the modular Hamiltonian of
$\mA$. The first-order logarithmic expansion of $D_\alpha^*D_\alpha$
gives, on an appropriate common quadratic-form domain,
\begin{equation}
    \boxed{
    \lim_{\alpha\downarrow0}Q_{\alpha,\mN\subset\mM}
    =\frac1{2\pi}
    \bigl(V^*K_\mM V-K_\mN\bigr).
    }
    \label{eq:Q-alpha-zero}
\end{equation}
For a common cyclic and separating vector this reduces to
$(K_\mM-K_\mN)/(2\pi)$. The limit is understood as a form limit
under the necessary differentiability and domain conditions; it is
not an automatic norm limit for unbounded QFT modular operators.

\item \emph{Invariant core and the unique ground state.}
The largest ambient-modular-invariant subalgebra of $\mN$ is
\begin{equation}
    \mN_{\rm core}^{\,\mM}
    :=\bigcap_{t\in\mathbb R}\sigma_t^{\omega;\mM}(\mN).
    \label{eq:Q-core-property}
\end{equation}
Theorem~\ref{thm:Q-kernel-core} below proves the \emph{full} kernel identity
\begin{equation}
    \boxed{\ker(1-\rho_\alpha)
    =\overline{\mN_{\rm core}^{\,\mM}\Omega},
    \qquad 0<\alpha<\tfrac12.}
    \label{eq:Q-full-kernel}
\end{equation}
For a full-support $\rho_\alpha$, the left-hand side is
$\ker Q_\alpha$. In particular, $\Omega$ is the unique ground vector
(up to phase) \emph{if and only if} the ambient invariant core is trivial.
This requires neither compactness nor nuclearity. By
Lemma~\ref{lem:standard-forgetful}, a standard inclusion of factors
(with $\Omega$ also cyclic and separating for $\mN'\cap\mM$)
is a sufficient case.

\item \emph{Three-lines majorization and the special midpoint.}
The Hadamard three-lines theorem implies that
\begin{equation}
    f_k(\alpha):=\log\left\|\bigwedge^k D_\alpha\right\|
    \quad\text{is convex and satisfies}\quad
    f_k(\tfrac12-\alpha)=f_k(\alpha).
    \label{eq:Q-three-lines-convex}
\end{equation}
Consequently, the sums of the lowest $k$ spectral levels of the
\emph{unnormalized} positive operator $-\log\rho_\alpha$ (when these
levels are discrete) increase on $(0,1/4)$ and attain their largest
value at the strip midpoint. The same is true of
\begin{equation}
    Q_\alpha^{\rm gap}
    :=-\frac12\log\rho_\alpha
    =\sin(2\pi\alpha)Q_\alpha,
    \qquad Q_{1/4}^{\rm gap}=Q_{1/4}.
    \label{eq:Q-gap-normalization}
\end{equation}
In particular, the gap of this \emph{gap-normalized} family is maximal
at $\alpha=1/4$. This is \emph{not} a claim that the gap of the
original sine-normalized $Q_\alpha$ is maximal there: the sine factor
removes the $\alpha$ dependence altogether for half-sided modular
inclusions and can change the direction of monotonicity in general.
Theorem~\ref{thm:Q-three-lines} gives the precise result and also
an alternative normalization with monotone spectral partial sums.
\end{enumerate}

\subsection{Ground states from generalized conditional expectations}
\label{subsec:Q-ground-space}

The absence of a nontrivial invariant core controls all ground-state
vectors, not only vectors of the form $a\Omega$ with bounded $a\in\mN$.
The distinction matters because cyclicity of $\Omega$ only asserts that
$\mN\Omega$ is dense, not that its intersection with a given closed
subspace is dense in that subspace. The additional input is a normal
state-preserving completely positive map.

\begin{theorem}[Ground-state space and modular invariant core]
\label{thm:Q-kernel-core}
Let $\mN\subset\mM$ be a unital inclusion, with $\Omega$
cyclic and separating for $\mM$, and let $\Delta_\mN$ act on
$\mH_\mN=\overline{\mN\Omega}$. For $0<\alpha<1/2$,
\begin{equation}
    \boxed{
    \ker(1-\rho_{\alpha,\mN\subset\mM})
    =\overline{\mN_{\rm core}^{\,\mM}\Omega}.}
    \label{eq:Q-kernel-core-theorem}
\end{equation}
The kernel of the inclusion parent Hamiltonian (on the support of
$\rho_\alpha$) is the same space. In particular, it is independent
of $\alpha$ in the open modular strip. Hence
\begin{equation}
    \boxed{\ker Q_{\alpha,\mN\subset\mM}=\mathbb C\Omega
    \quad\Longleftrightarrow\quad
    \mN_{\rm core}^{\,\mM}=\mathbb C1,}
    \label{eq:Q-core-unique-equivalence}
\end{equation}
when $\rho_\alpha$ has full support. A standard inclusion of factors
is a sufficient condition for a unique ground state by
Lemma~\ref{lem:standard-forgetful}.
\end{theorem}

\begin{proof}
First, the contraction $D(z)$ is analytic in the strip, and
$D(0)=V$. Equality in the contraction bound at one interior point,
$\|D_\alpha\xi\|=\|\xi\|$, implies by the vector-valued
maximum-modulus principle that $D(z)\xi=V\xi$ throughout the strip.
Thus
\begin{equation}
 \ker(1-\rho_\alpha)
 =\{\xi\in\mH_\mN:\ 
     \Delta_\mM^{it}V\xi=V\Delta_\mN^{it}\xi
       \text{ for all }t\in\mathbb R\},
 \label{eq:Q-kernel-modular-intertwining}
\end{equation}
independently of $\alpha$. In particular, for $a\in\mN$, the
intertwining identity and separatingness of $\Omega$ for $\mM$
show that $a$ belongs to the invariant core; the converse follows
from Takesaki's theorem applied to that core.

To obtain the closure of the bounded-observable zero modes, introduce
the symmetric $L^2$ embeddings
\begin{equation}
 j_\mA(a):=\Delta_\mA^{1/4}a\Omega,
 \qquad \mA=\mN,\mM,
 \label{eq:Q-symmetric-embedding}
\end{equation}
where $j_\mN(\mN)$ is dense in $\mH_\mN$.
The characteristic map obeys
$D_{1/4}j_\mN(a)=j_\mM(a)$ for $a\in\mN$.
Let $\mathsf E_\omega:\mM\to\mN$ be the
Accardi--Cecchini generalized conditional expectation associated
with the inclusion and the faithful reference state, and put
$\Phi=\mathsf E_\omega|_\mN$. It is a normal unital completely
positive map preserving $\omega|_\mN$, and its symmetric $L^2$
implementation is
\begin{equation}
 \boxed{\rho_{1/4}\,j_\mN(a)=j_\mN(\Phi(a)),
 \qquad a\in\mN.}
 \label{eq:Q-AC-L2}
\end{equation}
Indeed, for $a,b\in\mN$ the inner products of both sides
against $j_\mN(b)$ agree, by the defining symmetric pairing of
$\mathsf E_\omega$ and the identity
$D_{1/4}j_\mN(a)=j_\mM(a)$. By the
Accardi--Cecchini fixed-point theorem \cite{AccardiCecchini1982},
\begin{equation}
 \operatorname{Fix}(\Phi)
 :=\{a\in\mN:\Phi(a)=a\}
 =\mN_{\rm core}^{\,\mM}.
 \label{eq:Q-AC-fixed-point}
\end{equation}

Let $P_0$ denote the spectral projection of $\rho_{1/4}$
onto its eigenvalue $1$. Since $0\leq\rho_{1/4}\leq1$,
\begin{equation}
 P_0=\operatorname*{s-lim}_{n\to\infty}\rho_{1/4}^{\,n}.
\end{equation}
For $a\in\mN$ we have
$\rho_{1/4}^{\,n}j_\mN(a)=j_\mN(\Phi^n(a))$.
The iterates satisfy $\|\Phi^n(a)\|\leq\|a\|$; hence a subnet
converges ultraweakly to some $b\in\mN$. The symmetric embedding
$j_\mN$ is weakly continuous on norm-bounded sets, so
$P_0j_\mN(a)=j_\mN(b)$. Applying $\rho_{1/4}$ to this identity
and using \eqref{eq:Q-AC-L2} gives $\Phi(b)=b$. Thus
\begin{equation}
 P_0j_\mN(\mN)\subseteq
 j_\mN(\mN_{\rm core}^{\,\mM}).
\end{equation}
Conversely every fixed point of $\Phi$ gives an eigenvalue-one
vector by \eqref{eq:Q-AC-L2}. The density of $j_\mN(\mN)$
then implies
\begin{equation}
 \operatorname{Ran}P_0
 =\overline{j_\mN(\mN_{\rm core}^{\,\mM})}.
\end{equation}
The core is globally invariant under both modular groups, whose
restrictions to the core agree. Its $L^2$ closure therefore equals
its GNS closure:
$\overline{j_\mN(\mN_{\rm core}^{\,\mM})}
 =\overline{\mN_{\rm core}^{\,\mM}\Omega}$.
This proves the theorem for $\alpha=1/4$; the first step
extends it to all interior $\alpha$.
\end{proof}

\subsection{The inclusion parent Hamiltonian flow}

For $0<\alpha<1/2$ and under the standing full-support assumption, define
\begin{equation}
    U_t^{(\alpha)}:=e^{itQ_{\alpha,\mN\subset\mM}}.
\end{equation}
The reflection relation \eqref{eq:H-reflection} implies
\begin{equation}
    J_{\mN}U_t^{(\alpha)}J_{\mN}=U_{-t}^{(1/2-\alpha)}.
    \label{eq:H-quarter-time-reversal}
\end{equation}
At $\alpha=1/4$ this is an antiunitary time-reversal symmetry of a single flow. We identify $\mN$ with its faithful representation on
$\mH_{\mN}$ and define the inclusion parent Hamiltonian automorphism group by
\begin{equation}
    \gamma_t^\alpha(a)
    :=U_t^{(\alpha)}aU_t^{(\alpha)*},
    \qquad a\in B(\mH_{\mN}),\quad t\in\mathbb R.
    \label{eq:N-alpha-t-general}
\end{equation}
The transported subsystem is $\gamma_t^\alpha(\mN)$; the superscript
$\alpha$ labels the inclusion parent Hamiltonian, rather than a power of the map.
This flow is distinct from $\sigma_s^{\omega;\mM}$, which acts on the
ambient algebra $\mM$.
The algebra generated by the complete orbit is
\begin{equation}
    \mathcal A_\alpha
    :=
    \bigvee_{t\in\mathbb R}
    \gamma_t^\alpha(\mN)
    \subseteq B(\mH_{\mN}).
    \label{eq:A-alpha-general}
\end{equation}

To distinguish the two time orientations, define the future and past
inclusion parent Hamiltonian histories and cores by
\begin{align}
    \mathcal A_\alpha^\pm
    &:=\bigvee_{\pm t>0}\gamma_t^\alpha(\mN),
    &
    \mathcal C_\alpha^\pm
    &:=\bigcap_{\pm t>0}\gamma_t^\alpha(\mN),\nn\\
    \mathcal C_\alpha
    &:=\bigcap_{t\in\mathbb R}\gamma_t^\alpha(\mN).
    \label{eq:Halpha-one-sided-histories-cores}
\end{align}
We call the flow \emph{future generating} or \emph{past generating} on
$\mH_{\mN}$ when $\mathcal A_\alpha^+=B(\mH_{\mN})$ or
$\mathcal A_\alpha^-=B(\mH_{\mN})$, respectively, and say that its
corresponding one-sided core is trivial when
$\mathcal C_\alpha^\pm=\mathbb C1$. These cores refer to
$\gamma_t^\alpha$, whereas $\mN_{\rm core}^{\pm}$ in
Definition~\ref{def:history-memory} refers to the ambient modular flow.
The one-sided cores need only satisfy a one-sided inclusion under the
flow; the two-sided core $\mathcal C_\alpha$ is invariant under the
whole group.

This flow is fundamentally different from the ambient modular flow.
The latter is an automorphism group of $\mM$. By contrast,
$Q_{\alpha,\mN\subset\mM}$ is an operator on $\mH_{\mN}$, and there is
no reason for
\begin{equation}
    \gamma_t^\alpha(\mN)
\end{equation}
to have a realization inside $\mM$. The natural ambient realization
of its maximal algebra is instead
\begin{equation}
    B(\mH_{\mN})
    \simeq
    e_{\mN}B(\mH_{\mM})e_{\mN},
    \qquad
    a\longmapsto VaV^* .
    \label{eq:BHN-corner}
\end{equation}

The following lemma relates zero modes of the positive inclusion parent Hamiltonian
to observables retained by the original modular dynamics.

\begin{lemma}[Zero inclusion energy and modular invariant core]
\label{lem:zero-energy-modular-memory}
For $0<\alpha<1/2$ and $a\in\mN$,
\begin{equation}
 Q_{\alpha,\mN\subset\mM}a\Omega=0
 \quad\Longleftrightarrow\quad
 a\in\mN_{\rm core}^{\,\mM},
 \label{eq:Halpha-zero-a}
\end{equation}
where the left-hand side means that $\rho_\alpha a\Omega=a\Omega$.
In either case,
\begin{equation}
 \sigma_t^{\omega;\mM}(a)
 =\sigma_t^{\omega;\mN}(a)\in\mN,
 \qquad t\in\mathbb R.
 \label{eq:zero-energy-intertwining}
\end{equation}
\end{lemma}
\begin{proof}
This is the bounded-observable specialization of
Theorem~\ref{thm:Q-kernel-core}. If $a\Omega$ is a zero mode, the
modular-intertwining identity \eqref{eq:Q-kernel-modular-intertwining}
and separatingness of $\Omega$ imply equality of the two modular
actions on $a$, so $a$ lies in the invariant core. Conversely,
Takesaki's theorem for the core gives the modular-intertwining
identity for every $a$ in it.
\end{proof}

\paragraph{Positive-energy invariant algebras.}
We will also use the following elementary consequence of Tomita theory.
Let $\mathcal Q\subset B(\mH_{\mN})$ satisfy
\begin{equation}
    \gamma_t^\alpha(\mathcal Q)
    =
    \mathcal Q,
    \qquad
    t\in\mathbb R,
    \label{eq:Q-invariant-Halpha}
\end{equation}
and suppose $\Omega$ is separating for $\mathcal Q$. On
\begin{equation}
    \mathcal K_Q
    :=
    \overline{\mathcal Q\Omega}
\end{equation}
the vector $\Omega$ is cyclic and separating for the restricted algebra.
Since $U_t^{(\alpha)}\Omega=\Omega$ and
$U_t^{(\alpha)}$ preserves $\mathcal Q$, it commutes on
$\mathcal K_Q$ with the Tomita operator and hence with the modular
conjugation $J_Q$. Antiunitarity then gives
\begin{equation}
    J_Q
    Q_{\alpha,\mN\subset\mM}\big|_{\mathcal K_Q}
    J_Q
    =
    -
    Q_{\alpha,\mN\subset\mM}\big|_{\mathcal K_Q}.
\end{equation}
Since the restriction of $Q_{\alpha,\mN\subset\mM}$ is positive, it follows that
\begin{equation}
    \boxed{
    Q_{\alpha,\mN\subset\mM}\big|_{\mathcal K_Q}=0.
    }
    \label{eq:positive-H-invariant-Q}
\end{equation}
Thus a positive-energy flow fixing the reference vector can preserve a
standard subtheory only by acting trivially on the cyclic sector generated by
that subtheory.

\subsection{Trivial invariant core implies generation of $B(\mH_{\mN})$}

We can now formulate the main result of this section. Let
\begin{equation}
    \mN_{\rm core}^{\,\mM}
    :=
    \bigcap_{s\in\mathbb R}
    \sigma_s^{\omega;\mM}(\mN)
    \label{eq:ambient-memory-probing}
\end{equation}
be the two-sided invariant core of $\mN$ under the ambient modular flow.

\begin{proposition}[Trivial modular invariant core implies full inclusion parent Hamiltonian generation]
\label{prop:Halpha-full-generation}
Assume that the ambient modular action on $\mN$ has trivial invariant core,
\begin{equation}
    \boxed{
    \mN_{\rm core}^{\,\mM}
    =
    \bigcap_{s\in\mathbb R}
    \sigma_s^{\omega;\mM}(\mN)
    =
    \mathbb C1.
    }
    \label{eq:ambient-forgetful-probing}
\end{equation}
Let $0<\alpha<1/2$ and assume
$s(\rho_{\alpha,\mN\subset\mM})=1$. Then
\begin{equation}
    \boxed{
    \bigvee_{t\in\mathbb R}
    \gamma_t^\alpha(\mN)
    =
    B(\mH_{\mN}).
    }
    \label{eq:Halpha-full-BHN}
\end{equation}
\end{proposition}

\begin{proof}
By Theorem~\ref{thm:Q-kernel-core}, the trivial core implies
$\ker Q_\alpha=\mathbb C\Omega$ (using the full-support
hypothesis for the unitary group). Let $\mathcal A_\alpha$
be the algebra generated by the $Q_\alpha$-orbit of $\mN$.
It is invariant under $U_t^{(\alpha)}$, and its commutant
$\mathcal A_\alpha'$ is invariant as well. Since
$\mN\subset\mathcal A_\alpha$ and $\mN\Omega$ is dense
in $\mH_\mN$, $\Omega$ is separating for $\mathcal A_\alpha'$.
The positive-energy invariant-algebra argument
\eqref{eq:positive-H-invariant-Q}, applied to
$\mathcal A_\alpha'$, shows that
$Q_\alpha x\Omega=0$ for every $x\in\mathcal A_\alpha'$.
Uniqueness of the ground state forces $x\Omega=c_x\Omega$.
Separatingness for $\mathcal A_\alpha'$ then gives
$x=c_x1$. Hence $\mathcal A_\alpha'=\mathbb C1$ and
$\mathcal A_\alpha=B(\mH_\mN)$ by the bicommutant theorem.
\end{proof}

\begin{corollary}[Standard inclusions generate $B(\mH_{\mM})$]
\label{cor:standard-Halpha-full-generation}
Let $(\mN\subset\mM,\Omega)$ be a standard inclusion of factors,
let $0<\alpha<1/2$, and assume the standing full-support condition
$s(\rho_{\alpha,\mN\subset\mM})=1$ so that $Q_{\alpha,\mN\subset\mM}$ generates a
unitary group on the whole Hilbert space. Then
\begin{equation}
    \boxed{
    \bigvee_{t\in\mathbb R}
    \gamma_t^\alpha(\mN)
    =B(\mH_{\mM}).
    }
    \label{eq:standard-Halpha-full-BHM}
\end{equation}
\end{corollary}
\begin{proof}
Standardness implies that $\Omega$ is cyclic for $\mN$, hence
$\mH_{\mN}=\mH_{\mM}$ and $e_{\mN}=1$. It also implies a trivial ambient modular invariant core by Lemma~\ref{lem:standard-forgetful}. The result is therefore
Proposition~\ref{prop:Halpha-full-generation} with
$B(\mH_{\mN})=B(\mH_{\mM})$.
\end{proof}

This proposition makes precise the sense in which the inclusion parent Hamiltonian
probes more than the original ambient algebra. The modular orbit
\begin{equation}
    \sigma_s^{\omega;\mM}(\mN)
\end{equation}
always remains inside $\mM$. By contrast, under the same
trivial-core condition the orbit generated by $Q_{\alpha,\mN\subset\mM}$
fills the entire algebra of operators on the compressed GNS Hilbert space.

Embedding this algebra back into $\mH_{\mM}$ gives
\begin{equation}
    V B(\mH_{\mN})V^*
    =
    e_{\mN}B(\mH_{\mM})e_{\mN} .
    \label{eq:corner-BHN}
\end{equation}
For a genuine compression,
\begin{equation}
    e_{\mN}<1,
\end{equation}
this corner cannot be contained in $\mM$. Indeed, if
$e_{\mN}\in\mM$, then
\begin{equation}
    (1-e_{\mN})\Omega=0,
\end{equation}
and separatingness of $\Omega$ for $\mM$ would imply
$e_{\mN}=1$. Thus
\begin{equation}
    \boxed{
    e_{\mN}<1
    \quad\Longrightarrow\quad
    e_{\mN}B(\mH_{\mM})e_{\mN}
    \not\subset\mM.
    }
    \label{eq:corner-outside-M}
\end{equation}
In particular, the identity of the compressed algebra is represented on
$\mH_{\mM}$ by $e_{\mN}$, which itself already lies outside $\mM$.

Therefore the statement
\begin{equation}
    \bigvee_t\gamma_t^\alpha(\mN)
    =
    B(\mH_{\mN})
\end{equation}
should not be interpreted as generation of $\mM$. It is a stronger
statement in a different direction: after restricting to the GNS sector
created by the subsystem, the positive Hamiltonian extracted from imaginary
modular time explores every operator on that sector, including operators
which have no realization inside the original ambient algebra.

\paragraph{Future and past generation from pure point spectrum.}
The preceding proposition concerns the complete orbit. An additional
spectral hypothesis makes each time direction sufficient, without
requiring either one-sided ambient modular core to be trivial.

\begin{proposition}[Pure point spectrum implies future and past generation]
\label{prop:Halpha-pure-point-one-sided-generation}
Assume the hypotheses of Proposition~\ref{prop:Halpha-full-generation}.
Suppose, in addition, that $\mH_{\mN}$ is separable and
$Q_\alpha:=Q_{\alpha,\mN\subset\mM}$ has pure point spectrum, meaning
that its eigenvectors span $\mH_{\mN}$. Then
\begin{equation}
    \boxed{
    \mathcal A_\alpha^+
    =\mathcal A_\alpha^-
    =B(\mH_{\mN}).
    }
    \label{eq:Halpha-pure-point-one-sided-generation}
\end{equation}
In particular, this holds if
$\rho_{\alpha,\mN\subset\mM}$ is trace class with full support.
\end{proposition}

\begin{proof}
Choose an orthonormal eigenbasis $\{\psi_j\}$ of $Q_\alpha$, with
eigenvalues $E_j$. Simultaneous approximation of finitely many phases,
followed by a diagonal choice, gives times $t_n\to+\infty$ such that
$e^{it_nE_j}\to1$ for every $j$. Consequently
\begin{equation}
    U_{t_n}^{(\alpha)}\longrightarrow1,
    \qquad U_{-t_n}^{(\alpha)}\longrightarrow1
    \quad\text{strongly}.
    \label{eq:Halpha-pure-point-recurrence}
\end{equation}
For $a\in\mN$, $r\in\mathbb R$, and
$\varepsilon\in\{+1,-1\}$, it follows that
\begin{equation}
    \gamma_{\varepsilon t_n+r}^\alpha(a)
    \longrightarrow\gamma_r^\alpha(a)
    \quad\text{strongly}.
    \label{eq:Halpha-one-sided-recurrence-approximation}
\end{equation}
For large $n$, the times $\varepsilon t_n+r$ have sign $\varepsilon$.
The approximating operators therefore belong to
$\mathcal A_\alpha^+$ if $\varepsilon=+1$, and to
$\mathcal A_\alpha^-$ if $\varepsilon=-1$. Since von Neumann algebras
are strongly closed, each half-history contains
$\gamma_r^\alpha(\mN)$ for every real $r$. Hence
$\mathcal A_\alpha^+=\mathcal A_\alpha^-=\mathcal A_\alpha$, and
Proposition~\ref{prop:Halpha-full-generation} proves
\eqref{eq:Halpha-pure-point-one-sided-generation}.

If $\rho_\alpha$ is trace class with full support, it is a compact
positive operator with an orthonormal eigenbasis and positive
eigenvalues $\lambda_j$. The same basis diagonalizes $Q_\alpha$, with
$E_j=-[2\sin(2\pi\alpha)]^{-1}\log\lambda_j$, so the preceding argument
applies. Full support is essential here to define the unitary group
on all of $\mH_{\mN}$.
\end{proof}

For a standard inclusion of factors satisfying this spectral hypothesis,
$\mH_{\mN}=\mH_{\mM}$ and both half-histories therefore generate
$B(\mH_{\mM})$. In particular, the trace-class conditions discussed in
Section~\ref{sec:nuclearity}, together with full support, strengthen the
two-sided conclusion of Corollary~\ref{cor:standard-Halpha-full-generation}
to separate future and past generation. Pure point spectrum alone
does not replace the trivial ambient modular core hypothesis: recurrence
identifies the two half-histories with the complete history, while
Proposition~\ref{prop:Halpha-full-generation} identifies that history
with $B(\mH_{\mN})$.

\subsection{\NL{Interpolation tunnels, canonical towers and covariance}}
\label{subsec:Q-tower-covariance}

\subsubsection{\NL{Interpolation tunnels and canonical towers}}
\label{subsec:interpolation-tunnel}

An inclusion carries a natural partially ordered family of intermediate
observable algebras even when no state-preserving conditional expectation
exists. We first construct this family using the ambient modular flow;
we then add the discrete canonical tower when the reference vector is
cyclic and separating for both endpoint algebras. Throughout this
subsection, $\Omega$ is cyclic and separating for $\mM$, while it need
not be cyclic for $\mN$ on $\mH_{\mM}$.

\paragraph{The modular interpolation tunnel.}
For $T>0$, define
\begin{equation}
 \boxed{\quad
 \mI_{T,\mN\subset\mM}
 :=\bigvee_{|s|<T}\sigma_s^{\omega;\mM}(\mN)
 =\bigvee_{|s|<T}\Delta_{\mM}^{is}\mN\Delta_{\mM}^{-is}.
 \quad}
 \label{eq:interpolation-tunnel}
\end{equation}
We call $\{\mI_{T,\mN\subset\mM}\}_{T>0}$ the
\emph{interpolation tunnel} of the inclusion. Since the modular flow
preserves $\mM$, it satisfies
\begin{equation}
 \mN\subseteq\mI_{T_1,\mN\subset\mM}
 \subseteq\mI_{T_2,\mN\subset\mM}\subseteq\mM,
 \qquad 0<T_1<T_2,
 \label{eq:interpolation-nesting}
\end{equation}
and its union generates the full history algebra:
\begin{equation}
 \bigvee_{T>0}\mI_{T,\mN\subset\mM}=\mN_{\rm hist}.
 \label{eq:interpolation-limit}
\end{equation}
We adopt $\mI_{0,\mN\subset\mM}:=\mN$ as an endpoint
convention. No continuity at $T=0$ is implied: in a general inclusion
the algebra generated by arbitrarily short modular orbits can already
be larger than $\mN$. A simple, exact propagation property is
\begin{equation}
 \boxed{\quad
 \sigma_u^{\omega;\mM}(\mI_{T,\mN\subset\mM})
 \subseteq\mI_{T+|u|,\mN\subset\mM},\qquad u\in\mathbb R.
 \quad}
 \label{eq:interpolation-modular-propagation}
\end{equation}
This follows by translating the modular-time interval $(-T,T)$.
Thus modular time has a finite propagation bound in the filtration it
itself defines. The analogous statement for the inclusion parent
Hamiltonian is a separate dynamical question studied in
Section~\ref{subsec:Q-tower-covariance}.

The commutant gives a dual decreasing family,
\begin{equation}
 \mI_{T,\mN\subset\mM}'
 =\bigcap_{|s|<T}\sigma_s^{\omega;\mM}(\mN'),
 \label{eq:interpolation-commutant}
\end{equation}
which identifies the observables commuting with every member of the
modular history up to time $T$. In contrast to the canonical tower
below, the interpolation tunnel and its commutant exist for every
inclusion with the given faithful reference state.

\paragraph{Index of an inclusion:} The Pimsner--Popa index $\operatorname{Ind}(\mE)$ of a conditional expectation $\mE:\mM\to \mN$ is the infimum of $\lambda\geq1$ such that, for all positive operators $x_+\in\mM$,
\begin{align}
    \mE(x_+)\geq \lambda^{-1}x_+\ .
\end{align}
If no finite bound exists, $\operatorname{Ind}(\mE)=\infty$. For an inclusion of factors, the minimal index is
\begin{align}
    [\mM:\mN]_0=\inf_\mE \operatorname{Ind}(\mE)\ .
\end{align}
The index is a measure of the relative size of $\mM$ to $\mN$. When the relative commutant is infinite dimensional, or there are no normal conditional expectations, we say the index is infinite.

Finite index does not by itself specify a matrix relative commutant or conormality, and infinite index does not force the relative commutant to be infinite-dimensional.

\paragraph{The two-sided canonical tower in the common-cyclic case.}
Assume now that $\Omega$ is also cyclic and separating for $\mN$
on $\mH_{\mM}$. The characteristic function of
Definition~\ref{def:subsystem-flow} has the upper-strip endpoint
$D_{\mN\subset\mM}(i/2)=J_{\mM}J_{\mN}$. Set
\begin{equation}
 \Gamma:=D_{\mN\subset\mM}(i/2)^*
 =J_{\mN}J_{\mM},\qquad \Gamma\Omega=\Omega.
 \label{eq:interpolation-Gamma}
\end{equation}
Then the canonical endomorphism
\begin{equation}
 \gamma_{\rm can}(m):=\Gamma m\Gamma^*,\qquad
 \gamma_{\rm can}(\mM)=J_{\mN}\mM'J_{\mN}\subseteq\mN
 \label{eq:interpolation-canonical-endomorphism}
\end{equation}
produces a \emph{two-sided} tower of genuine von Neumann algebras.
For each $k\in\mathbb Z$, put
\begin{equation}
 \mA_{2k}:=\Gamma^k\mM\Gamma^{-k},\qquad
 \mA_{2k+1}:=\Gamma^k\mN\Gamma^{-k}.
 \label{eq:interpolation-tower-levels}
\end{equation}
The original inclusion is $\mA_1=\mN\subset\mA_0=\mM$,
and $\mA_{j+1}\subseteq\mA_j$ for every integer $j$:
\begin{equation}
 \boxed{\quad
 \cdots\supset\Gamma^{-1}\mM\Gamma
 \supset\Gamma^{-1}\mN\Gamma
 \supset\mM\supset\mN
 \supset\Gamma\mM\Gamma^*\supset\Gamma\mN\Gamma^*
 \supset\cdots .\quad}
 \label{eq:interpolation-two-sided-tower}
\end{equation}
The left side is a tower of \emph{extensions} of $\mM$,
not merely the commutants of the descending tunnel. Its first extension
is
\begin{equation}
 \mA_{-1}=\Gamma^{-1}\mN\Gamma
 =J_{\mM}\mN'J_{\mM}\supseteq\mM.
 \label{eq:interpolation-first-extension}
\end{equation}
Equivalently, applying commutants to the descending tower gives the
dual ascending chain
$\mM'\subseteq\mN'\subseteq(\Gamma\mM\Gamma^*)'
\subseteq\cdots$; the inversely shifted original tower gives the
extensions in \eqref{eq:interpolation-two-sided-tower}. This is the
canonical-endomorphism or Jones-tunnel construction familiar from
subfactor theory~\cite{revolutionizing,GuidoLongoWiesbrock1998}.

Every $\mA_j$ has $\Omega$ as a cyclic and separating vector.
Consequently, each neighboring pair has its own interpolation tunnel,
\begin{equation}
 \boxed{\quad
 \mI_{j,T}:=
 \bigvee_{|s|<T}
 \sigma_s^{\omega;\mA_j}(\mA_{j+1}),
 \qquad j\in\mathbb Z,\quad T>0.
 \quad}
 \label{eq:interpolation-tower-net}
\end{equation}
These algebras satisfy
\begin{equation}
 \mA_{j+1}\subseteq\mI_{j,T_1}
 \subseteq\mI_{j,T_2}\subseteq\mA_j,
 \qquad T_1<T_2,
 \label{eq:interpolation-tower-net-order}
\end{equation}
as well as $\mI_{j+1,S}\subseteq\mI_{j,T}$ for all $S,T>0$.
Thus the discrete depth $j$ and the continuous modular-history size $T$
organize a partially ordered net of von Neumann algebras on both
sides of the original inclusion. Since $\Gamma\Omega=\Omega$,
\begin{equation}
 \boxed{\quad
 \Gamma\mA_j\Gamma^*=\mA_{j+2},\qquad
 \Gamma\mI_{j,T}\Gamma^*=\mI_{j+2,T}.
 \quad}
 \label{eq:interpolation-shift-covariance}
\end{equation}
The shift by two is exact, but it is a \emph{discrete} covariance;
it does not make the endomorphism the continuous flow of a positive
generator. In half-sided modular inclusions a continuous geometric
interpolation is available, whereas for a general inclusion none is
implied by the tower alone.

\paragraph{\NL{Compression and the non-common-cyclic case.}}

\NL{For a general inclusion, use the coisometry $\Gamma_c=D(i/2)^*$ and the Petz dual $\mathsf E_\omega(m)=\Gamma_c m\Gamma_c^*$ of \eqref{eq:QEC-Petz-dual}. When $V=1$, $\Gamma_c=\Gamma$ and this is the canonical endomorphism; otherwise the map need not be multiplicative.}


The generalized expectation does not necessarily fix the retained
algebra. In fact, its perfectly recovered observables are exactly
\begin{equation}
 \{n\in\mN:\mathsf E_\omega(n)=n\}
 =\mN_{\rm core}^{\,\mM};
 \label{eq:interpolation-Petz-fixed-core}
\end{equation}
thus $\mathsf E_\omega|_{\mN}=\mathrm{id}_{\mN}$ precisely in the
modularly autonomous, state-preserving conditional-expectation case.
This connects the tower to the exact-QEC discussion of
Section~\ref{sec:QEC}. Away from exact correctability,
$\mathsf E_\omega$ is a canonical recovery \emph{candidate} for the
reference state, but approximate QEC requires separate quantitative
control of the recovery error or relative-entropy loss. The existence
of this map alone supplies no approximation bound.

There is a natural first extension even without common cyclicity:
\begin{equation}
 \mM_1:=J_{\mM}\mN'J_{\mM}\supseteq\mM,\qquad
 p:=\Gamma_c^*\Gamma_c=J_{\mM}e_{\mN}J_{\mM}\in\mM_1.
 \label{eq:interpolation-compressed-extension}
\end{equation}
Here $p$ is the support of $\omega$ on $\mM_1$. For
$e_{\mN}<1$, the state is not faithful on $\mM_1$, so its
modular conjugation does not automatically continue the two-sided
tower on the original Hilbert space. Compressing by $p$ instead gives
\begin{equation}
 \boxed{\quad
 \Gamma_c(p\mM p)\Gamma_c^*=\mathsf E_\omega(\mM)
 \ \subseteq\ 
 \Gamma_c(p\mM_1p)\Gamma_c^*=\mN .\quad}
 \label{eq:interpolation-compressed-system}
\end{equation}
The first term is a unital self-adjoint linear space, not generally a
von Neumann subalgebra; its ultraweak closure is an\emph{operator system}. These closures and the further images
under $\mathsf E_\omega$ form the descending order
\begin{equation}
 \mN\supseteq
 \overline{\mathsf E_\omega(\mM)}^{\,\sigma\mathrm w}
 \supseteq\overline{\mathsf E_\omega(\mN)}^{\,\sigma\mathrm w}
 \supseteq\overline{\mathsf E_\omega^{\,2}(\mM)}^{\,\sigma\mathrm w}
 \supseteq\cdots.
 \label{eq:interpolation-operator-systems}
\end{equation}
One may take the von Neumann algebras generated by these systems to
obtain a nested algebraic envelope inside $\mN$, but this envelope
need not have the multiplicative or standardness properties of the
Jones tower. Likewise, the compressed interpolation tunnel
\begin{equation}
 \mathsf I_T:=\overline{V^*\mI_{T,\mN\subset\mM}V}^{\,\sigma\mathrm w}
 \label{eq:interpolation-compressed-net}
\end{equation}
is a nested family of operator systems on $\mH_{\mN}$.
Hence every faithful inclusion has a canonical \emph{modular net of
von Neumann algebras} $\mI_T$, and a canonical first extension;
the full two-sided, standard von Neumann tower requires common
cyclicity (or additional input), while compression naturally leads
to operator systems. These qualifications will matter when discussing
the action of inclusion parent Hamiltonians.

\subsubsection{\NL{Covariance under inclusion parent Hamiltonians}}

We now ask how the inclusion parent Hamiltonian acts on the
interpolation tunnel and on the ascending and descending canonical
\NL{towers just constructed.} The key distinction
is between preserving the \emph{inclusion relations of an evolving net}
and preserving the \emph{original, fixed algebras} of that net.
The first statement is exact for the inclusion parent Hamiltonian,
but is also shared by any unitary evolution. The second is highly
restrictive and, in a two-sided form, is incompatible with nonzero
positive energy when the algebras are standard.

\begin{proposition}[Covariance of the interpolation net]
\label{prop:Q-interpolation-covariance}
Assume $e_{\mN}=1$ and $s(\rho_\alpha)=1$, and write
$U_t=e^{itQ_{\alpha,\mN\subset\mM}}$ and
$\gamma_t=\operatorname{Ad}U_t$. Define the transported pair
$\mN(t)=\gamma_t(\mN)\subseteq\mM(t)=\gamma_t(\mM)$.
Then, for all $t\in\mathbb R$ and $T>0$,
\begin{equation}
 \boxed{\quad
 \gamma_t(\mI_{T,\mN\subset\mM})
 =\mI_{T,\mN(t)\subset\mM(t)}.
 \quad}
 \label{eq:Q-interpolation-covariance}
\end{equation}
Let $\mA_j$ and $\mI_{j,T}$ be the two-sided canonical tower
and its interpolations. The corresponding tower of the evolved
inclusion satisfies
\begin{align}
 \Gamma(t)&:=J_{\mN(t)}J_{\mM(t)}=U_t\Gamma U_t^*,
 \label{eq:Q-evolved-Gamma}\\
 \mA_j(t)&=\gamma_t(\mA_j),&
 \mI_{j,T}(t)&=\gamma_t(\mI_{j,T}).
 \label{eq:Q-evolved-tower}
\end{align}
In particular, the full nesting and the $T$-filtration hold at every
time:
\begin{equation}
 \boxed{\quad
 \gamma_t(\mA_{j+1})\subseteq
 \gamma_t(\mI_{j,T_1})\subseteq
 \gamma_t(\mI_{j,T_2})\subseteq
 \gamma_t(\mA_j),\quad T_1<T_2.
 \quad}
 \label{eq:Q-preserves-all-inclusions}
\end{equation}
\end{proposition}
\begin{proof}
Since $Q_\alpha\Omega=0$, $U_t\Omega=\Omega$. For every
standard algebra $\mathcal C\subseteq B(\mH)$,
Tomita--Takesaki theory gives
\begin{equation}
 J_{\gamma_t(\mathcal C)}=U_tJ_{\mathcal C}U_t^*,
 \qquad
 \Delta_{\gamma_t(\mathcal C)}=U_t\Delta_{\mathcal C}U_t^*.
 \label{eq:Q-Tomita-transport}
\end{equation}
Conjugating the definition \eqref{eq:interpolation-tunnel}
with $U_t$ proves \eqref{eq:Q-interpolation-covariance}.
The first identity in \eqref{eq:Q-evolved-Gamma} follows from
modular-conjugation covariance; its powers give the transported
tower levels. Repeating the argument for the adjacent inclusion
$\mA_{j+1}\subset\mA_j$ proves the second identity in
\eqref{eq:Q-evolved-tower}. Finally, unitary conjugation preserves
every inclusion, giving \eqref{eq:Q-preserves-all-inclusions}.
\end{proof}

The proposition includes the tower of extensions beyond $\mM$.
For example,
\begin{equation}
 \gamma_t(\Gamma\mM\Gamma^*)
 \subseteq\gamma_t(\mN)
 \subseteq\gamma_t(\mI_{T,\mN\subset\mM})
 \subseteq\gamma_t(\mM)
 \subseteq\gamma_t(\Gamma^{-1}\mN\Gamma).
 \label{eq:Q-extension-chain-covariance}
\end{equation}
The same holds at arbitrary tower depth and for its intermediate
modular-history algebras. We emphasize that $\Gamma(t)$ is the
canonical unitary of the \emph{evolved} pair; we have not asserted
$[Q_{1/4},\Gamma]=0$ or
$\gamma_t(\mA_j)=\mA_j$. The midpoint identity
$J_{\mN}Q_{1/4}J_{\mN}=Q_{1/4}$ alone does not imply either condition.

\paragraph{What remains true under compression.}
For $e_\mN<1$, $\gamma_t$ is an automorphism of
$B(\mH_\mN)$ and does not intrinsically act on the ambient
algebras $\mI_T\subseteq\mM\subseteq B(\mH_\mM)$.
One convenient (but not physically distinguished) extension is
\begin{equation}
 \widetilde U_t:=VU_tV^*+(1-e_\mN),
 \qquad \widetilde\gamma_t:=\operatorname{Ad}\widetilde U_t.
 \label{eq:Q-ambient-extended-flow}
\end{equation}
It fixes $\Omega$ and preserves the embedded space $V\mH_\mN$.
The analogue of \eqref{eq:Q-interpolation-covariance} is
\begin{equation}
 \widetilde\gamma_t(\mI_{T,\mN\subset\mM})
 =\mI_{T,\widetilde\gamma_t(\mN)\subset
              \widetilde\gamma_t(\mM)}.
 \label{eq:Q-compressed-tunnel-transport}
\end{equation}
For the operator systems $\mathsf I_T$ of
\eqref{eq:interpolation-compressed-net} this reads
$\gamma_t(\mathsf I_T)=\mathsf I_T(t)$, with the latter obtained
by compressing the transported ambient tunnel. The Petz-dual map
is covariant as well:
\begin{equation}
 \mE_{\omega,t}\bigl(\widetilde\gamma_t(m)\bigr)
 =\gamma_t\bigl(\mE_\omega(m)\bigr),
 \qquad m\in\mM,
 \label{eq:Q-Petz-covariance}
\end{equation}
where $\mE_{\omega,t}$ belongs to the transported inclusion.
Thus all order relations of the compressed operator-system tower
are transported, but not generally fixed, by $\gamma_t$.

\paragraph{Why simultaneous covariance is not yet locality.}
Equation~\eqref{eq:Q-preserves-all-inclusions} follows for
\emph{any} unitary conjugation and therefore does not distinguish
local from nonlocal dynamics. Moreover, the stronger demand of
invariance of every original tower level would be inappropriate
for a nontrivial positive-energy flow:

\begin{lemma}[No-go theorem for two-sided invariance]
\label{lem:Q-fixed-standard-no-go}
Let $U_t=e^{itQ}$ on $\mH$, where $Q\geq0$ and
$U_t\Omega=\Omega$. If $\mathcal C\subseteq B(\mH)$ has
$\Omega$ cyclic and separating and
$U_t\mathcal C U_t^*=\mathcal C$ for every $t\in\mathbb R$,
then $Q=0$.
\end{lemma}
\begin{proof}
The Tomita operator of $(\mathcal C,\Omega)$ commutes with
$U_t$, and therefore $J_{\mathcal C}U_tJ_{\mathcal C}=U_t$.
Because $J_{\mathcal C}$ is antiunitary, this implies
$J_{\mathcal C}QJ_{\mathcal C}=-Q$. Both $Q$ and
$J_{\mathcal C}QJ_{\mathcal C}$ are positive, so $Q=0$.
\end{proof}

In the common-cyclic case every $\mA_j$ and $\mI_{j,T}$ is
standard for $\Omega$. Hence for nonzero $Q_{1/4}$ none can be
invariant under the entire two-sided flow $\gamma_t$.
This obstruction does \emph{not} exclude one-sided inclusions such
as $\gamma_t(\mA_j)\subseteq\mA_j$ for $t\geq0$.
Such inclusions are realized in half-sided modular examples.
For a chiral net with $\mA_j=\mathcal A((j,\infty))$ and
$Q_{1/4}=P\geq0$, the translation flow gives
\begin{equation}
 \gamma_t(\mA_j)=\mathcal A((j+t,\infty))
 \subseteq\mA_j,\qquad t\geq0,
 \label{eq:Q-half-sided-tower-locality}
\end{equation}
and the interpolations are explicitly
\begin{equation}
 \mI_{j,T}=\mathcal A((j+e^{-2\pi T},\infty)),\qquad
 \gamma_t(\mI_{j,T})\subseteq\mI_{j,T}\quad(t\geq0).
 \label{eq:Q-half-sided-interpolation-locality}
\end{equation}
This is an actual endomorphism property of the \emph{fixed}
net, not merely covariance of an evolving net.

These observations suggest a more substantial, purely algebraic
\emph{tower-locality} condition: for each finite time interval,
there should exist controlled functions $r(t)\in\mathbb Z_{\geq0}$
and $v(t)\geq0$ such that, for the \emph{fixed} reference tower,
\begin{equation}
 \gamma_t(\mI_{j,T})
 \subseteq\mI_{j-r(t),\,T+v(t)}
 \qquad\text{for all }j,T,
 \label{eq:Q-proposed-tower-locality}
\end{equation}
with corresponding bounds for the inverse flow.
A one-sided condition with $r=v=0$ is especially strong.
Equation~\eqref{eq:Q-proposed-tower-locality} is a proposed
\emph{additional} propagation axiom, not a consequence of the
construction of $Q_{1/4}$. It is analogous to finite-speed
propagation or a Lieb--Robinson bound, with inclusion depth and
modular-history width replacing geometric distance; it does not,
by itself, supply spacelike commutativity or a spacetime metric.
The modular flow of the ambient algebra already satisfies the
exact finite-history bound
\eqref{eq:interpolation-modular-propagation}, whereas determining
whether the inclusion parent Hamiltonian enjoys a corresponding
bound is a new dynamical problem.

\subsection{Specialization to common cyclic and separating vectors}
\label{subsec:probing-common-cyclic}

The construction simplifies considerably when $\Omega$ is cyclic and
separating for both $\mN$ and $\mM$ on the same Hilbert space.
Then
\begin{equation}
    \mH_{\mN}=\mH_{\mM}=\mathcal H,
    \qquad
    e_{\mN}=1,
    \qquad
    V=1,
\end{equation}
and the analytic map $D(z)$ becomes Borchers' characteristic function
\begin{equation}
    \boxed{
    D_{\mN\subset\mM}(z)
    =
    \Delta_{\mM}^{-iz}\Delta_{\mN}^{iz}
    :
    \mathcal H\to\mathcal H .
    }
    \label{eq:D-common-probing}
\end{equation}
It is bounded and analytic in
$0<\operatorname{Im}z<1/2$, contractive on the closed strip, and unitary on
both boundaries~\cite{revolutionizing}. The distinction between compressed
and uncompressed subsystem flow disappears.

At the upper endpoint,
\begin{equation}
    \boxed{
    D(i/2)=J_{\mM}J_{\mN} .
    }
    \label{eq:D-half-common}
\end{equation}
Thus
\begin{equation}
    \Gamma:=D(i/2)^*=J_{\mN}J_{\mM}
\end{equation}
is a unitary, and
\begin{equation}
    \gamma_{\rm can}(m)
    :=
    \Gamma m\Gamma^*,
    \qquad
    m\in\mM,
    \label{eq:canonical-endomorphism-common}
\end{equation}
maps $\mM$ into $\mN$:
\begin{equation}
    \gamma_{\rm can}(\mM)
    =
    J_{\mN}\mM'J_{\mN}
    \subseteq
    J_{\mN}\mN'J_{\mN}
    =
    \mN .
\end{equation}
Iteration gives the descending tower, introduced with its extensions
and modular-history interpolations in
Subsection~\ref{subsec:interpolation-tunnel}:
\begin{equation}
    \boxed{
    \mM
    \supset
    \mN
    \supset
    \gamma_{\rm can}(\mM)
    \supset
    \gamma_{\rm can}(\mN)
    \supset
    \gamma_{\rm can}^2(\mM)
    \supset\cdots .
    }
    \label{eq:canonical-discrete-tower}
\end{equation}
Together with its inverse images it is the two-sided $\mathbb Z$-tower
\eqref{eq:interpolation-two-sided-tower}; its adjacent pairs have the
interpolation tunnels \eqref{eq:interpolation-tower-net}.
One might be tempted to define a positive generator by taking the logarithm of $\Gamma$, but there are branch-cut ambiguities in taking the logarithm of a unitary operator. 
Even if we choose a prescription with a positive generator, we are not guaranteed that the unitary flow with this generator will result in a half-sided inclusion that fills in between steps of (\ref{eq:canonical-discrete-tower}) to give a half-sided translation. Indeed, covariance of a tower under the parent-Hamiltonian flow
does not imply that the tower is fixed or translated by that flow
(Proposition~\ref{prop:Q-interpolation-covariance}).

\paragraph{The dual inclusion and the two logarithmic orderings:} Common cyclicity also makes the commutant-dual structure completely explicit.
Taking commutants reverses the inclusion,
\begin{equation}
    \mM'\subset\mN',
\end{equation}
and
\begin{equation}
    \Delta_{\mM'}=\Delta_{\mM}^{-1},
    \qquad
    \Delta_{\mN'}=\Delta_{\mN}^{-1}.
\end{equation}
The characteristic function of the dual inclusion is therefore
\begin{align}
    D^\vee(z)
    &:=
    D_{\mM'\subset\mN'}(z)
    \nonumber\\
    &=
    \Delta_{\mN'}^{-iz}
    \Delta_{\mM'}^{iz}
    \nonumber\\
    &=
    \Delta_{\mN}^{iz}
    \Delta_{\mM}^{-iz}.
    \label{eq:D-dual-common}
\end{align}
On the real axis,
\begin{equation}
    D^\vee(t)=D(-t)^{-1}.
    \label{eq:D-dual-real}
\end{equation}
It is therefore natural to use
\begin{equation}
    D(-z)^{-1}
\end{equation}
as shorthand for the analytic characteristic function $D^\vee(z)$ of the
dual inclusion, rather than as an inverse obtained by analytically
continuing $D(-z)$ through the lower half-plane.

There are two possible logarithmic orderings:
\begin{align}
    G_z^{\mN\subset\mM}
    &:=
    \frac{1}{2\sin(2\pi iz)}
    \log\left[
        D^\vee(z)D(z)
    \right]
    \nonumber\\
    &=
    \frac{1}{2\sin(2\pi iz)}
    \log\left[
        D(-z)^{-1}D(z)
    \right],
    \label{eq:G-original-common}
    \\
    G_z^{\mM'\subset\mN'}
    &:=
    \frac{1}{2\sin(2\pi iz)}
    \log\left[
        D(z)D^\vee(z)
    \right]
    \nonumber\\
    &=
    \frac{1}{2\sin(2\pi iz)}
    \log\left[
        D(z)D(-z)^{-1}
    \right].
    \label{eq:G-dual-common}
\end{align}
Thus the ordering
\begin{equation}
    D(-z)^{-1}D(z)
\end{equation}
is naturally associated with the original inclusion
$\mN\subset\mM$, while
\begin{equation}
    D(z)D(-z)^{-1}
\end{equation}
is naturally associated with the commutant-dual inclusion
$\mM'\subset\mN'$.

Whenever $D(z)$ is invertible and the logarithmic branches are chosen
compatibly, the two products are similar:
\begin{equation}
    D(z)D^\vee(z)
    =
    D(z)
    \left[D^\vee(z)D(z)\right]
    D(z)^{-1}.
\end{equation}
Hence
\begin{equation}
    G_z^{\mM'\subset\mN'}
    =
    D(z)
    G_z^{\mN\subset\mM}
    D(z)^{-1}.
    \label{eq:G-similarity-common}
\end{equation}
They therefore carry the same spectral data, although for generic complex
$z$ the relation is a similarity rather than a unitary equivalence.

\paragraph{Boundary ambiguity.}
On the real boundary, $D(t)$ and $D^\vee(t)$ are unitary, so the products in
\eqref{eq:G-original-common} and \eqref{eq:G-dual-common} are unitary.
Their logarithms therefore have the usual winding ambiguity:
\begin{equation}
    \log U
    \longrightarrow
    \log U+2\pi i
    \sum_k n_kP_k,
    \qquad
    n_k\in\mathbb Z,
\end{equation}
on the spectral projections $P_k$ of the unitary. Since
\begin{equation}
    \sin(2\pi it)
    =
    i\sinh(2\pi t),
\end{equation}
a choice of anti-self-adjoint logarithm gives a self-adjoint $G_t$, but
there is no canonical branch.

The same ambiguity occurs on the upper boundary
$z=t+i/2$, where the characteristic functions are again unitary. At the
special point $z=i/2$ the denominator of
\eqref{eq:G-original-common}--\eqref{eq:G-dual-common} vanishes, so the
logarithmic generators themselves are not defined by direct substitution.
The meaningful endpoint datum is instead the unitary
$D(i/2)=J_{\mM}J_{\mN}$, which gives the discrete tower
\eqref{eq:canonical-discrete-tower}.

\paragraph{Imaginary modular time.}
For
\begin{equation}
    z=i\alpha,
    \qquad
    0<\alpha<\frac12,
\end{equation}
set
\begin{equation}
    X_\alpha
    :=
    D(i\alpha)
    =
    \overline{
    \Delta_{\mM}^\alpha
    \Delta_{\mN}^{-\alpha}} .
    \label{eq:X-alpha-common}
\end{equation}
The dual characteristic function is
\begin{equation}
    D^\vee(i\alpha)
    =
    \overline{
    \Delta_{\mN}^{-\alpha}
    \Delta_{\mM}^\alpha}
    =
    X_\alpha^* .
    \label{eq:Ddual-alpha-common}
\end{equation}
Therefore
\begin{align}
    D^\vee(i\alpha)D(i\alpha)
    &=
    X_\alpha^*X_\alpha,
    \label{eq:XX-original}
    \\
    D(i\alpha)D^\vee(i\alpha)
    &=
    X_\alpha X_\alpha^* .
    \label{eq:XX-dual}
\end{align}
Both are positive contractions. Since
\begin{equation}
    \sin(2\pi i\,i\alpha)
    =
    -\sin(2\pi\alpha),
\end{equation}
the logarithms now have a unique real spectral branch, and
\begin{align}
    \boxed{
    Q_{\alpha,\mN\subset\mM}
    =
    G_{i\alpha}^{\mN\subset\mM}
    =
    -\frac{1}{2\sin(2\pi\alpha)}
    \log(X_\alpha^*X_\alpha)
    \geq0,
    }
    \label{eq:Halpha-original-common}
    \\
    \boxed{
    Q_{\alpha,\mM'\subset\mN'}
    =
    G_{i\alpha}^{\mM'\subset\mN'}
    =
    -\frac{1}{2\sin(2\pi\alpha)}
    \log(X_\alpha X_\alpha^*)
    \geq0.
    }
    \label{eq:Halpha-dual-common}
\end{align}
The first is precisely the specialization of
\eqref{eq:H-alpha-general}. Thus the general definition based on
$D_\alpha^*D_\alpha$ already selects the ordering naturally associated with
the original inclusion.

If
\begin{equation}
    X_\alpha
    =
    U_\alpha|X_\alpha|
\end{equation}
is the polar decomposition, then on the supports
\begin{equation}
    X_\alpha X_\alpha^*
    =
    U_\alpha
    (X_\alpha^*X_\alpha)
    U_\alpha^*,
\end{equation}
and hence
\begin{equation}
    \boxed{
    Q_{\alpha,\mM'\subset\mN'}
    =
    U_\alpha
    Q_{\alpha,\mN\subset\mM}
    U_\alpha^* .
    }
    \label{eq:H-dual-polar}
\end{equation}
The inclusion and its commutant dual therefore have identical nonzero
inclusion spectra.

The geometric half-sided specialization, in which this positive operator
coincides with the Borchers translation generator, is discussed with explicit
QFT examples in the companion paper~\cite{companionQFT}.

\subsection{The canonical endomorphism and its symmetric $L^2$ transfer operator}
\label{subsec:emergent-transfer-time}

For a common cyclic and separating vector, the adjoint of the upper-strip
endpoint is the unitary $\Gamma=J_{\mN}J_{\mM}$, and the Petz dual is the
canonical endomorphism $\gamma_{\rm can}=\operatorname{Ad}\Gamma$.
By \eqref{eq:interpolation-shift-covariance}, this endomorphism shifts the
canonical tower and its interpolation net by two levels. It is therefore
natural to view $\Gamma$ as a discrete evolution preserving the nesting of
these observable algebras. This notion of locality is algebraic: it does
not imply the existence of a continuous local translation group.

\paragraph{Mixing of the discrete canonical evolution.}
For standard inclusions of factors that are also normal and conormal,
Longo proved that the canonical endomorphism is mixing
\cite{Longo1984FactorialStoneWeierstrass,Longo1985Pseudonormalcy}:
\begin{equation}
    \gamma_{\rm can}^{\,n}(m)
    =\Gamma^n m\Gamma^{-n}
    \xrightarrow[n\to\infty]{\rm WOT}
    \omega(m)1,\qquad m\in\mM.
    \label{eq:Longo-canonical-mixing}
\end{equation}
Here normality means $\mN=(\mN'\cap\mM)'\cap\mM$, while
conormality is $\mN\vee(\mN'\cap\mM)=\mM$. Since $\mM\Omega$
is dense and $\Gamma\Omega=\Omega$, the limit implies
\begin{equation}
    \Gamma^n\xrightarrow[n\to\infty]{\rm WOT}
    P_\Omega:=|\Omega\rangle\langle\Omega|.
    \label{eq:Longo-Gamma-weak}
\end{equation}
The projection here is onto the reference vector, whereas
\eqref{eq:Longo-canonical-mixing} converges to a multiple of the
\emph{identity observable}. The normality and conormality hypotheses
in the cited result should not be replaced by common cyclicity alone.
The same distinction between mixing and ergodicity of the canonical
endomorphism is relevant to the more general examples of
\cite{LongoTanimotoUeda2019}.

\paragraph{The symmetric transfer operator of the canonical tower.}
The midpoint $\alpha=1/4$ gives a distinguished positive representation
of this discrete endomorphism. Using the symmetric $L^2$ embeddings
$j_{\mA}(a)=\Delta_{\mA}^{1/4}a\Omega$ of
\eqref{eq:Q-symmetric-embedding}, the inclusion contraction and its
adjoint satisfy
\begin{equation}
 \begin{aligned}
 D_{1/4}\,j_{\mN}(a)&=j_{\mM}(a),&&a\in\mN,\\
 D_{1/4}^*j_{\mM}(m)&=j_{\mN}(\gamma_{\rm can}(m)),&&m\in\mM.
 \end{aligned}
 \label{eq:canonical-L2-adjoint}
\end{equation}
Thus $\rho_{1/4}=D_{1/4}^*D_{1/4}$ is the positive symmetric
$L^2$ transfer operator of the canonical endomorphism:
\begin{equation}
 \boxed{\quad
 e^{-2Q_{1/4}}j_{\mN}(a)
 =\rho_{1/4}j_{\mN}(a)
 =j_{\mN}(\Gamma a\Gamma^*),\qquad a\in\mN.\quad}
 \label{eq:canonical-L2-transfer}
\end{equation}
In particular, $\rho_{1/4}^{\,n}j_{\mN}(a)
=j_{\mN}(\Gamma^n a\Gamma^{-n})$ for $n\geq0$; at each discrete
Euclidean step it advances the canonical tower by two levels in this
representation. In contrast with the ambiguous logarithm of the unitary
$\Gamma$, $Q_{1/4}=-\tfrac12\log\rho_{1/4}$ is a canonically defined
positive self-adjoint generator when $\rho_{1/4}$ has full support.
The positive operator $\rho_{1/4}$ is therefore a canonical transfer operator for the discrete tower, while $Q_{1/4}$ gives a positive Hamiltonian on its support. When it has full support, the corresponding unitary group $e^{itQ_{1/4}}$ is well defined on the same Hilbert space. However, real-time conjugation by $e^{itQ_{1/4}}$ need not move the fixed canonical tower through its discrete levels; in particular,
\begin{equation}
 \Gamma\neq e^{2iQ_{1/4}}\qquad\text{in general}.
 \label{eq:canonical-not-unitary-interpolation}
\end{equation}
The additional properties that identify these unitary dynamics with geometric translations in special half-sided modular inclusions are reviewed and calculated in the companion QFT paper~\cite{companionQFT}.

\subsection{Three-lines majorization and the midpoint gap}
\label{subsec:Q-three-lines}

The reflection-symmetric midpoint has a second distinction besides
$J_\mN Q_{1/4}J_\mN=Q_{1/4}$. It maximizes the spectral
attenuation of the \emph{unnormalized} inclusion operator, as follows.

\begin{theorem}[Three-lines log-majorization]
\label{thm:Q-three-lines}
For each integer $1\leq k\leq\dim\mH_{\mN}$, define
\begin{equation}
 f_k(\alpha)
 :=\log\left\|\bigwedge^kD_\alpha\right\|,
 \qquad 0\leq\alpha\leq\tfrac12.
 \label{eq:Q-fk-definition}
\end{equation}
Then $f_k$ is convex, $f_k(0)=f_k(1/2)=0$, and
$f_k(1/2-\alpha)=f_k(\alpha)$. In particular,
\begin{equation}
 \boxed{f_k(\beta)\leq f_k(\alpha),
 \qquad 0<\alpha<\beta\leq\tfrac14.}
 \label{eq:Q-fk-monotone}
\end{equation}
When $D_\alpha$ is compact and has at least $k$ nonzero singular
values $s_1(\alpha)\geq\cdots\geq s_k(\alpha)$,
\begin{equation}
 f_k(\alpha)=\sum_{j=1}^k\log s_j(\alpha).
\end{equation}
Consequently, the sums of the $k$ lowest finite energies of
$F_\alpha:=-\log\rho_\alpha$ and
$Q_\alpha^{\rm gap}:=F_\alpha/2$ increase on $(0,1/4)$.
On the other hand, the alternative normalization
\begin{equation}
 \widehat Q_\alpha:=-\frac{1}{4\pi\alpha}\log\rho_\alpha
 =\frac{\sin(2\pi\alpha)}{2\pi\alpha}Q_\alpha
 \label{eq:Q-alpha-linear-normalization}
\end{equation}
has \emph{decreasing} sums of its lowest $k$ finite energies on
$0<\alpha<1/2$, whenever the relevant spectral levels are defined.
These are spectral statements; they do not imply a Loewner-order
inequality for the operators at different $\alpha$.
\end{theorem}

\begin{proof}
The exterior-power map $z\mapsto\bigwedge^k D(z)$ is bounded
and analytic on the modular strip. The relation
$D(t+i\alpha)=\Delta_\mM^{-it}D_\alpha\Delta_\mN^{it}$
shows that its norm is independent of the real part of $z$.
Hadamard's three-lines theorem \cite{BerghLofstrom1976}
therefore gives the convexity of $f_k$.
Both boundary characteristic maps are isometries, so $f_k(0)=
f_k(1/2)=0$. The strip-reflection identity
$D_{1/2-\alpha}=J_\mM D_\alpha J_\mN$
preserves exterior-power norms and implies the stated symmetry.
Convexity and midpoint symmetry imply \eqref{eq:Q-fk-monotone}.

For compact $D_\alpha$, the exterior-power norm equals the
product of its $k$ largest singular values. The corresponding
eigenvalues of $\rho_\alpha$ are $s_j(\alpha)^2$, hence the
lowest-energy partial sum of $F_\alpha$ is $-2f_k(\alpha)$,
and that of $Q_\alpha^{\rm gap}$ is $-f_k(\alpha)$.
Finally, convexity and $f_k(0)=0$ imply that
$f_k(\alpha)/\alpha$ is nondecreasing. The lowest-energy
partial sum of $\widehat Q_\alpha$ is
$-f_k(\alpha)/(2\pi\alpha)$, and is therefore nonincreasing.
\end{proof}

\paragraph{Maximizing the spectral gap.}
The reference vector is an eigenvector of every $\rho_\alpha$
with eigenvalue one, and $D(z)$ reduces the corresponding
orthogonal complements (as a map from $\mH_\mN$ to $\mH_\mM$).
Define the gap of the gap-normalized operator by
\begin{equation}
 g_{\rm gap}(\alpha)
 :=\inf\sigma\!\left(
 Q_\alpha^{\rm gap}\big|_{\Omega^\perp}\right)
 =-\log\left\|
 D_\alpha\big|_{\Omega^\perp}\right\|.
 \label{eq:Q-midpoint-gap}
\end{equation}
Applying the same three-lines argument to $D(z)$ restricted
to the orthogonal complement of the vacuum gives a concave,
midpoint-symmetric function $g_{\rm gap}(\alpha)$; therefore
\begin{equation}
 \boxed{
 g_{\rm gap}(\tfrac14)
 =\max_{0<\alpha<1/2}g_{\rm gap}(\alpha),
 \qquad Q_{1/4}^{\rm gap}=Q_{1/4}.}
 \label{eq:Q-gap-optimal-midpoint}
\end{equation}
If the invariant core is trivial and $D_{1/4}$ is compact,
then $\Omega$ is its unique eigenvector of singular value one,
so $g_{\rm gap}(1/4)>0$. In particular, midpoint
$L^2$-nuclearity provides a positive gap for the uniquely grounded
parent Hamiltonian. By contrast, the gap of
$Q_\alpha=Q_\alpha^{\rm gap}/\sin(2\pi\alpha)$ is not
in general maximized at the midpoint: the sine normalization was
chosen instead to make the half-sided modular translation generator
independent of $\alpha$. The two criteria meet at $\alpha=1/4$.

\paragraph{A finite-dimensional check of the normalization.}
The sine-normalized gap has no universal monotonic direction in
$\alpha$. For example, let
$\mN=B(\mathbb C^2)\otimes1\subset B(\mathbb C^4)=\mM$
and use the faithful two-qubit state
$\rho_{AB}^{(p)}=(1-p)|\Phi^+\rangle\langle\Phi^+|+p1_4/4$
with $0<p<1$. In its GNS representation, the nonzero excited
energy of $Q_\alpha$ (with multiplicity three) is
\begin{align}
 q_\alpha(p)&=-\frac{\log r_\alpha(p)}{2\sin(2\pi\alpha)},\nn\\
 r_\alpha(p)&=2b+a^{1-2\alpha}b^{2\alpha}
                   +b^{1-2\alpha}a^{2\alpha},
 \quad a=1-\frac{3p}{4},\quad b=\frac p4.
 \label{eq:Q-alpha-two-qubit-counterexample}
\end{align}
One finds $q_{0.1}(1/20)<q_{0.2}(1/20)$ but
$q_{0.1}(1/2)>q_{0.2}(1/2)$. This illustrates why midpoint
optimality is stated for $Q_\alpha^{\rm gap}$, which agrees
with the physically normalized parent Hamiltonian at $\alpha=1/4$,
rather than for the entire sine-normalized family.

\subsection{$L^2$-nuclearity and the canonical split flow}
\label{sec:nuclearity}

We now specialize to the common-cyclic case, so that
$\mH_{\mN}=\mH_{\mM}=: \mathcal H$ and
\begin{equation}
    X_\alpha=D_\alpha
    =\overline{\Delta_{\mM}^\alpha\Delta_{\mN}^{-\alpha}},
    \qquad 0<\alpha<\frac12 .
\end{equation}
Theorem~
ef{thm:Q-kernel-core} shows that core-freeness alone
is enough for a unique ground state, whereas nuclearity adds a spectral
gap. Theorem~
ef{thm:Q-three-lines} explains why the midpoint is
also the optimal point for the gap-normalized modular comparison family.
The standard $L^2$-nuclearity condition is the midpoint trace-class condition
\begin{equation}
    \boxed{\|X_{1/4}\|_1<\infty .}
    \label{eq:L2-nuclearity-alpha}
\end{equation}
It implies modular nuclearity and hence the split property~\cite{BDL2007}. More generally, at any exponent $0<\alpha<1/2$ for which $X_\alpha$ is trace class, the following spectral conclusions hold. In terms of the inclusion parent Hamiltonian,
\begin{equation}
    \rho_{\alpha,\mN\subset\mM}
    =X_\alpha^*X_\alpha
    =e^{-2\sin(2\pi\alpha)Q_{\alpha,\mN\subset\mM}}.
    \label{eq:rho-inclusion-vs-split}
\end{equation}
If $s_k>0$ are the nonzero singular values of $X_\alpha$, then on the support
of $\rho_\alpha$
\begin{equation}
    E_k^{(\alpha)}
    =-\frac{1}{\sin(2\pi\alpha)}\log s_k
\end{equation}
are the eigenvalues of $Q_{\alpha,\mN\subset\mM}$. Thus $L^2$-nuclearity makes
$\rho_\alpha$ compact and gives a discrete finite-energy spectrum for
$Q_\alpha$: every finite eigenvalue has finite multiplicity and the only
possible accumulation is at $+\infty$. If $\rho_\alpha$ has a kernel, that
kernel is the infinite-inclusion-energy sector discussed after
\eqref{eq:H-alpha-general}. At the midpoint the unrescaled convention used in
some of the literature is
\begin{equation}
    -\log\rho_{1/4}=2Q_{1/4,\mN\subset\mM}.
\end{equation}

At the midpoint, $L^2$-nuclearity implies the split property. If, in addition,
the split inclusion is standard, there is a canonical type-I factor
\begin{equation}
    \boxed{
    \mN\subset\mathcal F_{\rm split}\subset\mM
    }
    \label{eq:canonical-split-factor}
\end{equation}
and a canonical split unitary
$U_{\rm split}:\mathcal H\to\mathcal H_L\otimes\mathcal H_R$ such that
\begin{align}
    U_{\rm split}\mN U_{\rm split}^*
    &=\mN_L\otimes1, \nonumber\\
    U_{\rm split}\mathcal F_{\rm split}U_{\rm split}^*
    &=B(\mathcal H_L)\otimes1, \nonumber\\
    U_{\rm split}\mM'U_{\rm split}^*
    &=1\otimes\mM_R'.
    \label{eq:canonical-split-representation}
\end{align}
Writing $|\xi\rangle=U_{\rm split}|\Omega\rangle$, define
\begin{equation}
    \sigma_{\rm split}
    =\Tr_{\mathcal H_R}|\xi\rangle\langle\xi|,
    \qquad
    S_{\rm split}=-\Tr\sigma_{\rm split}\log\sigma_{\rm split}.
    \label{eq:sigma-split-definition}
\end{equation}
This gives a second natural flow associated with the same inclusion. Its inner
implementation on the original Hilbert space is
\begin{equation}
    V_t^{\rm split}
    =U_{\rm split}^*
      (\sigma_{\rm split}^{\,it}\otimes1)
      U_{\rm split}
    \in\mathcal F_{\rm split}.
    \label{eq:split-inner-dynamics}
\end{equation}
Consequently,
\begin{align}
    \Ad V_t^{\rm split}(\mathcal F_{\rm split})
    &=\mathcal F_{\rm split}, \nonumber\\
    \Ad V_t^{\rm split}(\mM)&=\mM, \nonumber\\
    \Ad V_t^{\rm split}(\mN)
    &\subset\mathcal F_{\rm split}\subset\mM.
    \label{eq:split-dynamics-localized}
\end{align}
Thus the split flow keeps the transported subsystem inside the canonical
localized type-I factor. It need not preserve $\mN$ itself; for a
proper common-cyclic inclusion, invariance of $\mN$ under the modular
flow of $\mathcal F_{\rm split}$ would, by Takesaki's theorem, produce an
$\Omega$-preserving expectation $\mathcal F_{\rm split}\to\mN$, which
is incompatible with proper common cyclicity.

This is sharply different from the positive inclusion parent Hamiltonian flow. For
a standard inclusion satisfying the standing full-support condition, Corollary~\ref{cor:standard-Halpha-full-generation}
gives
\begin{equation}
    \bigvee_{t\in\mathbb R}
    \gamma_t^\alpha(\mN)
    =B(\mH_{\mM}).
    \label{eq:Hinc-does-not-preserve-N}
\end{equation}
At any exponent for which $\rho_\alpha$ is trace class with full support,
Proposition~\ref{prop:Halpha-pure-point-one-sided-generation} gives the same
conclusion separately for $t>0$ and $t<0$.
Hence, whenever $\mM\neq B(\mH_{\mM})$, the $Q_\alpha$-orbit must
leave $\mM$ for some times. The two flows therefore probe the inclusion
in complementary ways: $Q_\alpha$ explores the full operator algebra on the
GNS Hilbert space, whereas the split flow is a localized inner dynamics
confined to $\mathcal F_{\rm split}\subset\mM$.

\paragraph{Inclusion entropy.}
Whenever $X_\alpha$ is trace class, the inclusion spectrum can be normalized to a density matrix. Define
\begin{equation}
    \widehat\rho_{\alpha,\mN\subset\mM}
    =\frac{\rho_{\alpha,\mN\subset\mM}}
           {\Tr\rho_{\alpha,\mN\subset\mM}},
    \qquad
    S_{\rm inc}^{(\alpha)}
    =-\Tr\widehat\rho_{\alpha,\mN\subset\mM}
          \log\widehat\rho_{\alpha,\mN\subset\mM}.
    \label{eq:rho-inclusion-normalized}
\end{equation}
Since $X_\alpha$ is trace class,
\begin{equation}
    S_{\rm inc}^{(\alpha)}
    \le
    2\log\|X_\alpha\|_1
    -\log\Tr(X_\alpha^*X_\alpha),
    \label{eq:Sinclusion-nuclearity-bound}
\end{equation}
so $S_{\rm inc}^{(\alpha)}$ is finite. This is not the same entropy as
$S_{\rm split}$: the latter is the entanglement entropy of the canonical split
purification, while $S_{\rm inc}^{(\alpha)}$ is the entropy of the normalized
squared singular values of the modular comparison operator.

\section{Discussion}\label{sec:discussion}

We have developed the subsystem flow as an object of ergodic theory and shown that it is the natural quantitative expression of quantum error correction for the dynamics generated by a state. Classically, every state on a Poisson manifold generates a flow, the relative transformation between the flows of a state and of its coarse-graining is a measure-preserving symplectic cocycle with generator $-\log(\rho/\rho_0)$, and its expectation is the Kullback--Leibler divergence, which is the classical sufficiency gap for the pair consisting of the state and the reference. Every classical coarse-graining is exactly correctable against erasure with respect to the reference, independently of whether the retained observables evolve autonomously. Quantum mechanically the two questions merge: by Takesaki's theorem, the subalgebras that admit a reference-preserving encoding correctable against erasure are exactly the modularly autonomous ones, and we proved that both are equivalent to triviality of the subsystem flow (Theorem~\ref{thm:correctable}). In the standard language of error correction this is a Knill--Laflamme condition with the Jones projection as code projection; for conormal inclusions it is literally the statement that the complementary channel into the relative commutant is forgetful on the code, and the information--disturbance trade-off survives in the modular setting at the level of the algebra but not pair by pair (Section~\ref{subsec:QEC-KL}). When the subsystem flow is not trivial, its generator, evaluated on a state, is the sum of two Petz sufficiency gaps (Proposition~\ref{prop:sufficiency}): the subsystem flow is a generating function for the failure of recoverability, with the second gap, absent classically, coming from the inclusion of commutants.

The modular ergodic hierarchy organizes the ways correctability fails. The history algebra and invariant core of a subalgebra are the smallest correctable algebra containing it and the largest correctable algebra inside it, and their exact duality (Lemma~\ref{lem:history-memory}) shows that a generating inclusion forces the environment to have trivial invariant core, the algebraic modular form of decoupling. Half-sided modular inclusions are the maximally non-correctable case: no expectation preserves the reference, the orbit generates the full algebra, the environment has trivial invariant core, the subsystem-flow generator is Borchers' positive translation generator, and it is transported by the modular flow at the exponent $2\pi$ that saturates the modular chaos bound. At the other end, the midpoint $L^2$-nuclearity condition $X_{1/4}\in S_1$ implies the split property and hence a canonical intermediate type-I factor with a localized split flow. More generally, trace class of $X_\alpha$ at a fixed interior exponent makes the finite spectrum of $Q_{\alpha,\mN\subset\mM}$ discrete with finite multiplicities. At the midpoint, $\|X_{1/4}\|_1=\Tr e^{-Q_{1/4,\mN\subset\mM}}$ in the normalization of \eqref{eq:H-alpha-general}. Modular $K$-systems instead have the continuous positive-energy spectrum characteristic of half-sided modular translations, so the chaotic and nuclear regimes represent sharply different spectral limits of the same imaginary-time construction. This is the sense in which the ergodic hierarchy of a many-body system, formulated subsystem by subsystem, is a hierarchy of correctability.

The results of Section~\ref{sec:probing-outside-M} give an algebraic meaning to the emergence of time in the title. The input is an inclusion $\mN\subset\mM$ and a reference vector $\Omega$; the output is a family of positive Hamiltonians fixing $\Omega$, constructed without an independently specified global Hamiltonian. For a standard inclusion of factors with $s(\rho_\alpha)=1$, the subsystem history under $\gamma_t^\alpha$ generates the entire algebra $B(\mH_{\mM})$ (Corollary~\ref{cor:standard-Halpha-full-generation}). This establishes the global algebraic reach of the dynamics, whose orbit can extend beyond $\mM$. When $Q_\alpha$ has pure point spectrum on the separable Hilbert space, Proposition~\ref{prop:Halpha-pure-point-one-sided-generation} strengthens this statement: either the future or the past history generates the full algebra. Trace class of $\rho_\alpha$ with full support is a sufficient spectral condition, linking this strengthening to the nuclearity regime.

Pure point spectrum enters this argument through recurrence: it identifies each half-history with the complete history. This is a statement about recurrent, reversible evolution and does not imply mixing or a thermodynamic arrow of time. The construction retains its dependence on $\Omega$ and the parameter $\alpha$; selecting a physical clock and establishing a spacetime interpretation require additional structure. Half-sided modular inclusions realize geometric null translations, while the Hamiltonian reconstruction examples of Ref.~\cite{chen2026local} identify inclusion parent Hamiltonians with physical global dynamics in special settings. These examples support the interpretation of the algebraic construction as emergent time. Determining when that identification extends to more general inclusions remains an open problem.

Several distinctions have been kept separate because they are frequently identified. Common cyclicity of the reference vector gives generation and excludes a reference-preserving expectation; it does not give mixing, asymptotic forgetfulness, a decay rate for connected correlations, or a trivial one-sided invariant core. The split property is statistical independence of $\mN$ and $\mM'$ and excludes every normal expectation onto a type $III$ subalgebra, but it is not a dynamical condition. Conormality concerns $\mN$ and $\mR=\mN'\cap\mM$, and together with a faithful normal expectation it forces a spatial tensor factorization, which is what fails for adjacent local algebras. Finite index does not specify a matrix relative commutant, and infinite index does not force an infinite-dimensional one. The absence of a reference-preserving expectation does not exclude expectations preserving other faithful states, nor error correction for other channels or families of states; the general criterion remains Takesaki's.

Several directions remain open. The sufficiency decomposition refers to the generator of the subsystem flow; the full flow $D_{\mN\subset\mM}(z)$ on the strip should admit an interpretation in terms of a one-parameter family of recovery channels interpolating between the two modular conjugations, and its matrix elements should control approximate recoverability in the sense of the Fawzi--Renner and Junge--Renner--Sutter--Wilde--Winter inequalities. When the two modular operators commute, the unrescaled midpoint logarithm $-\log\rho_{1/4}$ equals one half of $K_{\mM}-K_{\mN}$ (equivalently $Q_{1/4,\mN\subset\mM}=\tfrac14(K_{\mM}-K_{\mN})$ in the normalization of \eqref{eq:H-alpha-general}); the relation between its spectrum and the sufficiency gaps in the general case, and the level statistics of that spectrum for interacting theories (the spectral-statistics proposal of this section), are not known. The compressed subsystem flow of Section~\ref{subsec:QEC-takesaki} is a contraction whose conjugation action is completely positive but not in general unital; identifying the structure under which it is a channel on $\mN$ would extend the dictionary beyond the common-cyclic case. For modular $K$-systems the relation between half-sided nesting and a trivial invariant core is not settled by common cyclicity alone, while standardness is sufficient by Lemma~\ref{lem:standard-forgetful}. In the companion paper~\cite{companionQFT} the subsystem flows of wedges and balls in vacuum quantum field theory are computed explicitly, the modular chaos bound is derived from their analytic structure, and both ends of the hierarchy are realized geometrically.

The opposite of an invariant subsystem is one that has no nontrivial modular invariant core. For such inclusions, the positive parent Hamiltonians constructed from the analytic subsystem flow have a unique reference-vector ground state (Theorem~\ref{thm:Q-kernel-core}). Conversely, the ground-state space is precisely the GNS closure of the core. When the inclusion operator has full support, a trivial core is also sufficient for the unitary orbit of the retained algebra to generate the entire algebra $B(\mH_{\mN})$ (Proposition~\ref{prop:Halpha-full-generation}). This result does not require nuclearity; nuclearity and the three-lines theorem add spectral information and stronger one-sided generation statements under their own hypotheses.

The Petz dual and the inclusion parent Hamiltonian describe two aspects of the same modular comparison. On the upper boundary of the analytic strip the Petz dual is a generalized conditional expectation, and universal recovery is obtained by averaging its rotated versions. For common cyclic inclusions the dual becomes the canonical endomorphism $\operatorname{Ad}\Gamma$. Its action on the symmetric $L^2$ space is represented by the positive operator $\rho_{1/4}=e^{-2Q_{1/4}}$; Longo's mixing theorem for the discrete canonical evolution provides a particularly strong example under its additional hypotheses. We do not identify the general unitary group $e^{itQ_{1/4}}$ with a geometric translation. The corresponding QFT and holographic examples, including modular chaos and emergent time, are analyzed in the companion paper~\cite{companionQFT}.

\section*{Acknowledgements}
We thank \ldots

\appendix

\section{Additional operator-algebra background}\label{app:extra-operator-algebra}

\subsection{Review of general quantum systems (von Neumann algebras)}\label{app:vN}

We start by reviewing and fixing our notation for general quantum systems (von Neumann algebras). Readers familiar with standard definitions can skip this subsection.

\paragraph{Algebra of quantum observables:}
Quantum observables are described by a
\emph{von Neumann algebra} $\mM$ a unital $*$-subalgebra of the bounded operators
$B(\mathcal{H})$ on some separable Hilbert space $\mathcal{H}$, closed in the
weak operator topology.
A von Neumann algebra $\mM$ is a \emph{factor} if its center is trivial
$\mathcal{Z}(\mM) = \mM\cap\mM'=\mathbb{C} 1$, where $\mM' = \{b\in B(\mathcal{H}) : [b,a]=0\text{ for all }a\in\mM\}$
is the commutant in $B(\mH)$.  Physically, a factor algebra describes a system with no classical observables
 or superselection sectors \footnote{As we review in Appendix~\ref{sec:classical}, in classical physics, the commutator is replaced by the Poisson bracket and the commutant $\mM'$ is the quantum analog of the Poisson
centralizer introduced in Appendix~\ref{sec:classical}. It consists of all observables invisible to $\mM$. The classical analog of a factor is a non-degenerate Poisson bracket that restricts to a symplectic leaf.}.

\paragraph{Quantum states:}

A \emph{normal state} on $\mM$ is a positive, normalized, $\sigma$-weakly continuous
linear functional $\omega:\mM\to\mathbb{C}$, with $\omega(\id)=1$ \footnote{A normal state is the infinite-dimensional generalization of a standard quantum density matrix on the GNS Hilbert space. Infinite-dimensional von Neumann algebras can also have singular states that are not represented by trace-class density operators. For example, fixing a free ultrafilter $\mathcal U$ on $\mathbb N$, the ultralimit $\omega_{\mathcal U}(a)=\lim_{n\to\mathcal U}\braket{n|a|n}$ defines a singular state on $B(\mathcal H)$ that vanishes on compact operators. Requiring a state to be {\it normal} rules out such states.}.
For a finite matrix algebra it has the form $\omega(a)=\Tr(\rho_\omega a)$,
where $\rho_\omega$ is positive with unit trace \footnote{In certain finite-dimensional algebras (type $III$) there are no traces or intrinsic density matrices; therefore, the normal functional is the general definition.}. A state $\omega$ is \emph{faithful} if $\omega(a^*a)=0$ implies $a=0$.
In finite dimensions, this means that its density matrix has no zero eigenvalues (is full rank) \footnote{Faithfulness is the counterpart of a strictly positive classical probability density
and gives a separating vector in the GNS construction.}. A state $\omega$ is \emph{tracial} if $\omega(ab)=\omega(ba)$ for all $a,b\in\mM$.
\footnote{On a finite algebra, a faithful normal tracial state plays the role of the reference
Liouville measure of classical algebras. Its modular flow is trivial, despite faithfulness. In the type $III$
algebras relevant to QFT, we have faithful normal states but no trace.}

\paragraph{Dynamics, inner and outer:}
A dynamics is a flow, a pointwise $\sigma$-weakly continuous one-parameter group $t\mapsto\alpha_t\in\operatorname{Aut}(\mM)$. The automorphism $\alpha_t$ is \emph{inner} if it is $\operatorname{Ad}u$ for a unitary $u\in\mM$, and the flow is inner if every $\alpha_t$ is. Inner dynamics is generated by a Hamiltonian inside the algebra, and an outer one is generated by something outside the algebra \footnote{On a factor with separable predual an inner flow is always of the form $\alpha_t=\operatorname{Ad}e^{itH}$ with $H$ a self-adjoint operator affiliated with $\mM$, unique up to an additive constant~\cite{Kallman1971}.}. 

\paragraph{GNS Hilbert space and modular operators:}
Given a von Neumann algebra $\mM$ and a faithful normal state $\omega$, the
Gelfand--Naimark--Segal (GNS) construction produces a Hilbert space $\mathcal{H}_\omega$,
a cyclic unit vector $|\Omega\rangle\in\mathcal{H}_\omega$ with
$\omega(a)=\langle\Omega|a|\Omega\rangle$, and a $*$-representation
$\pi:\mM\to B(\mathcal{H}_\omega)$ \footnote{The classical analog is $\mathcal{H}=L^2(X,\mu)$
with the constant function $|\id\rangle$ as the reference vector.}.

The Tomita--Takesaki Modular theory assigns to the pair $(\mM,|\Omega\rangle)$ an
antilinear operator $S_\omega$ defined on the dense domain
$\mM|\Omega\rangle$ by
\begin{equation}
  S_\omega\, a|\Omega\rangle = a^*|\Omega\rangle, \qquad a\in\mM.
  \label{eq:Tomita-S}
\end{equation}
Its polar decomposition $S_\omega = J_\omega\,\Delta_\omega^{1/2}$ defines the
\emph{modular conjugation} $J_\omega$, an antiunitary involution, and the
\emph{modular operator} $\Delta_\omega$, a positive self-adjoint operator. 

\paragraph{Modular dynamics:}
The Tomita--Takesaki theorem~\cite{Takesaki1970} guarantees that the one-parameter group
of unitaries $\Delta_\omega^{it}$, $t\in\mathbb{R}$, implements an automorphism group of
$\mM$:
\begin{equation}
  \sigma_t^\omega(a) = \Delta_\omega^{it}\, a\, \Delta_\omega^{-it},
  \qquad a\in\mM,\quad t\in\mathbb{R}.
  \label{eq:modular-flow}
\end{equation}
This is the \emph{modular flow} (or \emph{modular automorphism group}) of $\omega$.
For matrix algebras, this is $\sigma^{\omega}_t(a)=\rho_\omega^{it}a\rho_\omega^{-it}$.
The modular operator $\Delta_\omega$ and $J_\omega$ encode the non-commutativity of
$\mM$. While in Abelian algebras, modular flow is trivial, in Appendix~\ref{subsec:classical-state-flow}, we argue that the
Poisson flow associated with classical states is the correct replacement of modular flow to formulate subsystem flow and ergodic theory\footnote{This construction uses
the Poisson bracket and differs from the trivial modular automorphism group of
an abelian von Neumann algebra. In both constructions the reference state is stationary:
$\omega\circ\sigma_t^\omega=\omega$.}.

With the convention \eqref{eq:modular-flow}, the \emph{KMS condition} reads
\begin{equation}
 F_{a,b}(t)=\omega\bigl(a\,\sigma_t^\omega(b)\bigr),
 \qquad F_{a,b}(t-i)=\omega\bigl(\sigma_t^\omega(b)a\bigr).
 \label{eq:KMS-function}
\end{equation}
The function is bounded and analytic on $-1<\operatorname{Im}z<0$, with continuous
boundary values. Equivalently, $\beta_t^\omega:=\sigma_{-t}^\omega
=\operatorname{Ad}\Delta_\omega^{-it}$ gives the upper-strip condition at inverse
temperature one. Its generator is $K_\omega=-\log\Delta_\omega$, and on its domain
$\partial_t\sigma_t^\omega(b)|_{t=0}=i[\log\Delta_\omega,b]$.

\paragraph{Relative modular flow and relative entropy:}
For faithful normal states $\omega,\psi$ on a von Neumann algebra $\mathcal C$, represented by vectors $\Omega,\Psi$ in standard form, the relative Tomita operator $S_{\psi|\omega;\mathcal C}$ is defined by
\begin{align}
 S_{\psi|\omega;\mathcal C} a\ket{\Omega}=a^*\ket{\Psi},  \forall a\in\mM
\end{align}
and the relative modular operator is
\begin{align}
 \Delta_{\psi|\omega;\mathcal C}=S_{\psi|\omega;\mathcal C}^\dagger S_{\psi|\omega;\mathcal C},
 \qquad
 \langle\Delta_{\psi|\omega;\mathcal C}^{1/2}a\Omega|\Delta_{\psi|\omega;\mathcal C}^{1/2}b\Omega\rangle=\psi(ba^*),\qquad a,b\in\mathcal C,
 \label{eq:relative-modular}
\end{align}
so that $\Delta_{\omega|\omega;\mathcal C}=\Delta_{\omega;\mathcal C}$. Araki's relative entropy is
\begin{align}
 S_{\mathcal C}(\omega\|\psi)
 =-\langle\Omega|\log\Delta_{\psi|\omega;\mathcal C}|\Omega_\rangle
 =i\left.\frac{d}{dt}\langle\Omega|\Delta_{\psi|\omega;\mathcal C}^{it}|\Omega\rangle\right|_{t=0},
 \label{eq:relative-entropy}
\end{align}
with the derivative understood when the logarithmic spectral moment is integrable~\cite{Araki1976,Araki1977}. In matrix algebras, the relative modular super operator acts as  $\Delta_{\psi|\omega}(b)=\rho_\psi b \rho_\omega^{-1}$ and $S_{\mathcal C}(\omega\|\psi)$ reduces to the usual quantum relative entropy. Relative entropy is nonnegative and monotonic under restriction to a von Neumann subalgebra; these properties underlie the sufficiency gaps used below \cite{Araki1976}.

\subsection{Algebra types and ergodicity of modular flow}\label{subsec:types-modular}

The inner/outer character and the ergodic properties of modular flows distinguish the types of infinite-dimensional factors. To start, we need a few definitions.
The set of operators that are invariant under modular flow (conserved modular charges) forms an algebra of modular charges (fixed point subalgebra)\footnote{The subalgebra of modular charges is called {the centralizer} in mathematical jargon.}. Ergodicity of modular dynamics requires the absence of conserved modular charges (trivial centralizer) \cite{longo1979notes,ouseph2024local}.
In ergodic theory, we also often study the connected modular two-point functions:
\begin{align}
    C^{\omega}_{a,b}(t)=\omega\bigl(a^*\sigma^\omega_t(b)\bigr)-\omega(a^*)\omega(b)\ .\label{eq:modular-correlation}
\end{align}
If there are classical observables (the algebra has a center), they are inside the intersection of modular charges of all states. Since our focus is modular ergodicity, we focus on von Neumann factors so that there are no operators that are conserved under the modular flow of all states. von Neumann factors can be classified as
\begin{enumerate}
    \item {\bf Type I$_\infty$:} Every faithful normal state has a trace-class density matrix with pure-point spectrum. In other words, the modular unitary group is almost periodic and cannot be ergodic. There is a trace; when it is not finite, i.e., $\tr(\mathbb{I})=\infty$, we say the algebra is type I$_\infty$.

    \item {\bf Type II:} Every faithful normal state or weight (unnormalized state) has the form $\omega=\tau(\rho_\omega\,\cdot\,)$ with $\tau$ the tracial weight and $\sigma_t^\omega=\operatorname{Ad}\rho_\omega^{it}$ for a positive operator $\rho_\omega$ affiliated with the algebra $\mM$. The spectral projections of $\rho_\omega$ lie in the centralizer, so a type-II factor has no ergodic modular state. A finite trace exists in type II$_1$ and has trivial modular flow; type II$_\infty$ has a trace, but it is not finite: $\tr(\mathbb{I})=\infty$.
    
    \item {\bf Type III:} Type III algebras admit no trace and no intrinsic density matrices. Within modular ergodic theory, type III$_1$ is the only factor type that can support a faithful normal state with no modular charges (trivial centralizer). Under the standard assumptions of relativistic AQFT, local algebras are typically type III$_1$ factors. The relation between modular mixing, information loss, and emergent type III$_1$ structure was emphasized in~\cite{furuya2023information}.

\end{enumerate}

\subsection{Modular flows versus arbitrary unitary dynamics}\label{sec:modular-vs-dynamics}

By the Stone-von Neumann theorem, every continuous unitary flow has a self-adjoint generator; i.e, $U_t=e^{i H t}$. If $e^{-H}$ is normalizable and belongs to the algebra, the unitary flow is the same the modular flow of the density matrix $\rho=e^{-H}/\tr(e^{-H})$. In this case, the unitary dynamics can be viewed as modular flow of some state\footnote{The modular flow of a state is the unique one-parameter group of automorphisms of $\mM$ for which $\omega$ is KMS at $\beta=1$~\cite{Takesaki1970}.}. In a general von Neumann algebra, the modular flow of any state is an example of dynamics. However, it is no longer true that any dynamics can be viewed as the modular flow of a state \footnote{For example, the observable algebra of an infinite spin chain in the infinite-temperature state is the hyperfinite $\mathrm{II}_1$ factor $R=\bigotimes_n(M_2,\operatorname{tr})$; the global rotation $\alpha_t=\bigotimes_n\operatorname{Ad}e^{it\sigma_z}$ is a continuous flow whose would-be generator, the total magnetization $\sum_n\sigma_z^{(n)}$, is not affiliated with $R$ \footnote{This flow is outer for all $t\notin\pi\mathbb Z$, since $\bigotimes_n\operatorname{Ad}u_n$ is inner iff $\sum_n(1-|\operatorname{tr}u_n|)<\infty$ and here $|\operatorname{tr}e^{it\sigma_z}|=|\cos t|$.}.}. In this subsection, we ask, given some dynamics $\alpha_t$ on an infinite-dimensional quantum system $\mM$, when is it the modular flow of some state? The answer depends on the type of $\mM$, as we explain below.

The first observation is that the modular flow of any state in type I and type II algebras (semifinite algebras) is always inner. 
Therefore, in a type I or II algebra, a unitary flow can be viewed as the modular flow of a weight (unnormalized state) if and only if it is inner. 

\paragraph{Type I:} The modular flow of unnormalized states (weights) includes all unitary flows.
Type I algebras have no outer automorphisms and every flow $\operatorname{Ad}e^{itH}$ is the modular flow of a weight, but in type I$_\infty$ the weight might not be normalizable; e.g., it might be that $e^{-H}=\infty$.

\paragraph{Type II:} The modular flow of unnormalized states (weights) includes all inner unitary flows. 
In type $\mathrm{II}_1$ outer automorphisms exist (e.g., global symmetry rotations of a spin chain or Bernoulli shifts). They do not correspond to any modular flows.
In type $\mathrm{II}_\infty$ algebras, a new kind of outer automorphism appears: the trace-scaling automorphisms $\tau\circ\theta_\lambda=\lambda\tau$, which exist for every $\lambda>0$.

\paragraph{Type III:} The modular flow of weight form a canonical outer class.
No inner flow is modular for all $t$. Every type III factor has outer automorphisms to times $t\notin T(\mM)$, where $T(\mM):=\{t:\sigma^\varphi_t\in\operatorname{Inn}(\mM)\}$ is Connes' $T$-invariant independent of the choice of state. All faithful normal semifinite weights determine the same canonical outer modular flow $\delta_{\mM}(t):=[\sigma^\varphi_t]\in\operatorname{Out}(\mM)$, because any two modular flows differ by an inner cocycle, $\sigma^\psi_t=\operatorname{Ad}[D\psi:D\varphi]_t\circ\sigma^\varphi_t$~\cite{Connes1973,ConnesTakesaki1977}. For $\mathrm{III}_\lambda$, $0<\lambda<1$, one has $T(\mM)=(2\pi/|\log\lambda|)\mathbb Z$, and one may choose faithful normal periodic weights with this modular period \footnote{For $\mathrm{III}_0$, the $T$-invariant is not fixed by the type; the finer classification is encoded by the flow of weights.}. For $\mathrm{III}_1$, $T(\mM)=\{0\}$: $\delta_{\mM}$ is injective, and every modular automorphism of every faithful normal semifinite weight is outer for every $t\ne0$; its continuous core is a type $\mathrm{II}_\infty$ algebra carrying the dual trace-scaling flow. $\mathrm{III}_1$ is also the only type on which a state can have a trivial centralizer, an ergodic modular flow~\cite{Connes1973,Longo1982}, though not every state of a $\mathrm{III}_1$ factor does. For a physicist interested in the thermodynamic or continuum limit, $\mathrm{III}_1$ algebras are important. The local algebras in broad classes of relativistic quantum field theories are type III$_1$. 
    
    Type III algebras allow more general outer dynamics that do not correspond to any modular flows. Physics examples include spacetime or internal symmetries of the theory. For example, translations in perpendicular spatial directions are unitary flows of the wedge algebra that do not correspond to a modular flow. In fact, any geometric symmetry whose generator is not always time-like cannot be generated by modular flow \cite{BenyKempfKribs2007}. Another example is the unitary intertwiners that act as outer automorphisms of the additive algebras in the DHR theory of $0$-form symmetries \cite{casini2020entanglement}. 

\paragraph{Type III$_1$ modular flows can be ergodic and strongly mixing:}

When the algebra $\mM$ has a non-trivial center $Z\subset \mM$, then this center is inside the centralizer of all modular flows:
\begin{align}
    \forall z\in Z,\forall t,\omega:\qquad \Delta_{\omega}^{it}z \Delta_\omega^{-it}=z
\end{align}
This means that these observables are conserved charges, and an obstruction to ergodicity in modular dynamics of arbitrary states. That is why, in quantum ergodic theory, it is natural to restrict to algebras with trivial center, i.e., factors. For a fixed faithful normal state, modular ergodicity means that the centralizer is trivial:
\begin{align}
    \label{eq:modular-ergodicity}
    \mM^\omega\equiv \{a\in \mM:\Delta_{\omega;\mM}^{it}a\Delta_{\omega;\mM}^{-it}=a\qquad \forall t\}=\mathbb{C} 1
\end{align}
For a nontrivial factor, this requires type $III_1$, but it does not follow from the
type alone, and depends on the state \footnote{For example, the Araki-Woods type III$_1$ algebras are always realized in a state with a non-trivial centralizer. Consider an infinite tensor product of identical three-level
systems with faithful site density matrix
$\rho=\sum_{i=1}^3p_i\ket{i}\bra{i}$.
If $\log(p_1/p_3)$ and $\log(p_2/p_3)$ are rationally independent, the resulting
Araki--Woods factor is type $III_1$~\cite{ArakiWoods1968}.
Nevertheless, the site spectral projections $\ket{i}\bra{i}_n$ belong to the
centralizer of the product state.}. In the ergodic hierarchy of modular flows, the intuition of information loss can be formalized as the condition of strong mixing (clustering in modular time), which is strictly stronger than ergodicity~\cite{furuya2023information}. Type III$_1$ algebras realized in the local algebra of QFT also satisfy this condition.

\subsection{\NL{Modular comparison and contraction estimates}}
\label{app:modular-comparison}

\NL{We give the derivation and proof of Lemma~\ref{lem:modular-monotonicity}.}

The modular operators of $\mM$ and of $\mN$ act on different spaces, $\mathcal H$ and $\mH_{\mN}=\overline{\mN\Omega}$, but they are tied together by the Tomita operators, which require no cyclicity of $\Omega$ for $\mN$: since $\Omega$ is separating for $\mM$ it is separating for $\mN$, and it is cyclic for $\mN$ on $\mH_{\mN}$. Let $S_{\omega;\mM}=J_{\omega;\mM}\Delta_{\omega;\mM}^{1/2}$ on $\mathcal H$ and $S_{\omega;\mN}=J_{\omega;\mN}\Delta_{\omega;\mN}^{1/2}$ on $\mH_{\mN}$ be the closures of $a\Omega\mapsto a^*\Omega$, $a\in\mM$, and of $n\Omega\mapsto n^*\Omega$, $n\in\mN$. For $n\in\mN$,
\begin{equation}
  S_{\omega;\mM}\,Vn\Omega=n^*\Omega=V\,S_{\omega;\mN}\,n\Omega ,
  \label{eq:Tomita-extension}
\end{equation}
so the Tomita operator of $\mM$ extends that of $\mN$ along $V$: $S_{\omega;\mM}V\supseteq VS_{\omega;\mN}$, because $S_{\omega;\mM}$ is closed. Inserting the polar decompositions and using that the modular conjugations are antiunitary involutions,
\begin{equation}
  \Delta_{\omega;\mM}^{1/2}\,V\xi=J_{\omega;\mM}VJ_{\omega;\mN}\,\Delta_{\omega;\mN}^{1/2}\xi ,
  \qquad \xi\in\operatorname{Dom}\Delta_{\omega;\mN}^{1/2},
  \label{eq:modular-extension}
\end{equation}
where $J_{\omega;\mM}VJ_{\omega;\mN}$ is a linear isometry of $\mH_{\mN}$ into $\mathcal H$. \NL{This identity gives the inclusion-contraction estimates of Lemma~\ref{lem:modular-monotonicity}.}

\begin{proof}
(i) is \eqref{eq:modular-extension}, since $J_{\omega;\mM}VJ_{\omega;\mN}$ preserves norms; the form statement is the definition of the order of closed forms, with equality on the smaller domain. (ii) is the Heinz--Kato inequality~\cite{Kato1952}: if $A$ on $\mathcal H$ and $B$ on $\mH_{\mN}$ are positive self-adjoint and $T:\mH_{\mN}\to\mathcal H$ is a contraction with $T\operatorname{Dom}B\subseteq\operatorname{Dom}A$ and $\|AT\xi\|\le\|B\xi\|$ on $\operatorname{Dom}B$, then $T\operatorname{Dom}B^{s}\subseteq\operatorname{Dom}A^{s}$ and $\|A^{s}T\xi\|\le\|B^{s}\xi\|$ for $0\le s\le1$. Apply it with $A=\Delta_{\mM}^{1/2}$, $B=\Delta_{\mN}^{1/2}$, $T=V$ and $s=2\alpha$, using (i). For $V=1$ this is the L\"owner--Heinz theorem, the operator monotonicity of $z\mapsto z^{s}$ for $0\le s\le1$. The endpoints are $D_0=V$ and \eqref{eq:modular-extension}.
\end{proof}

\subsection{\NL{Conditional expectations and product structure}}
\label{app:expectation-proofs}

\NL{\paragraph{Proof of Theorem~
ef{thm:split-no-expectation}.}}
\begin{proof}
Let $\mN\subset\mathcal F\subset\mM$ be an interpolating type I
factor and suppose that $\mE:\mM\to \mN$ is normal.  Its restriction
\begin{equation}
 \mE_{\mathcal F}:=\mE|_{\mathcal F}:\mathcal F\longrightarrow\mN
\end{equation}
is a normal norm-one projection with range $\mN$.  Tomiyama's type
theorem states that a normally expected von Neumann subalgebra of a type I von
Neumann algebra is type I \cite{Tomiyama1959}\cite{Takesaki1972}. This contradicts $\mN$ being a type III factor.
\end{proof}

\NL{\paragraph{Proof of Lemma~
ef{lem:expectation-slice}.}}
\begin{proof}
For $b\in\mN$ and $c\in\mR$,
\begin{equation}
 b \mE(c)=\mE(b c)=\mE(c b)=\mE(c)b.
\end{equation}
Thus $\mE(c)\in Z(\mN)=\mathbb C\mathbbm 1$, so
$\mE(c)=\chi_{\mR}(c)\mathbbm 1$ for a normal state $\chi_{\mR}$.  If $c\geq0$ and
$\chi_{\mR}(c)=0$, faithfulness of $\mE$ implies $c=0$; hence $\chi_{\mR}$ is faithful.  The
second identity is the bimodule property.
\end{proof}

For a faithful normal state $\varphi$ on $\mN$, the state
$\varphi_0:=\varphi\circ \mE$ is faithful and normal and satisfies
\begin{equation}
 \varphi_0(b c)=\varphi_0(b)\chi_{\mR}(c).
\end{equation}

\NL{\paragraph{Proof of Proposition~
ef{prop:conormal-factorization}.}}
\begin{proof}
Multiplication is a $*$-homomorphism on $\mN\odot\mR$ because
the two algebras commute.  Moreover,
\begin{equation}
 \varphi_0\bigl(m(b\otimes c)\bigr)
 =(\varphi\otimes\chi_{\mR})(b\otimes c).
\end{equation}
The GNS representation of the faithful normal product state
$\varphi\otimes\chi_{\mR}$ is the spatial tensor-product representation.  Hence
multiplication extends normally to $\mN\overline\otimes\mR$.
Faithfulness of the product state gives injectivity, while surjectivity follows
from the definition of $\mN\vee\mR$.  Conormality then yields
\eqref{eq:conormal-factorization}; compare
\cite{Faulkner2020,FaulknerSperanza2024}.
\end{proof}

\subsection{\NL{Compression identities for subsystem flow}}
\label{app:compressed-flow}

\NL{With $C(t)$ and $e_{\mN}$ as in Definition~\ref{def:subsystem-flow}, the following identity records the effect of compression.}

To display the compression algebraically, extend the intrinsic unitary to
$U_{\mN}(t)=V\Delta_{\mN}^{it}V^*+(1-e_{\mN})$ on $\mathcal H$ and set
$u_t=\Delta_{\mM}^{-it}U_{\mN}(t)$, $\widehat C(t)=e_{\mN}u_te_{\mN}=VC(t)V^*$.
Writing $\widetilde\sigma^{\omega;\mN}_{-t}=\operatorname{Ad}U_{\mN}(-t)$ gives
\begin{align}
 \widehat C(t+s)&=\widehat C(t)\widetilde\sigma^{\omega;\mN}_{-t}(\widehat C(s))
                  +\mathcal L(t,s),\nn\\
 \mathcal L(t,s)&=e_{\mN}u_t(1-e_{\mN})\widetilde\sigma^{\omega;\mN}_{-t}(u_s)e_{\mN}.
\end{align}
The second term records passage through the orthogonal complement. This algebraic
identity does not itself specify an open-system master equation.
On vectors $\psi\in\operatorname{Dom}K_{\mN}$ with $V\psi\in\operatorname{Dom}K_{\mM}$,
\begin{equation}
 -i\left.\partial_t C_{\mN\subset\mM}(t)\psi\right|_{t=0}
 =(V^*K_{\mM}V-K_{\mN})\psi,\qquad K_{\mathcal C}=-\log\Delta_{\mathcal C}.
\end{equation}
This compares modular generators on the stated domain, rather than asserting
a strong derivative on the full Hilbert space or a dissipative generator.
The contraction is a Hilbert-space operator. Conjugation
$a\mapsto C(t)^*aC(t)$ is CP on $B(\mH_{\mN})$ but is unital only if
$C(t)^*C(t)=1$; preservation of $\mN$ would require a further condition.
An interpretation as a quantum channel therefore needs additional structure.

\subsection{\NL{Characterization by the subsystem flow}}
\label{app:characteristic-function}

\NL{For a common cyclic and separating vector, Borchers\textquotesingle{} characterization uses the following properties~\cite{revolutionizing}.}

\begin{enumerate}
    \item $D(0)=1$, and $D_{\mN\subset\mM}(t)$ and $D_{\mN\subset\mM}(t +\frac{i}{2})$ are unitary and strongly continuous in $t$;
    \item $D_{\mN\subset\mM}(z)$ is analytic on $0<\operatorname{Im}z<1/2$ and contractive on its closure;
    \item $D_{\mN\subset\mM}(t)\ket{\Omega} = \ket{\Omega}$;
    \item The cocycle condition holds: $D(t+s)=\Delta_{\mM}^{-it}D(s)\Delta_{\mM}^{it}D(t)$. Here conjugation by $\Delta_{\mM}^{it}$ acts on $B(\mH)$, since $D(s)$ need not lie in $\mM$;
    \item The intrinsic boundary symmetry is
    \[
     D(t+i/2)^*J_{\mM} D(t)=D(t)^*J_{\mM} D(t+i/2),
    \]
    independent of $t$. For the given inclusion its value is $J_{\mN}$, equivalently
    $D(t+i/2)=J_{\mM} D(t)J_{\mN}$. In particular, $D_{\mN\subset\mM}(\frac{i}{2}) = J_{\mM}J_{\mN}$. This unitary operator is crucial in subfactor theory creating a chain of algebras:
    \begin{equation}
        \mM\supset\mN\supset D_{\mN\subset\mM}(\frac{i}{2})^*\mM D_{\mN\subset\mM}(\frac{i}{2})\supset D_{\mN\subset\mM}(\frac{i}{2})^*\mN D_{\mN\subset\mM}(\frac{i}{2})\supset...
    \end{equation}
    \item For all $t\in\mathbb{R}$, $D_{\mN\subset\mM}(t)D_{\mN\subset\mM}(\frac{i}{2})^*\mM D_{\mN\subset\mM}(\frac{i}{2})D_{\mN\subset\mM}(t)^*\subset\mM$ holds.
\end{enumerate}

\begin{theorem}
    Given a von Neumann algebra $\mM$ and a cyclic and separating vector $\ket{\Omega}$, each subalgebra $\mN\subset\mM$ for which $\ket{\Omega}$ is cyclic is uniquely determined by its subsystem flow $D_{\mN\subset\mM}$, and every operator-valued function with the properties listed above is the subsystem flow of such a subalgebra.
\end{theorem}

\section{Additional quantum error-correction examples}\label{app:extra-qec-examples}

\subsection{\NL{Knill--Laflamme and the complementary channel}}
\label{app:finite-KL}

\paragraph{The finite-dimensional dictionary.}
Let $\Phi(\rho)=\sum_kE_k\rho E_k^\dagger$ be a channel on $B(\mathcal H_S)$ with Stinespring isometry $W\xi=\sum_kE_k\xi\otimes\ket{k}_E$, so that $\Phi=\tr_E\circ\Ad W$. The complementary channel is $\Phi^c=\tr_S\circ\Ad W$, explicitly
\begin{equation}
 \Phi^c(\rho)=\sum_{k,l}\tr\!\bigl(E_l^\dagger E_k\,\rho\bigr)\,\ket{k}\!\bra{l}.
 \label{eq:complementary-channel}
\end{equation}
A code subspace $P\mathcal H_S$ is exactly correctable for $\Phi$ if and only if~\cite{KnillLaflamme1997}
\begin{equation}
 P\,E_l^\dagger E_k\,P=\lambda_{lk}\,P\qquad\hbox{for all }k,l,
 \label{eq:KL}
\end{equation}
for some matrix $\lambda$. For a state $\rho$ supported on $P$ the matrix elements of \eqref{eq:complementary-channel} are $\Phi^c(\rho)_{kl}=\tr(PE_l^\dagger E_kP\,\rho)$. Hence \eqref{eq:KL} holds if and only if
\begin{equation}
 \Phi^c(\rho)=\Lambda\quad\hbox{for every state $\rho$ supported on }P,\qquad\Lambda_{kl}=\lambda_{lk},
 \label{eq:forgetful}
\end{equation}
because the code states span all operators on $P\mathcal H_S$. Condition \eqref{eq:forgetful} says that the complementary channel is constant, or forgetful, on the code: the environment ends in the same state whatever logical state was encoded. Conditions \eqref{eq:KL} and \eqref{eq:forgetful} are therefore one statement read from two sides. From the system's side, no pair of errors can be told apart on the code except through the code-independent Gram matrix $\lambda$; from the environment's side, the environment, whose state records which error occurred, carries no information about the logical input. The content of the Knill--Laflamme theorem is the further implication that \eqref{eq:forgetful} suffices for recovery by a channel acting on the system alone. Its approximate version is the information--disturbance trade-off~\cite{SchumacherWestmoreland2002,KretschmannSchlingemannWerner2008}, and for a logical \emph{algebra} rather than a subspace the scalar $\lambda_{lk}$ is replaced by an element of the commutant of the logical algebra on the code~\cite{BenyKempfKribs2007}.

\paragraph{Stabilizer codes: symmetries, syndromes and modular autonomy.}
Stabilizer codes give an explicit finite-dimensional example of
Eq.~\eqref{eq:symmetry-Haar-expectation}~\cite{Gottesman1997,KnillLaflamme1997}.
Let $\mathsf S$ be an Abelian Pauli stabilizer group not containing $-1$.
Its adjoint action on $\mM=B((\mathbb C^2)^{\otimes n})$ has fixed-point
algebra
\begin{align}
 \mN_{\mathsf S}&=\mathsf S'\cap\mM,
 &\mE_{\mathsf S}(x)&=\frac{1}{|\mathsf S|}
                    \sum_{s\in\mathsf S}sxs^\dagger .
 \label{eq:stabilizer-conditional-expectation}
\end{align}
The simultaneous eigenspace projectors $P_\chi$, indexed by characters
$\chi\in\widehat{\mathsf S}$ that occur in the representation, decompose
this algebra as
\begin{align}
 \cH=\bigoplus_{\chi}\cH_\chi,
 \qquad
 \mN_{\mathsf S}=\bigoplus_{\chi}B(\cH_\chi),
 \qquad
 \mE_{\mathsf S}(x)=\sum_\chi P_\chi xP_\chi .
 \label{eq:stabilizer-syndrome-decomposition}
\end{align}
Thus Haar averaging removes coherences \emph{between} syndromes, but
retains arbitrary observables \emph{within} each syndrome sector. The
ordinary stabilizer code is only the trivial-character sector
$P=P_{\mathbf1}$, with $sP=P$ for all $s\in\mathsf S$. Its projector $P$
should not be confused with the Jones projection
$e_{\mN}=[\mN_{\mathsf S}\Omega]$ in the GNS representation of a faithful
reference state: $P$ acts on the physical Hilbert space, whereas $e_{\mN}$
projects onto the invariant observable-generated GNS subspace.

For example, the three-qubit repetition code is stabilized by
$\mathsf S=\langle Z_1Z_2,Z_2Z_3\rangle$, and
\begin{align}
 P&=\tfrac14(1+Z_1Z_2)(1+Z_2Z_3),
 &P\cH&=\operatorname{span}\{|000\rangle,|111\rangle\}.
 \label{eq:three-qubit-code}
\end{align}
There are four syndrome spaces, each of dimension two, so that
$\mN_{\mathsf S}\simeq\bigoplus_{\chi\in\mathbb Z_2^2}M_2(\mathbb C)$.
The set $\{1,X_1,X_2,X_3\}$ obeys the standard Knill--Laflamme
conditions for the code $P\cH$,
\begin{align}
 P X_i^\dagger X_jP=\delta_{ij}P,
 \qquad X_0:=1,
 \quad i,j\in\{0,1,2,3\} .
 \label{eq:repetition-KL}
\end{align}
The four error spaces are distinguished by their stabilizer syndromes,
which can be measured and used to choose a recovery. This is \emph{active}
error correction; the Haar map in
Eq.~\eqref{eq:stabilizer-conditional-expectation} is the associated
syndrome dephasing, not by itself the full recovery channel.

Choose now a faithful density matrix block diagonal in syndromes,
$\rho=\bigoplus_\chi p_\chi\rho_\chi$, with
$p_\chi>0$, $\rho_\chi>0$ and $\sum_\chi p_\chi=1$. Its modular
flow obeys
\begin{align}
 \sigma_t^\rho(\mN_{\mathsf S})=\mN_{\mathsf S},
 \qquad
 \sigma_t^\rho(x_\chi)=\rho_\chi^{it}x_\chi\rho_\chi^{-it},
 \quad x_\chi\in B(\cH_\chi).
 \label{eq:stabilizer-modular-flow}
\end{align}
Hence the syndrome algebra forms an autonomous modular subsystem, and
the reference-preserving code $\mathcal C_{\mE_{\mathsf S}}$ consists of
all syndrome-block-diagonal states. Neither the modular-invariance
condition nor this family of states is identical to the subspace
correctability condition~\eqref{eq:repetition-KL}. In particular, the
maximally mixed code state $P/\Tr P$ is not faithful on $\mM$; a
faithful reference is obtained by assigning positive weight to all
syndromes. Taking different full-rank $\rho_\chi$ makes the modular
flow nontrivial even though it preserves the syndrome algebra.

\paragraph{Topological stabilizers and noncontractible logical operators.}
The toric code gives a geometrical version of the same distinction
between a stabilizer-fixed algebra and a correctable subspace
\cite{Kitaev2003}. On a lattice with periodic boundary conditions,
take commuting star and plaquette stabilizers
$A_v=\prod_{e\ni v}X_e$ and $B_p=\prod_{e\in\partial p}Z_e$.
Their common $+1$ eigenspace has projector
$P=\prod_v(1+A_v)/2\prod_p(1+B_p)/2$, with redundant global
constraints understood. Every operator $o$ supported on a sufficiently
small contractible disk obeys $PoP=c_oP$, expressing local
indistinguishability and correction of erasure of that disk. In contrast,
noncontractible Wilson-loop operators commute with the stabilizers but
act nontrivially on the logical ground-state space. The stabilizer Haar
map on the full matrix algebra still projects onto all syndrome blocks,
not just onto the logical operators of $P\cH$. A pure toric-code ground
state is not a faithful reference state on that full algebra; modular
autonomy of a faithful stabilizer-invariant reference is therefore a
different statement from local correctability of the ground-state code.

\paragraph{Noiseless and operator-algebra quantum subsystems.}
A second example makes the relation to the earlier tensor-product
discussion direct~\cite{KribsLaflammePoulin2005,BenyKempfKribs2007}.
Let $\cH=\cH_L\otimes\cH_G$, where $L$ is a logical system and $G$ a
gauge subsystem, and set
\begin{align}
 \mN=B(\cH_L)\otimes1_G\subset
 \mM=B(\cH_L\otimes\cH_G),\qquad
 \mE_\tau(x)=(\operatorname{id}_L\otimes\tau)(x)\otimes1_G .
 \label{eq:subsystem-gauge-expectation}
\end{align}
Here $\tau$ is a faithful state of the gauge factor. A faithful
reference state $\rho_{LG}$ has autonomous modular dynamics on $\mN$
precisely when
\begin{align}
 \rho_{LG}=\rho_L\otimes\tau_G,
 \qquad
 \sigma_t^\rho(b\otimes1_G)
   =\rho_L^{it}b\rho_L^{-it}\otimes1_G .
 \label{eq:subsystem-gauge-product}
\end{align}
Then $\mE_\tau$ is reference preserving, and erasure of the gauge
system is exactly correctable for the family $\{\varphi_L\otimes\tau_G\}$.
The gauge factor may evolve nontrivially, but its state remains fixed as
the logical input changes. For a conventional subsystem code supported
on a subspace $P\cH\simeq\cH_L\otimes\cH_G$, error correction against
a noise channel with Kraus operators $E_a$ instead requires
\begin{align}
 P E_a^\dagger E_bP=1_L\otimes\Lambda_{ab}
 \qquad\text{on }P\cH,
 \label{eq:subsystem-KL}
\end{align}
which permits recoverable errors acting on the gauge factor. The
reference-preserving expectation describes a particular family of
encoded states and a particular erasure channel; it must not be taken as
a criterion for correcting every noise model.

\paragraph{A non-Abelian noiseless subsystem.}
Collective $SU(2)$ noise on three spin-$\frac12$ particles provides an
example in which the multiplicity space of a symmetry representation
encodes a logical qubit. Angular-momentum addition gives
\begin{align}
 (\mathbb C^2)^{\otimes3}
 &\simeq V_{3/2}\oplus\bigl(V_{1/2}\otimes\mathbb C^2\bigr),\nonumber\\
 B((\mathbb C^2)^{\otimes3})^{SU(2)}
 &\simeq\mathbb C1_{V_{3/2}}\oplus
       \bigl(1_{V_{1/2}}\otimes M_2(\mathbb C)\bigr).
 \label{eq:collective-noiseless-subsystem}
\end{align}
Collective rotations act on $V_{1/2}$ but trivially on the second
$\mathbb C^2$ factor: the latter is a noiseless logical subsystem.
For any faithful rotationally invariant reference density matrix, the
fixed-point algebra in Eq.~\eqref{eq:collective-noiseless-subsystem}
is modular covariant. Its intrinsic modular dynamics can act
nontrivially on the multiplicity qubit, even though collective noise
cannot distinguish its logical states. This illustrates how protected
information may arise from representation multiplicities rather than
from pointwise fixed physical vectors.

\section{Classical subsystem flows and ergodic theory}\label{sec:classical}

\begin{table}[htbp]
\centering
\label{tab:classical-quantum-dictionary}
\begingroup
\footnotesize
\setlength{\tabcolsep}{4pt}
\renewcommand{\arraystretch}{1.22}
\begin{tabular}{@{}>{\raggedright\arraybackslash}p{0.19\linewidth}
                >{\raggedright\arraybackslash}p{0.365\linewidth}
                >{\raggedright\arraybackslash}p{0.365\linewidth}@{}}
\toprule
\textbf{Concept} & \textbf{Classical} & \textbf{Quantum} \\
\midrule
Observable algebra &
$\mM=L^\infty(X,\Sigma,\mu)$, commutative. &
$\mM\subset B(\mH_{\mM})$, generally noncommutative \\
\addlinespace[2pt]
Coarse-graining &
Sub-$\sigma$-algebra $\Sigma_0\subset\Sigma$; for a subsystem, a projection $\pi:X\to Y$. &
Inclusion $\mN\subset\mM$; no tensor factorization is required. \\
\addlinespace[2pt]
Reduced state &
$\rho_0=\mE_\mu(\rho)$ (in a product, $\rho_0(y,z)=\rho_Y(y)$). &
Restriction $\omega|_{\mN}$, with intrinsic modular operator $\Delta_{\mN}$. \\
\addlinespace[2pt]
State-generated flow &
Poisson flow of $K_\rho=-\log\rho$:
$\delta_\rho(f)=\{K_\rho,f\}$. &
Modular flow $\sigma_t^{\omega;\mM}=\operatorname{Ad}\Delta_{\mM}^{it}$,
$K_{\mM}=-\log\Delta_{\mM}$. \\
\addlinespace[2pt]
Subsystem flow &
$\mathcal V_t=\alpha_t^\rho\circ\alpha_{-t}^{\rho_0}$,
comparing full and lifted reduced motions. &
$D_{\mN\subset\mM}(t)=\Delta_{\mM}^{-it}V\Delta_{\mN}^{it}$,
$V:\mH_{\mN}\hookrightarrow\mH_{\mM}$. \\
\addlinespace[2pt]
Infinitesimal comparison &
$K_{\mathrm{rel}}=K_\rho-K_0
=-\log(\rho/\rho_0)$ generates the relative Poisson motion. &
$V^\dagger K_{\mM}V-K_{\mN}$ is the comparison quadratic form on $\mH_{\mN}$. \\
\addlinespace[2pt]
Information loss &
$S(\rho\|\rho_0)=-\langle K_{\mathrm{rel}}\rangle_\rho\ge0$. &
$\langle V^\dagger K_{\mM}V-K_{\mN}\rangle_\Psi
=g_{\rm ret}+g_{\rm dual}\ge0$. \\
\addlinespace[2pt]
Conditional expectation &
A state-preserving expectation onto every measurable coarse algebra exists. &
An $\omega$-preserving normal expectation exists iff
$\sigma_t^{\omega;\mM}(\mN)=\mN$ (Takesaki). \\
\addlinespace[2pt]
Autonomous subsystem &
$\alpha_t(\mN)=\mN$;
for the product Poisson model, $\rho=\rho_Y\rho_Z$ ensures agreement of full and reduced flows. &
$\sigma_t^{\omega;\mM}(\mN)=\mN$ iff
$D_{\mN\subset\mM}(t)=V$; equivalently, reference-preserving QEC. \\
\addlinespace[2pt]
Ergodicity and mixing &
$\mM^\alpha=\mathbb C1$;
clustering of correlations is stronger. &
$\mM^\omega=\mathbb C1$ (trivial centralizer);
modular mixing implies this and requires type $\mathrm{III}_1$. \\
\addlinespace[2pt]
History and invariant core &
$\mN_{\rm hist}=\bigvee_t\alpha_t(\mN)$,
$\mN_{\rm core}=\bigcap_t\alpha_t(\mN)$. &
$\mN_{\rm hist}=\bigvee_t\sigma_t^{\omega;\mM}(\mN)$,
$\mN_{\rm core}=\bigcap_t\sigma_t^{\omega;\mM}(\mN)$. \\
\addlinespace[2pt]
One-sided nesting &
$\alpha_t(\mN_+)\subseteq\mN_+$ for $t\ge0$;
full history and trivial tail give the $K$-property. &
$\sigma_t^{\omega;\mM}(\mN)\subseteq\mN$ for $t\le0$;
for an HSMI, $K_{\mM}-K_{\mN}=2\pi P\ge0$. \\
\bottomrule
\end{tabular}
\endgroup
\vspace{0.3em}

\caption{Dictionary between the classical constructions of this appendix and the
quantum subsystem flows studied in the main text. The classical state-generated
Poisson flow is an analogy, not the Tomita--Takesaki modular flow of an
abelian von Neumann algebra. Here $g_{\rm ret}=S_{\mM}(\varphi\|\omega)-S_{\mN}(\varphi\|\omega)$ and
$g_{\rm dual}=S_{\mN'}(\varphi\|\omega)-S_{\mM'}(\varphi\|\omega)$, with
$S_{\mN'}$ evaluated on the support of the reference state when $e_{\mN}<1$.
The expectation $\mE_\mu$ defining $\rho_0$ preserves $\mu$; a
$\rho$-preserving expectation also exists, but need not coincide with it.
The classical $\alpha_t$ in the ergodic hierarchy is a specified
measure-preserving flow; the Poisson state flow $\alpha_t^\rho$ conserves
$\rho$ and need not be ergodic.}
\end{table}

In this appendix, we review notions from classical ergodic theory to motivate and build intuition about subsystem flows.
We begin by reviewing the measure-theoretic definition of a classical dynamical system in Appendix~\ref{subsec:classical-definitions}, the subspaces and subsystems of observables in Appendix~\ref{subsec:classical-region-subsystem}, and their role in classical ergodic theory in Appendix~\ref{subsec:classical-ergodic-hierarchy}. In Appendix~\ref{subsec:classical-state-flow}, we argue that when the phase space $X$ is equipped with a Poisson bracket (as suggested by the phase space intuition from physics) {\it the Poisson flow of a state} is a Poisson-geometric avatar of state-generated modular dynamics. It is distinct from the trivial Tomita--Takesaki modular automorphism of an abelian algebra, but it generates an automorphism leaving the state invariant and satisfies a classical avatar of the KMS condition \footnote{In the quantum case, the KMS condition at inverse temperature one characterizes the modular automorphism group; stationarity alone does not~\cite{Takesaki1970}.}. However, crucially for us, as opposed to the quantum relative modular data and the associated inclusion contractions, the Poisson Hamiltonian of a state does not satisfy a monotonicity condition under coarse-graining as in the quantum case \footnote{Recall that in the quantum case, this monotonicity originates from the fact that the Tomita operator of the algebra is an extension of its restriction to the subalgebra. In the classical case, there is no analog of the Tomita operator.}. This monotonicity is at the heart of non-commutative results on the reconstruction of geometry from a net of subsystems, such as half-sided modular inclusions and intersections as we review in the companion QFT paper~\cite{companionQFT}. This appendix is not needed for the quantum development in the main text.

For each coarse-graining, we define a notion of a coarse-graining flow. For a general coarse-graining, the relative phase $(\rho/\rho_0)^{it}$ of Appendix~\ref{subsec:classical-definitions} is a generating function for the entropy of the coarse-graining. However, the focus of this work is {\it subsystem flows} which are richer.
Classical many-body systems and field theory come equipped with a net of subsystems, and the quantum generalization of subsystem flows plays a crucial role in the emergence of geometry from a net of abstract subsystems.

\subsection{Review of classical dynamical systems}\label{subsec:classical-definitions}

\paragraph{Algebra of observables:} In classical mechanics, one often begins with a phase space that we take to be an orientable smooth manifold \(X\) that admits a strictly positive volume measure. Points \(x\in X\) represent microstates of the system. But from the algebraic viewpoint, the primary object is not the set of points but the commutative von Neumann algebra of observables (complex measurable functions)
\[
\mM=L^\infty(X,\Sigma,\mu),
\]
where \((X,\Sigma,\mu)\) is a standard probability space. When $X$ is a smooth phase space of finite volume, we take $\mu$ to be the normalized volume measure\footnote{$\Sigma$ is the completion of the Borel $\sigma$-algebra with respect to $\mu$. Elements of $L^\infty(X,\Sigma,\mu)$ are equivalence classes of essentially bounded measurable functions that agree $\mu$-almost everywhere.}.

\paragraph{States:} The normalized measure $\mu$ above is interpreted as a distinguished reference state\footnote{A $\sigma$-finite volume measure of infinite total volume instead defines a distinguished weight. Weights are introduced in Section~\ref{sec:modular-vs-dynamics}.}. General normal states correspond to probability measures $\nu$ on $X$ that are absolutely continuous with respect to $\mu$ (denoted $\nu\ll \mu$), meaning that for every measurable set $E\subseteq X$, if $\mu(E)=0$ then $\nu(E)=0$.
There exists a $\mu$-measurable function $\rho_\nu:X\to [0,\infty)$ such that $\nu(E)=\int_E \rho_\nu d\mu$ \footnote{In mathematical terminology, $\rho_\nu$ is called a Radon-Nikodym derivative.}. In other words, every normal state can be viewed as a probability density function $\rho_\nu$ on $X$, with expectation values of an observable $f$ in this state captured by integrals on $X$:
\begin{align}
\braket{f}_\nu=\int_{x\in X}f(x) d\nu(x)=\int_{x\in X} \rho_\nu(x) f(x) d\mu(x)\ .
\end{align}

\paragraph{GNS Hilbert space:} The GNS Hilbert space of the reference state $\mu$ is
\[
  \cH=L^2(X,\Sigma,\mu),
  \qquad
  \mM=L^\infty(X,\Sigma,\mu).
\]
Observables act by multiplication, and the reference state $\mu$ is represented by the constant unit vector $\ket{1}$. A general normal state with density $\rho$ is represented by
\[
  \ket{\sqrt{\rho}}\in\cH,
  \qquad
  \|\ket{\sqrt{\rho}}\|_2^2=\int_X\rho\,\dd\mu=1.
\]
The GNS construction builds a connection between observables and states.

\paragraph{Dynamics:} Dynamics will be described by a measurable one-parameter group of invertible, measure-preserving transformations $\Phi_t:X\to X$. General dynamics need not be smooth. In physics, we are often interested in smooth flows described by a one-parameter family of volume-preserving diffeomorphisms of $X$. Taking a $t$-derivative of how a point $x\in X$ moves under this flow defines a vector field $V$, which preserves the volume measure if and only if $\mathcal{L}_V\mu=0$; i.e., it is divergence-free with respect to $\mu$.

\paragraph{Coarse-grainings:} Classical coarse-graining retains a subalgebra of observables and discards the information needed to distinguish states that agree on this subalgebra. Two useful constructions arise from a subregion of the space of microstates and from a subsystem of observed variables. Both can be formulated in measure theory, before introducing a Hamiltonian or a Poisson bracket. We first describe these constructions and then ask when the retained observables evolve autonomously.

In the measure-theoretic setting above, a coarse description is specified by a sub-$\sigma$-algebra $\Sigma_0\subseteq\Sigma$, with associated unital von Neumann subalgebra
\begin{align}
    \mN=L^\infty(X,\Sigma_0,\mu)\subseteq\mM\ .
\end{align}
The reference measure $\mu$ determines the coarse-graining map
\begin{align}
    \mE_{\mN}(f)=\mathbb \mE_\mu[f\mid\Sigma_0]\ .
\end{align}
This is the unique $\mu$-preserving normal conditional expectation onto $\mN$. On $L^2(X,\mu)$ it acts as the orthogonal projection onto $L^2(X,\Sigma_0,\mu)$. A general normal state $\omega_\rho(f)=\int_X\rho f\,d\mu$ is coarse-grained to $\omega_0=\omega_\rho\circ \mE_{\mN}$, whose density is $\rho_0=\mE_{\mN}(\rho)$, using the $L^1$ extension of conditional expectation. The two states agree on $\mN$. Thus the subalgebra specifies the retained information, while the chosen expectation specifies how a state on this subalgebra is extended back to the full algebra.

\paragraph{Entropy of a coarse-graining:} Consider a $\mu$-preserving coarse-graining $\mE$, a normal state with density $\rho$, and its coarse-grained density $\rho_0=\mE(\rho)$. Then, the relative entropy (KL-divergence) of these two states associates an entropy to this coarse-graining:
\begin{align}
    \label{eq:KL-classical}
    S(\rho\|\rho_0)=\int_X \rho \log(\rho/\rho_0)d\mu=H(\rho_0)-H(\rho)
\end{align}
where $H(\rho)=-\int_X\rho\log\rho\,d\mu$ is the entropy relative to the chosen reference measure, and the equality assumes that the two entropies are finite. The nonnegativity of relative entropy implies that coarse-graining cannot decrease this entropy. This compares a state with its coarse-grained version; it does not by itself imply monotonic entropy growth under time evolution~\cite{Lashkari2019EntanglementScale}.

\paragraph{Generating function for entropy:} For faithful densities $\rho,\rho_0>0$ almost everywhere, the multiplicative relative phase
\[
u_t
=
\left(\frac{\rho}{\rho_0}\right)^{it}
\]
is the generating function for coarse-graining entropy because
\begin{align}
&-i\left.\frac{\dd}{\dd t}u_t\right|_{t=0}
=
\log\frac{\rho}{\rho_0},\nn\\
&S(\rho\|\rho_0)=-i \int_X \rho (\partial_t u_t\big|_{t=0})d\mu\ .
\end{align}
Here the differentiation under the integral is justified when $\rho\log(\rho/\rho_0)$ is integrable. More generally, for every faithful density $\rho$ we can define the strongly continuous unitary multiplication group
\[
  U_t^\rho:=\Mop_{\rho^{it}}
  =\exp\!\bigl(it\Mop_{\log\rho}\bigr),
  \qquad
  (U_t^\rho\psi)(x)=e^{it\log\rho(x)}\psi(x).
\]
Similarly, the coarse-grained density defines
\[
  U_t^{\rho_0}:=\Mop_{\rho_0^{it}}.
\]
The parameter $t$ is an auxiliary phase parameter. Nothing in this construction identifies it with laboratory time. In particular, $U_t^\rho$ does not transport points of $X$; it only rotates the phase of a Hilbert space vector pointwise\footnote{For faithful normal states on the same abelian algebra, $u_t=(\rho/\rho_0)^{it}$ is the Connes Radon-Nikodym cocycle. Relative modular operators remain nontrivial, even though the modular automorphism group of each state is trivial.}.

\subsection{Subspace (single-particle) and subsystem (multi-particle) coarse-grainings}
\label{subsec:classical-region-subsystem}

Subspace and subsystem coarse-grainings organize incomplete information differently. An example of a subspace coarse-graining resolves microstates within a selected region of the phase space and identifies those outside it. A microstate is either inside the coarse-graining region or outside it. The set of microstates adds up in this case; that is why, in anticipation of the generalization to Hilbert space, we refer to them as {\it subspaces} so that addition of the number of independent microstates becomes a direct sum. As we will see, there is a net of subspace algebras on the phase space of a single particle. As an example of such a net of coarse-grainings, consider the reduced state to the part of the phase space of a pendulum where the angle $\theta$ is in some limited range $\theta\in (-\theta_0,\theta_0)$, and the set of all $\theta_0\in (0,\pi)$.

Next, consider a microstate described using several (potentially independent) degrees of freedom. In this case, the set of microstates multiplies, and in anticipation of generalization to Hilbert spaces, we refer to this case as {\it subsystems}, so that multiplication is the tensor product. As a motivating example, consider two particles with the phase space $(q_1,p_1,q_2,p_2)$, and the restriction to the first particle as a coarse-graining on the second particle's degrees of freedom.

Note that since subspaces could have the same dimension as the entire phase space, they can be discussed even in a single-particle setup. A subsystem is instead described by a coarse space $Y$ and a map $\pi: X\to Y$ where $Y$ need not be an embedded submanifold of $X$. There is a net of subsystems on the phase space of many-body systems and field theory.

\paragraph{Subspace coarse-grainings:}
For a measurable region $O\subset X$, define the algebra generated by observables that are constant outside of $O$:
\begin{align}
    \mM_O=\{f\in\mM:\ f=f \cdot 1_O +c 1_{O^c}\text{ for some }c\in\mathbb C\}
\end{align}
where $1_O$ is the indicator function on $O$, and $(f\cdot 1_O)(x)=f(x)1_O(x)$ is point by point multiplication.
This algebra resolves the microstate completely within $O$ and retains only the event ``outside $O$'' on its complement. Assuming $0<\mu(O)<1$, its conditional expectation is
\begin{align}
    \mE_O(f)=1_O f+
    1_{O^c}\mathbb{E}_{O^c}(f)
\end{align}
where
\begin{align}
    \mathbb{E}_{O^c}(f)=\frac{1}{\mu(O^c)}\int_{O^c} f d\mu
\end{align}
is the average of $f$ on $O^c$.
For $O_1\subseteq O_2$, we have $\mM_{O_1}\subseteq\mM_{O_2}$, so this construction gives a net of coarse observable algebras. Since a microstate $x\in X$ is either in $O$ or $O^c$ we have
\begin{align}
    \mM=1_O\mM\oplus 1_{O^c}\mM
    \simeq L^\infty(O,\mu|_O)\oplus L^\infty(O^c,\mu|_{O^c})\ .
\end{align}
Our unital subspace algebras have the form $\mM_O=1_O\mM\oplus\mathbb C1_{O^c}$, so $\mM_O$ and $\mM_{O^c}$ are not themselves the two direct summands. In fact,
\begin{align}
    \mM_O\cap\mM_{O^c}=\operatorname{span}\{1_O,1_{O^c}\}\ .
\end{align}

\paragraph{Subsystem coarse-grainings:}
Suppose instead that the probability space has a product decomposition
\begin{align}
    (X,\Sigma,\mu)=(Y\times Z,\Sigma_Y\otimes\Sigma_Z,\mu_Y\otimes\mu_Z)\ .
\end{align}
A microstate now specifies both an observed variable $y$ and a hidden variable $z$. The observable algebra factorizes as
\begin{align}
    \mM\simeq\mM_Y\,\bar\otimes\,\mM_Z,
    \qquad \mM_Y=L^\infty(Y,\mu_Y),\quad
    \mM_Z=L^\infty(Z,\mu_Z)\ .
\end{align}
Retaining the subsystem $Y$ means retaining $\mN_Y=\mM_Y\otimes 1$. The corresponding conditional expectation averages over the hidden variables:
\begin{align}
    (\mE_Y f)(y,z)=\int_Z f(y,z')\,d\mu_Z(z')\ .
\end{align}
This operation forgets the hidden variables at every microstate. The tensor factors describe simultaneously specified degrees of freedom. Although the reference measure is a product, a general state with density $\rho(y,z)$ can have correlations between them.

More generally, a measurable map $\pi:X\to Y$ between standard probability spaces specifies coarse observables through
\begin{align}
    &\mN_Y=\pi^*\mM_Y\subset \mM,
    \qquad \mu_Y=\mu\circ \pi^{-1}\nn\\
    &\mM_Y=L^\infty(Y,\mu_Y)
\end{align}

Here $\mN_Y=\pi^*\mM_Y$ is the pullback of $\mM_Y$, and $\mu_Y=\pi_*\mu$ is the
push-forward measure. Under the standard Borel assumptions, disintegration gives
conditional probability measures $\mu_y$ concentrated on the fibers $\pi^{-1}(y)$.

The conditional expectation is
\begin{align}
    (\mE_\pi f)(x)=\int_{\pi^{-1}(\pi(x))}f(x')\,d\mu_{\pi(x)}(x')\ .
\end{align}
The fiber measures may depend on $y$, and no global product decomposition is required.

Note that both subspace and subsystem conditional expectations are special cases of a common measure-theoretic notion of coarse-graining.

\paragraph{Dynamically closed and invariant coarse-grainings:}
Let $\Phi_t:X\to X$ be a measurable one-parameter group of invertible, measure-preserving transformations. It evolves observables according to
\begin{align}
    \alpha_t(f)=f\circ\Phi_{-t}\ .
\end{align}
For any coarse algebra $\mN\subseteq\mM$, we can follow the family $\alpha_t(\mN)$. We call the coarse description dynamically closed when
\begin{align}\label{eq:classical-coarse-invariance}
    \alpha_t(\mN)=\mN,\qquad t\in\mathbb R\ .
\end{align}
For nonatomic $\mu$ and $0<\mu(O)<1$, invariance of the subspace algebra $\mM_O$ violates ergodicity, as we show below.
For the subsystem $Y$ in a product space $Y\times Z$, dynamical closure means that the $Y$-component of $\Phi_t(y,z)$ depends only on $y$ \footnote{Note that this does not follow from preservation of the product measure.}. For a coarse-graining specified by $\pi$, this means that there is a measure-preserving evolution $\Psi_t$ on the coarse space such that
\begin{align}\label{eq:classical-factor-intertwining}
    \pi\circ\Phi_t=\Psi_t\circ\pi\ ,
\end{align}
understood modulo null sets. Evolving the microstate and then retaining its coarse variables gives the same result as evolving the coarse variables themselves. In ergodic theory such a coarse description is called a {\it factor system}, and we say the coarse-grained variables evolve autonomously. Intuitively, this means that the future of the macroscopic (or observed) variables can be determined from the present macroscopic state, without the need to know the hidden microscopic details. In other words, there is an effective theory that describes the coarse-grained variables.

\paragraph{Effective theory (Compatibility of coarse-graining and evolution):}
For the measure-preserving dynamics above, invariance of $\mN$ is equivalent to
\begin{align}\label{eq:classical-expectation-equivariance}
    \mE_{\pi}\circ\alpha_t=\alpha_t\circ \mE_{\pi},
    \qquad t\in\mathbb R\ .
\end{align}
Indeed, when $\mN_Y$ is invariant, $\alpha_t\circ \mE_{\pi}\circ\alpha_{-t}$ is another $\mu$-preserving conditional expectation onto $\mN_Y$, so uniqueness gives the equality. Conversely, this equality ensures that $\alpha_t$ maps the range of $\mE_{\pi}$ onto itself. The physical interpretation is macroscopic predictability: the future observed variables can be determined from their present values. This compatibility supplies an effective dynamics on the coarse algebra.

\paragraph{Correlation functions of observables and subalgebras:}
For complex functions $f,g\in\mM$, define the correlation function and its connected part as
\begin{align}
    &C_{f,g}(t)=\mu\bigl(f^*\alpha_t(g)\bigr)\nn\\
    &C^{con}_{f,g}(t)
    :=\mu\bigl(f^*\alpha_t(g)\bigr)-\mu(f^*)\mu(g),
    \qquad \mu(f)=\int_X f\,d\mu\ .
\end{align}
Ergodicity is equivalent to the vanishing of the time-averaged connected correlation for every pair of observables:
\begin{align}
    \lim_{T\to\infty}\frac{1}{T}\int_0^T C^{con}_{f,g}(t)\,dt=0\ .
\end{align}
A stronger condition, {\it weak mixing}, requires
\begin{align}
    \lim_{T\to\infty}\frac{1}{T}\int_0^T |C^{con}_{f,g}(t)|^2\,dt=0\ .
\end{align}
The flow is {\it mixing} if $C^{con}_{f,g}(t)\to0$ as $t\to+\infty$. These conditions distinguish cancellation after time averaging from actual loss of correlation at late times. Mixing does not specify a decay rate and does not imply that there are no autonomous subsystems.

Quantitative measures of the correlation between two subspace algebras, or between two subsystems, and their relation to the total variation distance from the product measure, are collected in Appendix~\ref{app:classical-correlations}.

\paragraph{Subsystem coarse-grainings of a free relativistic field theory:} Classical field theory provides a natural setting for a net of subsystem algebras with infinitely many degrees of freedom. For an unconstrained real scalar field on a Cauchy surface $\Sigma$, a phase-space point is a pair of initial data $(\phi,\pi)$. An open spatial region $U\subset\Sigma$ selects the unital polynomial algebra
\begin{align}
    \mN_U
    :=\operatorname{Alg}_{\mathbb R}\{1,\phi(f),\pi(g):f,g\in C_c^\infty(U)\},
    \qquad
    \{\phi(f),\pi(g)\}=\int_\Sigma fg\,d\Sigma\ .
\end{align}
Here $\phi(f)=\int_\Sigma f\phi\,d\Sigma$, and similarly for $\pi(g)$; the other brackets between generators vanish. Thus $\mN_U$ is a Poisson subalgebra, and algebras associated with disjoint spatial regions Poisson commute. Retaining $\mN_U$ resolves the field data in $U$ while discarding the exterior data. Although labeled by a spatial region, this is a subsystem coarse-graining: $U$ selects degrees of freedom, whereas a region of phase space selects entire field configurations.

Choose an ambient Poisson algebra $\mM$ containing these observables and stable under a complete Hamiltonian evolution. With our convention for the pullback, the transported subsystem is
\begin{align}
    \alpha_t^H(F)=F\circ\Phi_{-t}^H,
    \qquad \partial_t\alpha_t^H(F)=\{H,\alpha_t^H(F)\},
    \qquad \alpha_t^H(\mN_U)\subseteq\mM\ .
\end{align}
Hamiltonian evolution preserves Poisson brackets, so every $\alpha_t^H(\mN_U)$ remains a Poisson subalgebra. Invariance for all $t$ gives an autonomous reversible subsystem. More generally, propagation can make the transported observables depend on initially discarded data. Whether their time translates generate $\mM$, or instead form a nested family, is a property of the chosen subsystem and dynamics together.

\subsection{Correlation coefficients for classical subspaces and subsystems}\label{app:classical-correlations}

This appendix records the measures of correlation between two coarse algebras of Appendix~\ref{subsec:classical-region-subsystem}, for a subspace pair and for a subsystem pair, and their relation to the total variation distance from the product measure. They quantify the distinction between the two kinds of coarse-graining but are not used in the main text.

\paragraph{Correlations between subspaces:} Write $C(f,g)=C_{f,g}(0)$ and similarly for its connected part. Of course, the definition of correlation above changes if we multiply the function by a constant. To account for this and define a measure of correlations between two subalgebras, we define
\begin{align}
    C(O_1,O_2)\equiv \sup_{\text{supp}(f_i)\in O_i,|f_i|_\infty\leq 1}|C(f_1,f_2)|
\end{align}
where $|f|_\infty=\sup_{x\in X}|f(x)|$. It follows from the definition above that
\begin{align}
    C(O_1,O_2)\leq \int_{O_1\cap O_2} 1=\mu(O_1\cap O_2)=C(1_{O_1},1_{O_2})
\end{align}
Therefore, we have
\begin{align}
    C(O_1,O_2)=C(1_{O_1},1_{O_2})\ .
\end{align}

For connected correlations, define the event coefficient
\begin{equation}
 \alpha(O_1,O_2)=
 \sup_{\substack{A\subset O_1\\B\subset O_2}}
 |\mu(A\cap B)-\mu(A)\mu(B)|.
\end{equation}
In particular,
\begin{equation}
 C^{con}_{1_{O_1},1_{O_2}}(t)
 =\mu\bigl(O_1\cap\Phi_t(O_2)\bigr)-\mu(O_1)\mu(O_2).
\end{equation}
This measures the correlation of an event with a transported event.

\paragraph{Correlations between subsystems:}
Let $X=Y\times Z$ carry a joint probability measure $\mu_{YZ}$ with marginals
$\mu_Y,\mu_Z$. The unconnected correlation has supremum one over functions bounded
by one, because constant functions attain it. For connected correlations, distinguish
\begin{align}
 \alpha(Y,Z)&=\sup_{\substack{A\subset Y\\B\subset Z}}
 |\mu_{YZ}(A\times B)-\mu_Y(A)\mu_Z(B)|,\nn\\
 \Gamma(Y,Z)&=\sup_{\substack{f,g\ {\rm real}\\\|f\|_\infty,\|g\|_\infty\leq1}}
 |\mu_{YZ}(fg)-\mu_Y(f)\mu_Z(g)|=4\alpha(Y,Z).
\end{align}
The factor four follows by rescaling real functions from $[-1,1]$ to $[0,1]$
and writing the latter as integrals of indicator functions. With
$\nu=\mu_{YZ}-\mu_Y\otimes\mu_Z$, total variation satisfies
\begin{equation}
 \alpha(Y,Z)\leq\operatorname{TV}(\mu_{YZ},\mu_Y\otimes\mu_Z)
 :=\sup_E|\nu(E)|=\tfrac12\|\nu\|_1,
 \qquad \Gamma(Y,Z)\leq2\operatorname{TV}.
\end{equation}
Here $\|\nu\|_1$ is the total variation norm of the signed measure. All these
coefficients vanish precisely when the joint measure is a product.

\subsection{Ergodic hierarchy, subspaces and subsystems}
\label{subsec:classical-ergodic-hierarchy}

The ergodic hierarchy describes how much information about a state remains accessible at late times~\cite{CornfeldFominSinai1982}. In the algebraic language, it can be studied both through correlations of observables and through the evolution of coarse observable algebras. For our purposes, the relevant distinction is between conserved events, autonomous subsystems, and subalgebras that are preserved only in one time direction.

\paragraph{Ergodicity:}
Let $(\mM,\mu,\alpha_t)$ be the measure-preserving dynamical system introduced above. The flow is {\it ergodic} if every measurable set invariant under all $\Phi_t$ has measure zero or one. Equivalently, its fixed-point algebra is trivial:
\begin{align}
    \mM^\alpha
    :=\{f\in\mM:\alpha_t(f)=f\text{ for all }t\in\mathbb R\}
    =\mathbb C1\ .
\end{align}
For nonatomic $\mu$ and $0<\mu(O)<1$, ergodicity forbids invariant subspace algebras. If
\begin{align}
    \alpha_t(\mM_O)=\mM_O,\qquad t\in\mathbb R,
\end{align}
then $\Phi_t(O)=O$ modulo null sets and $1_O\in\mM^\alpha$ (it is a {\it conserved event}), which contradicts ergodicity. However, ergodicity does not exclude an invariant subsystem $\mN\subset\mM$: the condition $\alpha_t(\mN)=\mN$ permits observables in $\mN$ to evolve into one another. Such a subsystem is an autonomous factor of the dynamics (conserved subsystem). Taking two identical independent copies of a classical dynamical system always results in autonomous factors; however, a standard result in ergodic theory implies that if the original system is weakly mixing, the combined dynamical system is ergodic \footnote{For an invertible discrete-time example, take two independent bilateral Bernoulli shifts on $\{0,1\}^{\mathbb Z}$ with their product probability measures. Each shift, and their product, is mixing.}.

\paragraph{Generating property:}
Given a coarse algebra $\mN$, the family $\alpha_t(\mN)$ describes the observables accessible as this coarse description is transported by the dynamics. If $\alpha_t(\mN)=\mN$, evolution remains within the original description. Otherwise, the transported observables can depend on variables that were initially hidden. The algebra
\begin{align}
    \mN_{\rm hist}:=\bigvee_{t\in\mathbb R}\alpha_t(\mN)
\end{align}
contains the information accessible from all time translates of the coarse observables. Here $\bigvee$ denotes the generated von Neumann algebra. The condition $\mN_{\rm hist}=\mM$ says that these translates generate the full observable algebra. This is called the {\it generating property} and is distinct from ergodicity of the dynamics.

\paragraph{Kolmogorov property:}
Of particular interest are subalgebras satisfying
\begin{align}\label{eq:classical-one-sided-inclusion}
    \alpha_t(\mN_+)\subseteq\mN_+,\qquad t\geq0\ .
\end{align}
The family $\alpha_t(\mN_+)$ is then decreasing: $\alpha_{t+s}(\mN_+)\subseteq\alpha_t(\mN_+)$ for $s\geq0$. To interpret this inclusion, one can take $\mN_+$ to be generated by observations on a half-infinite time interval, with the orientation chosen so that positive $t$ removes the earliest part of the record.

The Kolmogorov, or $K$, property can be formulated through the existence of such a one-sided algebra whose translates generate the full algebra and whose tail is trivial:
\begin{align}\label{eq:classical-K-algebra}
    \bigvee_{t\in\mathbb R}\alpha_t(\mN_+)=\mM,
    \qquad
    \bigcap_{t\geq0}\alpha_t(\mN_+)=\mathbb C1\ .
\end{align}
The first condition says that the complete record resolves all observables. The second says that no nonconstant observable remains accessible arbitrarily far into the retained half of the record. The analogous conditions with integer time describe a $K$-automorphism. They provide an inclusion-based formulation of loss of memory and motivate the future and past algebras used in quantum ergodic theory~\cite{ouseph2024local}.

\paragraph{Example: a Bernoulli shift:}
For independent uniform bits $x_n$, $n\in\mathbb Z$, with the shift $\alpha(x_n)=x_{n+1}$ and $\mN_+$ the algebra generated by $\{x_n:n\geq0\}$, one has $\alpha^k(\mN_+)=\operatorname{vN}\{x_n:n\geq k\}\subsetneq\mN_+$ for $k>0$, the translates generate $\mM$, and the tail is trivial by the Kolmogorov zero--one law: the retained description forgets one coordinate per step although the shift is invertible and measure preserving. This is a subsystem coarse-graining, the retained and discarded coordinates being tensor factors. Within the finite-measure invertible measure-preserving setting used here, strict one-sided nesting of this kind is possible for subsystem algebras but not for the region-type subspace algebras above. In a finite measure space $\Phi_t(O)\subseteq O$ and measure preservation force $\mu(O\setminus\Phi_t(O))=0$, so a region, and with it a subspace algebra $\mM_O$ in the nonatomic setting of Appendix~\ref{subsec:classical-region-subsystem}, cannot shrink strictly under a measure-preserving flow.

\paragraph{Coarse-graining along a decreasing family:}
Each $\alpha_t(\mN_+)$ has a $\mu$-preserving conditional expectation $\mE_t:\mM\to\alpha_t(\mN_+)$. For $s\geq t$, the tower property gives
\begin{align}
    \mE_s\mE_t=\mE_s\ .
\end{align}
For a fixed state with density $\rho$, its coarse-grained densities $\rho_t=\mE_t(\rho)$ therefore satisfy
\begin{align}
    H(\rho_s)\geq H(\rho_t),
    \qquad
    S(\rho\|\rho_s)\geq S(\rho\|\rho_t),
    \qquad s\geq t,
\end{align}
whenever the entropies are finite. Here the loss of information follows from the nesting of the retained algebras. It does not require irreversible microscopic dynamics. For a trivial tail, the reverse martingale theorem further gives $\mE_t(f)\to\mu(f)1$ in $L^2$ for every $f\in L^2(X,\mu)$.

\subsection{Classical avatar of modular flow in Hamiltonian systems}
\label{subsec:classical-state-flow}

So far, the state and the dynamics have been independent inputs. A probability density assigns weights to microstates, but measure theory alone does not turn these weights into a motion of the microstates. A Poisson bracket supplies the missing structure: it converts a function into a Hamiltonian vector field. Applied to the logarithm of a density, it gives a distinguished state-dependent dynamics. This construction will let us compare the dynamics associated with a state and with its coarse description.

\paragraph{Poisson brackets and Hamiltonian motion:}
A Poisson structure on a smooth manifold $X$ is an antisymmetric bivector $\Pi$ whose bracket
\begin{align}
    \{f,g\}=\Pi(df,dg),\qquad f,g\in C^\infty(X),
\end{align}
satisfies the Jacobi identity. The bracket obeys the Leibniz rule in each argument, so fixing a smooth real Hamiltonian $H$ defines a vector field. We use the convention
\begin{align}\label{eq:classical-Hamiltonian-convention}
    X_H(f)&=\{f,H\},\qquad
    \alpha_t^H(f)=f\circ\Phi_{-t}^H,\nn\\
    \delta_H(f)&:=\left.\partial_t\alpha_t^H(f)\right|_{t=0}=\{H,f\},
\end{align}
where $\Phi_t^H$ is the flow of $X_H$. Thus the generator on observables has the opposite sign from the vector field transporting points. The Jacobi identity ensures that Hamiltonian evolution preserves the Poisson bracket. In canonical coordinates, $\{q_i,p_j\}=\delta_{ij}$ and the point trajectories satisfy the usual Hamilton equations.

\paragraph{Casimirs and symplectic leaves:}
A Casimir is a function $c$ with $\{c,f\}=0$ for every smooth $f$. It is conserved by every Hamiltonian, rather than by one chosen dynamics. The points accessible from a given microstate by successive Hamiltonian motions form a symplectic leaf. Every Casimir is constant on each leaf, although global Casimirs need not distinguish all leaves. A leaf is therefore a classical analogue of an irreducible sector for Hamiltonian motion. This analogy concerns the Poisson bracket: the pointwise algebra $L^\infty(X,\mu)$ remains abelian, and restricting it to a leaf does not turn it into a nonabelian factor.

\paragraph{Liouville measure and the state Hamiltonian:}
For the explicit constructions below, take $X$ to be a connected compact symplectic manifold without boundary, with symplectic form $\omega$ and normalized Liouville measure
\begin{align}
    d\mu=\frac{\omega^n/n!}{\int_X\omega^n/n!},\qquad \dim X=2n.
\end{align}
Hamiltonian flows preserve this measure. Equivalently,
\begin{align}\label{statepreserve}
    \int_X\{f,g\}\,d\mu=0
\end{align}
for smooth $f,g$. Compactness makes smooth Hamiltonian flows complete and avoids boundary terms. On a noncompact phase space these properties, together with the required integrability, must be checked separately.

Let $\rho$ be a smooth strictly positive probability density relative to $\mu$. Define its {\it state Hamiltonian} and the associated evolution by
\begin{align}\label{eq:classical-state-Hamiltonian}
    K_\rho=-\log\rho,\qquad
    \Phi_t^\rho:=\Phi_t^{K_\rho},\qquad
    \alpha_t^\rho:=\alpha_t^{K_\rho}.
\end{align}
The infinitesimal evolution of an observable is
\begin{align}\label{eq:classical-state-derivation}
    \delta_\rho(f)=\{K_\rho,f\}=\{f,\log\rho\}.
\end{align}
The flow moves along level sets of the density, since $X_{K_\rho}(\rho)=0$. It preserves both $\mu$ and the state $\omega_\rho(f)=\int_X\rho f\,d\mu$. For a Gibbs density $\rho=Z^{-1}e^{-\beta H}$, we have $K_\rho=\beta H+\log Z$ and $\Phi_t^\rho=\Phi_{\beta t}^H$. Thus the state recovers the Hamiltonian motion, with a rescaled parameter; for a general density there is no prior identification of this parameter with physical time.

\paragraph{The classical KMS identity:}
The logarithm is distinguished by an integration-by-parts identity. Using \eqref{statepreserve},
\begin{align}
    0=\int_X\{g,f\rho\}\,d\mu
    =-\omega_\rho(\{f,g\})+
    \omega_\rho\bigl(f\{g,\log\rho\}\bigr).
\end{align}
Consequently,
\begin{align}\label{KMSclassical}
    \omega_\rho(\{f,g\})=\omega_\rho(f\delta_\rho(g)).
\end{align}
For a Gibbs state this reads $\omega_\beta(\{f,g\})=\beta\omega_\beta(f\delta_H(g))$. This is the classical, or static, KMS condition, written with the time orientation in \eqref{eq:classical-Hamiltonian-convention}~\cite{vanDeVen2022}. It is stronger than stationarity: for a given smooth positive density, testing against all smooth $f$ determines $\delta_\rho(g)$ for every $g$. Conversely, on a connected symplectic manifold, imposing this identity for a specified Hamiltonian dynamics determines a smooth positive density to be Gibbs, up to normalization. The identity also motivates formulations on algebras of local observables in infinite systems, where existence of the dynamics and equilibrium states requires additional hypotheses; a global Gibbs density need not exist.

\paragraph{Poisson modular flow:}
There is a geometric formulation of the same construction. For a smooth positive measure $\nu$, define its Poisson modular vector field by
\begin{align}
    \mathfrak m_\nu(f):=\operatorname{div}_\nu X_f.
\end{align}
This measures the failure of Hamiltonian motions to preserve $\nu$~\cite{Weinstein1997}. For $d\nu=\rho\,d\mu$ and Liouville $\mu$, the divergence formula gives
\begin{align}
    \mathfrak m_\nu(f)=X_f(\log\rho)
    =\{f,-\log\rho\}=X_{K_\rho}(f).
\end{align}
Thus $\Phi_t^\rho$ is precisely this Poisson modular flow; its action on observables uses the inverse pullback, so $\delta_\rho=-\mathfrak m_\nu$. On a general Poisson manifold and for an arbitrary smooth reference measure, the formula instead reads $\mathfrak m_{\rho\mu}=\mathfrak m_\mu+X_{-\log\rho}$. A reference measure invariant under every Hamiltonian need not exist. Our symplectic setting removes this additional obstruction.

The construction uses the Poisson structure in an essential way. It is distinct from the multiplication group $U_t^\rho$ introduced above and from the trivial modular automorphism group of the abelian von Neumann algebra. It also supplies no automatic ergodic or entropy-production property. A nonconstant $\rho$ is itself a conserved observable, so the flow cannot be ergodic on the full probability space. Questions about ergodicity and mixing must be posed on invariant sectors with suitable invariant measures. The comparison of full and coarse state Hamiltonians, however, is meaningful without assuming ergodicity.

\subsection{Subsystem Poisson flows in classical many-body systems and field theory}
\label{subsec:classical-subalgebra-Hamiltonian}

A coarse description determines a reduced state and hence, when the coarse space is Poisson, its own state-generated dynamics. We can then ask whether the reduced dynamics agrees with what the full dynamics does to the observed variables. Comparing the two motions gives a classical model of the subsystem flows studied in this paper. The comparison separates evolution already visible within the retained algebra from evolution involving the discarded variables.

\paragraph{Poisson subsystems:} We say a subalgebra $\mN_Y\subset \mM$ of smooth observables is a Poisson subsystem if for all $b_1,b_2\in \mN_Y$, we have $\{b_1,b_2\}\in \mN_Y$. In the geometric setting below, $\pi:X\to Y$ is a smooth surjective Poisson submersion:
\begin{align}\label{eq:classical-Poisson-subsystem}
    \{\pi^*f,\pi^*g\}_X=\pi^*\{f,g\}_Y,
    \qquad f,g\in C^\infty(Y).
\end{align}
The smooth coarse observables $\mN=\pi^*C^\infty(Y)$ are closed under the bracket. This condition allows us to speak of Hamiltonian dynamics within the coarse description.

The simplest setting makes both the geometric and probabilistic structures explicit. Take a product $X=Y\times Z$ of connected compact symplectic manifolds, with product symplectic structure and normalized Liouville measure $\mu=\mu_Y\otimes\mu_Z$. The projection to $Y$ is Poisson, and the associated von Neumann subalgebra is $\mN_Y=L^\infty(Y,\mu_Y)\otimes1$. For a smooth positive density $\rho$, restriction of the state gives the marginal density
\begin{align}\label{eq:classical-marginal-Hamiltonian}
    \rho_Y(y)=\int_Z\rho(y,z)\,d\mu_Z(z),\qquad
    K_Y=-\log\rho_Y.
\end{align}
The expectation $\mE_Y$ from Appendix~\ref{subsec:classical-region-subsystem} extends this marginal back to $X$ as $\rho_0(y,z)=\rho_Y(y)$. Here, $\rho_Y=\pi_*\rho$ is the classical analog of the reduced state on $Y$, and $\rho_0=\pi^*\rho_Y=\pi^*\pi_* \rho$ is the pullback of $\rho_Y$ to $X$, which is the coarse-grained version of $\rho$. Its Hamiltonian is $K_0=\pi^*K_Y$, whose flow is
\begin{align}
    \Phi_t^0(y,z)=(\Phi_t^Y(y),z).
\end{align}
This provides a definite lift of the reduced motion to the full phase space. For a general Poisson submersion, a specified coarse Hamiltonian can likewise be pulled back, but the marginal measure need not be Liouville and a complementary subsystem need not exist. We use the product setting below to keep the marginal, its reference measure, and its lifted Hamiltonian unambiguous.

\paragraph{Agreement between reduced and full Poisson dynamics:}
Agreement means that the observed trajectory is determined by its initial coarse state and follows the flow generated by $K_Y$:
\begin{align}\label{eq:classical-state-intertwining}
    \pi\circ\Phi_t^\rho=\Phi_t^Y\circ\pi.
\end{align}
For our complete smooth flows, this is equivalent to the infinitesimal condition
\begin{align}\label{eq:classical-state-defect}
    \{\pi^*g,K_\rho\}_X-\pi^*\{g,K_Y\}_Y
    =\{\pi^*g,K_\rho-K_0\}_X=0,
    \qquad g\in C^\infty(Y).
\end{align}
The left side measures the part of the observed velocity missed by the reduced state Hamiltonian. A Hamiltonian invisible to all coarse observables belongs to the Poisson centralizer
\begin{align}
    \mN^{\mathrm c}_{\mathrm P}
    :=\{h\in C^\infty(X):\{h,b\}=0\text{ for every }b\in\mN\}.
\end{align}
Thus agreement is precisely the condition $K_\rho-K_0\in\mN^{\mathrm c}_{\mathrm P}$. This is an analogue of a relative commutant defined using Poisson brackets. The ordinary operator commutant of a classical subalgebra contains the whole abelian observable algebra and cannot express this condition.

In the product setting, nondegeneracy of the symplectic form and connectedness of $Y$ imply that this centralizer consists of functions of $z$ alone. Hence \eqref{eq:classical-state-intertwining} is equivalent to
\begin{align}\label{eq:classical-product-state-closure}
    K_\rho(y,z)=K_Y(y)+K_Z(z),
    \qquad
    \rho(y,z)=\rho_Y(y)\rho_Z(z).
\end{align}
Normalization of the marginal fixes $\int_Z\rho_Z\,d\mu_Z=1$. The condition \eqref{eq:classical-state-defect} implies that the coarse observables have the same Poisson evolution under $\rho_0$ and $\rho$. In other words:
\begin{align}
    \forall g\in C^\infty(Y):\qquad \alpha^\rho_t\circ\alpha_{-t}^0(\pi^*g)=\pi^*g\ .
\end{align}
This compares full and reduced observable evolution, as in Section~\ref{subsec:quantum-coarse}; the state-preserving expectation criterion itself differs in the classical and quantum theories.

\paragraph{Comparing the two flows:}
Define the relative transformation by first undoing the lifted coarse motion and then applying the full motion:
\begin{align}\label{eq:classical-relative-transformation}
    R_t=\Phi_t^\rho\circ\Phi_{-t}^0,
    \qquad
    \mathcal V_t(f):=f\circ R_t^{-1}
    =\alpha_t^\rho\circ\alpha_{-t}^0(f).
\end{align}
For each $t$, this is a symplectic transformation and a measure-preserving automorphism, respectively. Since the lifted coarse flow preserves $\mN_Y$,
\begin{align}\label{eq:classical-relative-algebra-flow}
    \mathcal V_t(\mN_Y)=\alpha_t^\rho(\mN_Y).
\end{align}
The relative transformation therefore follows the same moving subalgebra as the full dynamics, while canceling the chosen internal coarse evolution. It fixes every coarse observable exactly when \eqref{eq:classical-state-intertwining} holds. For a product state, it can still move hidden variables (the relative commutant): $R_t(y,z)=(y,\Phi_t^Z(z))$.

The infinitesimal generator of $R_t$ is the difference of the state Hamiltonians:
\begin{align}\label{eq:classical-relative-Hamiltonian}
    K_{\mathrm{rel}}&=K_\rho-K_0=-\log\frac{\rho}{\rho_0},\nn\\
    \left.\partial_tR_t\right|_{t=0}&=X_{K_{\mathrm{rel}}},\qquad
    \left.\partial_t\mathcal V_t(f)\right|_{t=0}=\{K_{\mathrm{rel}},f\}.
\end{align}

Every normal state of the hidden factor, with density $\rho_Z$ relative to $\mu_Z$,
gives an expectation $\mE_{\rho_Z}:\mM\to\mN_Y$:
\begin{equation}
 (\mE_{\rho_Z}f)(y)=\int_Z\rho_Z(z)f(y,z)\,d\mu_Z(z).
\end{equation}
These states parameterize the expectations with $y$-independent averaging.
General normal expectations allow measurable probability kernels $\nu_y(dz)$
depending on $y$, with $\mu_Y(dy)\nu_y(dz)\ll\mu_Y\otimes\mu_Z$:
$\mE_\nu f(y)=\int_Z f(y,z)\nu_y(dz)$.

\paragraph{Example: a uniform marginal with nonautonomous observables:}
Take $Y$ and $Z$ to be two-dimensional tori with angular coordinates $(q,p)$ and $(Q,P)$, each of period $2\pi$, and the product symplectic form $dq\wedge dp+dQ\wedge dP$. Consider
\begin{align}
    \rho(q,p,Q,P)=\frac{e^{\epsilon\cos(q-Q)}}{\mathcal Z_\epsilon},
    \qquad
    \mathcal Z_\epsilon=\int_X e^{\epsilon\cos(q-Q)}\,d\mu,
    \qquad \epsilon\ne0.
\end{align}
Translation invariance of the angular integral gives $\rho_Y=1$, so $K_0=0$ and the reduced state predicts no motion. The full state Hamiltonian is $K_\rho=-\epsilon\cos(q-Q)+\log\mathcal Z_\epsilon$, giving
\begin{align}
    \dot q=\dot Q=0,\qquad
    \dot p=-\epsilon\sin(q-Q),\qquad
    \dot P=\epsilon\sin(q-Q).
\end{align}
For example, the bounded coarse observable $e^{ip}$ evolves as
\begin{align}
    \alpha_t^\rho(e^{ip})
    =\exp\bigl(i[p+\epsilon t\sin(q-Q)]\bigr).
\end{align}
Its future value depends on the hidden angle $Q$, so it leaves $\mN_Y$ for $t\ne0$. The full state and its uniform marginal are stationary, but the retained observables do not evolve autonomously. Here $R_t=\Phi_t^\rho$ directly displays the motion that the marginal state Hamiltonian misses. The example also exhibits the failure of monotonicity announced at the beginning of this section: the relative Hamiltonian $K_{\mathrm{rel}}=K_\rho-K_0=-\epsilon\cos(q-Q)+\log\mathcal Z_\epsilon$ is indefinite for every $\epsilon\ne0$, since $\mathcal Z_\epsilon=I_0(\epsilon)<e^{|\epsilon|}$; its $\mu$-average $\log\mathcal Z_\epsilon$ is positive while its $\rho$-average $-S(\rho\|\rho_0)$ is negative. Its quantum counterpart $K_{\mM}-K_{\mN}$ is a nonnegative quadratic form, by \eqref{eq:modular-inequality} and Proposition~\ref{prop:sufficiency}.

\paragraph{Relative entropy and the information discarded:}
The same relative logarithm appeared in the phase construction at the beginning of this section. Here it generates a vector field, while its expectation gives
\begin{align}\label{eq:classical-relative-entropy-Hamiltonian}
    S(\rho\|\rho_0)=-\omega_\rho(K_{\mathrm{rel}})
    =H(\rho_0)-H(\rho).
\end{align}
The vector field alone does not determine this number: on a connected symplectic manifold it determines its Hamiltonian only up to a constant. The normalized densities fix that constant. On a general Poisson manifold, Casimirs give a larger ambiguity.

Information loss and failure of autonomous coarse dynamics are also distinct. For a product state, the relative transformation acts only on $Z$, yet
\begin{align}
    S(\rho_Y\rho_Z\|\rho_Y\otimes1)
    =\int_Z\rho_Z\log\rho_Z\,d\mu_Z
\end{align}
can be positive. Coarse-graining has discarded the hidden distribution even without correlations. More generally,
\begin{align}
    S(\rho\|\rho_Y\otimes1)
    =S(\rho\|\rho_Y\rho_Z)+S(\rho_Z\|1),
\end{align}
where $\rho_Z$ is the other marginal. The first term is the mutual information and the second measures the hidden marginal relative to the chosen reference. Neither this entropy nor the relative transformation supplies a monotone entropy-production law without an additional coarse-graining or dynamical assumption.

\paragraph{Subregions and the passage to field theory:}
The measurable region algebras $\mM_O$ remain valid coarse descriptions. What can fail is the smoothness needed to assign their coarse densities a Hamiltonian flow. For the region expectation $\mE_O$,
\begin{align}
    \rho_0=\mE_O(\rho)=\chi_O\rho+c_O\chi_{O^c},\qquad
    c_O=\frac{1}{\mu(O^c)}\int_{O^c}\rho\,d\mu.
\end{align}
For an open region with smooth boundary this density is generally discontinuous there. In the interior of $O$, $K_0=K_\rho$; in the interior of the complement, $K_0$ is constant. These local statements do not specify a smooth global Hamiltonian flow across the boundary. Additional regularization or boundary conditions are needed. This is a limitation of the Hamiltonian construction, not of the measure-theoretic net.

Finally, every classical coarse algebra still admits a state-preserving conditional expectation for a faithful probability state, even when its Poisson state flow fails to preserve that algebra. In the product example this expectation is explicitly
\begin{align}
    \mE_{\rho,Y}(f)(y)=\int_Z f(y,z)\frac{\rho(y,z)}{\rho_Y(y)}\,d\mu_Z(z).
\end{align}
Its existence does not imply \eqref{eq:classical-state-intertwining}. In quantum modular theory the two questions merge: by Takesaki's theorem (Section~\ref{subsec:quantum-coarse} and Theorem~\ref{thm:correctable}), a reference-preserving expectation onto a subalgebra exists if and only if the subalgebra is invariant under the modular flow of the reference. This is the additional structure that the quantum theory contributes to the comparison of full and coarse dynamics.

The classical construction thus motivates studying how relative state evolutions transport retained observables. In QFT a spacetime region selects local field observables, whereas a region $O\subset X$ above selects possible microstates. The connection to the subregion algebras of QFT, developed in the companion paper~\cite{companionQFT}, is through inclusions of observable algebras and their state-dependent flows; it does not identify these two notions of region.

\paragraph{A classical half-sided inclusion:}
A simple example of strict nesting is supplied by a classical chiral current $j(x)$ on the line. On rapidly decreasing smooth configurations, choose
\begin{align}
    \{j(x),j(y)\}=-\partial_x\delta(x-y),
    \qquad P=\frac12\int_{\mathbb R}j(x)^2\,dx,
    \qquad \tau_a j(x)=j(x+a)\ .
\end{align}
The nonnegative Hamiltonian $P$ generates $\tau_a$ through $\{P,j(x)\}=\partial_xj(x)$. Let $\mN_b$ be generated by $1$ and $j(f)$ with $f\in C_c^\infty((b,\infty))$. Since $\tau_a j(f)=j(f(\,\cdot-a))$,
\begin{align}
    \tau_a(\mN_b)=\mN_{b+a}\subsetneq\mN_b\quad(a>0),
    \qquad
    \bigcup_{b\in\mathbb R}\mN_b=\mM_{\mathrm{loc}},
    \qquad
    \bigcap_{b\geq0}\mN_b=\mathbb R1\ .
\end{align}
Here $\mM_{\mathrm{loc}}$ is the polynomial algebra of compactly smeared currents; $P$ is a global functional generating its automorphisms. The full translation group is reversible, but its positive half acts on $\mN_0$ by injective, nonsurjective Poisson endomorphisms. The decreasing family progressively discards accessible field degrees of freedom.

The dilation Hamiltonian $D=\frac12\int_{\mathbb R}xj(x)^2\,dx$ gives a further connection with half-sided modular inclusions. Its flow $\delta_s$ satisfies
\begin{align}
    \{D,P\}=P,
    \qquad \delta_s(\mN_b)=\mN_{e^s b},
    \qquad \delta_s\tau_a\delta_{-s}=\tau_{e^s a}\ .
\end{align}
Thus dilations preserve $\mN_0$ and act half-sidedly on $\mN_b$ for $b>0$. These are classical Hamiltonian realizations of the translation--dilation relations. Their construction uses the Poisson structure; identifying the dilation with a state-generated modular evolution requires additional input.

\paragraph{Classical error correction and the classical sufficiency gap.}
The constructions of this section can be read in the language that Section~\ref{sec:QEC} will use for the quantum theory. Restriction of a density to the coarse algebra $\mN=L^\infty(X,\Sigma_0,\mu)$ is an erasure channel. The reference-preserving expectation $\mE_{\mN}$ always exists, and the densities it preserves,
\begin{align}
 \mathcal C_{\mN}=\{\rho:\ \mE_{\mN}(\rho)=\rho\},
\end{align}
are exactly the $\Sigma_0$-measurable densities; in the product setting these are the densities $\rho_Y(y)\,1$ whose hidden marginal is the reference $\mu_Z$. Every such density is recovered from its restriction by $\mE_{\mN}$ itself, so $\mathcal C_{\mN}$ is an exactly correctable code for the erasure, and its logical algebra is $\mN$. For a general density, the chain rule $S(\rho\|1)=S(\rho\|\rho_0)+S(\rho_0\|1)$ with $\rho_0=\mE_{\mN}(\rho)$, which uses only that $\log\rho_0$ is $\Sigma_0$-measurable, shows that
\begin{align}\label{eq:classical-sufficiency-gap}
 S(\rho\|\rho_0)=S(\rho\|1)-S(\rho_0\|1)=-\omega_\rho(K_{\mathrm{rel}})\ \ge0
\end{align}
is the loss of relative entropy to the reference under restriction to $\Sigma_0$. By the classical equality case of monotonicity, this vanishes if and only if $\Sigma_0$ is a sufficient statistic for the pair $(\rho,\mu)$, that is, if and only if $\rho\in\mathcal C_{\mN}$~\cite{OhyaPetz1993}. The expectation value of the relative Hamiltonian, which generates the classical subsystem flow \eqref{eq:classical-relative-transformation}, is therefore the classical sufficiency gap. Two features distinguish this from the quantum case. The correctable code exists for every coarse algebra, independently of the dynamics; whether the retained observables evolve autonomously is a separate question answered by \eqref{eq:classical-product-state-closure}. And there is a single gap, because a commutative algebra has no commutant distinct from itself. In the quantum theory both features change: the existence of a reference-preserving expectation becomes equivalent to autonomy (Takesaki), and the subsystem-flow generator acquires a second gap for the inclusion of commutants (Proposition~\ref{prop:sufficiency}).

\bibliographystyle{unsrt}
\bibliography{subsystemflow}

\end{document}
